% !TEX program = xelatex
% ===========================================================================
%  《主题》—— 中文数学习题课讲义（Loom 模板 · 课前自学版）
%  必须用 XeLaTeX 编译两遍；QA 用 scripts/qa_check.py。
%  第一行的魔法注释让 VS Code (LaTeX Workshop) 走 XeLaTeX——别删（默认是
%  pdfLaTeX，会报 xeCJK requires XeTeX）。
% ===========================================================================
\documentclass{loom}

% 挑战题 ★ 用 Zapf Dingbats 绘制，不依赖系统字体（Optima 等）是否含 U+2605
\usepackage{pifont}
\usepackage{adjustbox}

% 页眉线（可选）；老师偏好极简时可整行注释掉
%\runningthread{《主题》 · 课前自学讲义}

% ---------------------------------------------------------------------------
%  讲义专用组件：题面（茜红结）、解析（靛蓝结）、三段式解答标签
%  —— 这些宏定义住在 main.tex（不在 cls），以便随讲义微调。
% ---------------------------------------------------------------------------
% 题目框：breakable=false 保证一题不跨页。默认不留答题空白区（老师偏好）。
% —— 若要做「打印作答版」，取消下面 \answerspace 定义与 problem 框 after upper
%    一行的注释即可，题末会追加一块纯白书写面板（默认 5 行，可 \answerspace[8]）。
% \newcommand{\answerspace}[1][5]{\par\medskip\noindent
%   \colorbox{white}{\parbox[t][#1\baselineskip][t]{\dimexpr\linewidth-2\fboxsep\relax}{\strut}}}
\newtcolorbox{problem}[1]{knotmadder, breakable=false,
  before upper={{\headingfont\bfseries\color{madder}#1}\par\smallskip}%
  % , after upper={\answerspace}   % 打印作答版：取消本行注释
}
\newtcolorbox{solution}[1]{knotindigo,
  before upper={{\headingfont\bfseries\color{indigo}#1}\par\smallskip}}

\newcommand{\think}{\par\smallskip\noindent
  {\headingfont\bfseries\color{indigo}〔怎么想〕}\enspace}
\newcommand{\solve}{\par\smallskip\noindent
  {\headingfont\bfseries\color{inkiron}〔解〕}\enspace}
\newcommand{\solvehead}[1]{\par\smallskip\noindent
  {\headingfont\bfseries\color{inkiron}〔#1〕}\enspace}
\newcommand{\insight}[1]{\par\smallskip\noindent
  {\headingfont\bfseries\color{weld!60!inkiron}〔方法点睛〕}\enspace
  {\strandfont\small #1}\par}
\newcommand{\hint}[1]{\par\smallskip\noindent
  {\strandfont\small{\raisebox{-0.8pt}{\loomtile{weld}}\hspace{4pt}}%
  {\bfseries\color{weld!60!inkiron}入手提示}\enspace{\color{inkiron!85!white}#1}}\par}
% 页边旁批：放"想一想/追问"、∑/∏ 记号说明，不打断正文
\newcommand{\sidenote}[1]{\marginnote{\raggedright\strandfont\footnotesize
  \color{selvage}#1}}
\newcommand{\starred}{{\color{weld!70!inkiron}\ding{72}}}

\setlist[enumerate]{leftmargin=2.3em, itemsep=2pt, topsep=3pt}
\setlist[itemize]{leftmargin=1.6em, itemsep=2pt, topsep=3pt}


\usetikzlibrary{calc,angles,quotes,arrows.meta}
\definecolor{esblue}{HTML}{355C86}
\definecolor{esorange}{HTML}{B76543}
\DeclareRobustCommand{\escond}[1]{\tikz[baseline=(n.base)]
  \node[draw,circle,inner sep=.65pt,line width=.35pt,font=\footnotesize] (n) {#1};}
\newcommand{\esdiagram}[1]{\par\smallskip{\centering
  \adjustbox{max width=.84\linewidth,max height=5.8cm}{#1}
  \par}\smallskip}
\newif\ifteachernotes\teachernotestrue
\newif\ifshowhints\showhintstrue
\let\originalhint\hint
\renewcommand{\hint}[1]{\ifshowhints\originalhint{#1}\fi}
\newcommand{\diagram}[1]{\par\smallskip\begin{center}#1\end{center}\smallskip}
\let\originalblock\block
\renewcommand{\block}[1]{\needspace{4\baselineskip}\originalblock{#1}}
\hypersetup{pdftitle={geometry Models},pdfauthor={Hence},pdfsubject={初中几何模型培优讲义：夹半角、手拉手、双等腰、八字形}}
\renewcommand{\loomcovermotto}{从图形中发现关系，在推理中理解几何}
\renewcommand{\contentsname}{目录}
% BEGIN GENERATED: chapter-diagrams
% 可编辑的几何题图；所有坐标与数值已核验。
\newcommand{\figOne}{%
\begin{tikzpicture}[scale=0.95,every node/.style={font=\small,inner sep=2pt},line width=.65pt]
\coordinate (A) at (0,3);
\coordinate (B) at (0,0);
\coordinate (C) at (3,0);
\coordinate (D) at (3,3);
\coordinate (E) at (1,0);
\coordinate (F) at (3,1.5);
\draw[inkiron] (A) -- (B) -- (C) -- (D) -- (A);
\draw[inkiron] (A) -- (E) -- (F) -- (A);
\pic[draw=madder,angle radius=4.2mm,angle eccentricity=1.45,"$45^\circ$"] {angle=E--A--F};
\node[above left] at (A) {$A$};
\node[below left] at (B) {$B$};
\node[below right] at (C) {$C$};
\node[above right] at (D) {$D$};
\node[below] at (E) {$E$};
\node[right] at (F) {$F$};
\end{tikzpicture}}

\newcommand{\figIntroA}{%
\begin{tikzpicture}[scale=0.73,every node/.style={font=\small,inner sep=2pt},line width=.65pt]
\coordinate (A) at (0,3);
\coordinate (B) at (0,0);
\coordinate (C) at (3,0);
\coordinate (D) at (3,3);
\coordinate (E) at (1,0);
\coordinate (F) at (3,1.5);
\draw[inkiron] (A) -- (B) -- (C) -- (D) -- (A);
\draw[inkiron] (A) -- (E) -- (F) -- (A);
\pic[draw=madder,angle radius=4.2mm,angle eccentricity=1.45,"$45^\circ$"] {angle=E--A--F};
\node[above left] at (A) {$A$};
\node[below left] at (B) {$B$};
\node[below right] at (C) {$C$};
\node[above right] at (D) {$D$};
\node[below] at (E) {$E$};
\node[right] at (F) {$F$};
\node[font=\small] at (current bounding box.south) [below=6pt] {图 A 顶点夹半角};
\end{tikzpicture}}

\newcommand{\figOneSolve}{%
\begin{tikzpicture}[scale=0.76,every node/.style={font=\small,inner sep=2pt},line width=.65pt]
\coordinate (A) at (0,3);
\coordinate (B) at (0,0);
\coordinate (C) at (3,0);
\coordinate (D) at (3,3);
\coordinate (E) at (1,0);
\coordinate (F) at (3,1.5);
\coordinate (G) at (3,4);
\draw[inkiron] (A) -- (B) -- (C) -- (D) -- (A);
\draw[inkiron] (A) -- (E) -- (F) -- (A);
\draw[indigo,dashed] (D) -- (G) -- (A);
\draw[indigo,dashed] (F) -- (G);
\pic[draw=madder,angle radius=4.2mm,angle eccentricity=1.45,"$45^\circ$"] {angle=E--A--F};
\pic[draw=madder,angle radius=4.2mm,angle eccentricity=1.45,"$45^\circ$"] {angle=F--A--G};
\node[above left] at (A) {$A$};
\node[below left] at (B) {$B$};
\node[below right] at (C) {$C$};
\node[above right] at (D) {$D$};
\node[below] at (E) {$E$};
\node[right] at (F) {$F$};
\node[above right] at (G) {$G$};
\path (D) -- node[right,fill=white] {$BE$} (G);
\end{tikzpicture}}

\newcommand{\figIntroB}{%
\begin{tikzpicture}[scale=0.73,every node/.style={font=\small,inner sep=2pt},line width=.65pt]
\coordinate (A) at (0,3);
\coordinate (B) at (0,0);
\coordinate (C) at (3,0);
\coordinate (D) at (3,3);
\coordinate (E) at (1,0);
\coordinate (F) at (0,1.5);
\coordinate (P) at (0.6,1.2);
\draw[inkiron] (A) -- (B) -- (C) -- (D) -- (A);
\draw[inkiron] (A) -- (E);
\draw[inkiron] (F) -- (C);
\pic[draw=madder,angle radius=4.2mm,angle eccentricity=1.45,"$45^\circ$"] {angle=E--P--C};
\node[above left] at (A) {$A$};
\node[below left] at (B) {$B$};
\node[below right] at (C) {$C$};
\node[above right] at (D) {$D$};
\node[below] at (E) {$E$};
\node[left] at (F) {$F$};
\node[above right] at (P) {$P$};
\node[font=\small] at (current bounding box.south) [below=6pt] {图 B 交叉角与平行转移};
\end{tikzpicture}}

\newcommand{\figIntroC}{%
\begin{tikzpicture}[scale=0.73,every node/.style={font=\small,inner sep=2pt},line width=.65pt]
\coordinate (A) at (0,3);
\coordinate (B) at (0,0);
\coordinate (C) at (3,0);
\coordinate (D) at (3,3);
\coordinate (E) at (1,0);
\coordinate (F) at (1.5,3);
\coordinate (P) at (0.6,1.2);
\draw[inkiron] (A) -- (B) -- (C) -- (D) -- (A);
\draw[inkiron] (A) -- (E);
\draw[inkiron] (B) -- (F);
\pic[draw=madder,angle radius=4.2mm,angle eccentricity=1.45,"$45^\circ$"] {angle=F--P--A};
\node[above left] at (A) {$A$};
\node[below left] at (B) {$B$};
\node[below right] at (C) {$C$};
\node[above right] at (D) {$D$};
\node[below] at (E) {$E$};
\node[above] at (F) {$F$};
\node[left] at (P) {$P$};
\node[font=\small] at (current bounding box.south) [below=6pt] {图 C 交叉角与垂直转移};
\end{tikzpicture}}

\newcommand{\figIntroD}{%
\begin{tikzpicture}[scale=0.73,every node/.style={font=\small,inner sep=2pt},line width=.65pt]
\coordinate (A) at (0,3);
\coordinate (B) at (0,0);
\coordinate (C) at (3,0);
\coordinate (D) at (3,3);
\coordinate (E) at (0.3,3);
\coordinate (F) at (1.3,0);
\coordinate (M) at (0,2.2);
\coordinate (N) at (3,0.7);
\coordinate (P) at (0.68,1.86);
\draw[inkiron] (A) -- (B) -- (C) -- (D) -- (A);
\draw[inkiron] (E) -- (F);
\draw[inkiron] (M) -- (N);
\pic[draw=madder,angle radius=4.2mm,angle eccentricity=1.45,"$45^\circ$"] {angle=F--P--N};
\node[above left] at (A) {$A$};
\node[below left] at (B) {$B$};
\node[below right] at (C) {$C$};
\node[above right] at (D) {$D$};
\node[above] at (E) {$E$};
\node[below] at (F) {$F$};
\node[left] at (M) {$M$};
\node[right] at (N) {$N$};
\node[above right] at (P) {$P$};
\node[font=\small] at (current bounding box.south) [below=6pt] {图 D 两条线同时平移};
\end{tikzpicture}}

\newcommand{\figTwo}{%
\begin{tikzpicture}[scale=0.36,every node/.style={font=\small,inner sep=2pt},line width=.65pt]
\coordinate (A) at (0,9);
\coordinate (B) at (0,0);
\coordinate (C) at (9,0);
\coordinate (D) at (9,9);
\coordinate (M) at (3,9);
\coordinate (N) at (9,4.5);
\draw[inkiron] (A) -- (B) -- (C) -- (D) -- (A);
\draw[inkiron] (B) -- (M) -- (N) -- (B);
\node[above left] at (A) {$A$};
\node[below left] at (B) {$B$};
\node[below right] at (C) {$C$};
\node[above right] at (D) {$D$};
\node[above] at (M) {$M$};
\node[right] at (N) {$N$};
\end{tikzpicture}}

\newcommand{\figTwoSolve}{%
\begin{tikzpicture}[scale=0.34,every node/.style={font=\small,inner sep=2pt},line width=.65pt]
\coordinate (A) at (0,9);
\coordinate (B) at (0,0);
\coordinate (C) at (9,0);
\coordinate (D) at (9,9);
\coordinate (M) at (3,9);
\coordinate (N) at (9,4.5);
\coordinate (P) at (9,-3);
\coordinate (Q) at (5.4,7.2);
\draw[inkiron] (A) -- (B) -- (C) -- (D) -- (A);
\draw[inkiron] (B) -- (M) -- (N) -- (B);
\draw[indigo,dashed] (C) -- (P) -- (B);
\draw[indigo,dashed] (P) -- (N);
\draw[indigo,dashed] (B) -- (Q);
\pic[draw=madder,angle radius=4.2mm,angle eccentricity=1.45,"$45^\circ$"] {angle=N--B--M};
\pic[draw=madder,angle radius=4.2mm,angle eccentricity=1.45,"$45^\circ$"] {angle=P--B--N};
\pic[draw=madder,angle radius=4.2mm,angle eccentricity=1.45,"$90^\circ$"] {angle=B--Q--N};
\node[above left] at (A) {$A$};
\node[below left] at (B) {$B$};
\node[below right] at (C) {$C$};
\node[above right] at (D) {$D$};
\node[above] at (M) {$M$};
\node[right] at (N) {$N$};
\node[below right] at (P) {$P$};
\node[above right] at (Q) {$Q$};
\end{tikzpicture}}

\newcommand{\figThree}{%
\begin{tikzpicture}[scale=0.36,every node/.style={font=\small,inner sep=2pt},line width=.65pt]
\coordinate (A) at (0,9);
\coordinate (B) at (0,0);
\coordinate (C) at (9,0);
\coordinate (D) at (9,9);
\coordinate (M) at (3,9);
\coordinate (N) at (9,4.5);
\draw[inkiron] (A) -- (B) -- (C) -- (D) -- (A);
\draw[inkiron] (B) -- (M) -- (N);
\pic[draw=madder,angle radius=4.2mm,angle eccentricity=1.45,"$\alpha$"] {angle=B--M--N};
\pic[draw=madder,angle radius=4.2mm,angle eccentricity=1.45,"$\alpha$"] {angle=C--B--M};
\node[above left] at (A) {$A$};
\node[below left] at (B) {$B$};
\node[below right] at (C) {$C$};
\node[above right] at (D) {$D$};
\node[above] at (M) {$M$};
\node[right] at (N) {$N$};
\end{tikzpicture}}

\newcommand{\figFour}{%
\begin{tikzpicture}[scale=0.78,every node/.style={font=\small,inner sep=2pt},line width=.65pt]
\coordinate (A) at (0,4);
\coordinate (B) at (0,0);
\coordinate (C) at (4,0);
\coordinate (D) at (4,4);
\coordinate (E) at (1,3);
\coordinate (N) at (4,2);
\draw[inkiron] (A) -- (B) -- (C) -- (D) -- (A);
\draw[inkiron] (A) -- (C);
\draw[inkiron] (B) -- (E) -- (N);
\pic[draw=madder,angle radius=4.2mm,angle eccentricity=1.45,"$90^\circ$"] {angle=B--E--N};
\node[above left] at (A) {$A$};
\node[below left] at (B) {$B$};
\node[below right] at (C) {$C$};
\node[above right] at (D) {$D$};
\node[above right] at (E) {$E$};
\node[right] at (N) {$N$};
\end{tikzpicture}}

\newcommand{\figFourSolve}{%
\begin{tikzpicture}[scale=0.78,every node/.style={font=\small,inner sep=2pt},line width=.65pt]
\coordinate (A) at (0,4);
\coordinate (B) at (0,0);
\coordinate (C) at (4,0);
\coordinate (D) at (4,4);
\coordinate (E) at (1,3);
\coordinate (N) at (4,2);
\draw[inkiron] (A) -- (B) -- (C) -- (D) -- (A);
\draw[inkiron] (A) -- (C);
\draw[inkiron] (B) -- (E) -- (N);
\draw[indigo,dashed] (E) -- (D);
\pic[draw=madder,angle radius=4.2mm,angle eccentricity=1.45,"$90^\circ$"] {angle=B--E--N};
\node[above left] at (A) {$A$};
\node[below left] at (B) {$B$};
\node[below right] at (C) {$C$};
\node[above right] at (D) {$D$};
\node[above right] at (E) {$E$};
\node[right] at (N) {$N$};
\end{tikzpicture}}

\newcommand{\figFive}{%
\begin{tikzpicture}[scale=0.47,every node/.style={font=\small,inner sep=2pt},line width=.65pt]
\coordinate (A) at (0,6);
\coordinate (B) at (-3,0);
\coordinate (D) at (0,0);
\coordinate (C) at (2,0);
\draw[inkiron] (A) -- (B) -- (C) -- (A);
\draw[inkiron] (A) -- (D);
\pic[draw=madder,angle radius=4.2mm,angle eccentricity=1.45,"$45^\circ$"] {angle=B--A--C};
\pic[draw=madder,angle radius=4.2mm,angle eccentricity=1.45,"$90^\circ$"] {angle=C--D--A};
\node[above] at (A) {$A$};
\node[below left] at (B) {$B$};
\node[below] at (D) {$D$};
\node[below right] at (C) {$C$};
\end{tikzpicture}}

\newcommand{\figFiveSolve}{%
\begin{tikzpicture}[scale=0.4,every node/.style={font=\small,inner sep=2pt},line width=.65pt]
\coordinate (A) at (0,6);
\coordinate (B) at (-3,0);
\coordinate (D) at (0,0);
\coordinate (C) at (2,0);
\coordinate (E) at (-6,0);
\coordinate (F) at (6,0);
\coordinate (H) at (6,6);
\coordinate (G) at (6,3);
\draw[inkiron] (A) -- (B) -- (C) -- (A);
\draw[inkiron] (A) -- (D);
\draw[indigo,dashed] (E) -- (A) -- (F);
\draw[indigo,dashed] (E) -- (B);
\draw[indigo,dashed] (C) -- (F) -- (H) -- (A);
\draw[indigo,dashed] (A) -- (G) -- (C);
\pic[draw=madder,angle radius=4.2mm,angle eccentricity=1.45,"$45^\circ$"] {angle=B--A--C};
\node[above] at (A) {$A$};
\node[below left] at (B) {$B$};
\node[below] at (D) {$D$};
\node[below right] at (C) {$C$};
\node[below left] at (E) {$E$};
\node[below right] at (F) {$F$};
\node[above right] at (H) {$H$};
\node[right] at (G) {$G$};
\end{tikzpicture}}

\newcommand{\figFiveSquareSolve}{%
\begin{tikzpicture}[scale=0.42,every node/.style={font=\small,inner sep=2pt},line width=.65pt]
\coordinate (A) at (0,6);
\coordinate (B) at (-3,0);
\coordinate (D) at (0,0);
\coordinate (C) at (2,0);
\coordinate (M) at (-6,0);
\coordinate (N) at (-6,6);
\coordinate (G) at (-6,4);
\draw[inkiron] (A) -- (B) -- (C) -- (A);
\draw[inkiron] (A) -- (D);
\draw[indigo,dashed] (A) -- (N) -- (M) -- (B);
\draw[indigo,dashed] (A) -- (G) -- (B);
\pic[draw=madder,angle radius=4.2mm,angle eccentricity=1.45,"$45^\circ$"] {angle=B--A--C};
\pic[draw=madder,angle radius=4.2mm,angle eccentricity=1.45,"$45^\circ$"] {angle=G--A--B};
\node[above] at (A) {$A$};
\node[below left] at (B) {$B$};
\node[below] at (D) {$D$};
\node[below right] at (C) {$C$};
\node[below left] at (M) {$M$};
\node[above left] at (N) {$N$};
\node[left] at (G) {$G$};
\path (M) -- node[left,fill=white] {$6-x$} (G);
\path (G) -- node[left,fill=white] {$x$} (N);
\end{tikzpicture}}

\newcommand{\figFold}{%
\begin{tikzpicture}[scale=0.44,every node/.style={font=\small,inner sep=2pt},line width=.65pt]
\coordinate (A) at (0,6);
\coordinate (B) at (-3,0);
\coordinate (D) at (0,0);
\coordinate (C) at (2,0);
\coordinate (E) at (-4.8,2.4);
\coordinate (F) at (3.6,1.2);
\coordinate (P) at (-1.2,-2.4);
\draw[inkiron] (A) -- (B) -- (C) -- (A);
\draw[inkiron] (A) -- (D);
\draw[indigo,dashed] (A) -- (E) -- (P) -- (F) -- (A);
\pic[draw=madder,angle radius=4.2mm,angle eccentricity=1.45,"$90^\circ$"] {angle=E--A--F};
\pic[draw=madder,angle radius=4.2mm,angle eccentricity=1.45,"$90^\circ$"] {angle=F--P--E};
\node[above] at (A) {$A$};
\node[below left] at (B) {$B$};
\node[below] at (D) {$D$};
\node[below right] at (C) {$C$};
\node[left] at (E) {$E_1$};
\node[right] at (F) {$F_1$};
\node[below] at (P) {$P$};
\end{tikzpicture}}

\newcommand{\figSix}{%
\begin{tikzpicture}[scale=0.38,every node/.style={font=\small,inner sep=2pt},line width=.65pt]
\coordinate (A) at (0,6);
\coordinate (B) at (12,6);
\coordinate (C) at (12,0);
\coordinate (D) at (0,0);
\coordinate (F) at (3,0);
\coordinate (E) at (12,2);
\draw[inkiron] (A) -- (B) -- (C) -- (D) -- (A);
\draw[inkiron] (A) -- (E);
\draw[inkiron] (A) -- (F);
\pic[draw=madder,angle radius=4.2mm,angle eccentricity=1.45,"$45^\circ$"] {angle=F--A--E};
\node[above left] at (A) {$A$};
\node[above right] at (B) {$B$};
\node[below right] at (C) {$C$};
\node[below left] at (D) {$D$};
\node[below] at (F) {$F$};
\node[right] at (E) {$E$};
\end{tikzpicture}}

\newcommand{\figSixSolve}{%
\begin{tikzpicture}[scale=0.38,every node/.style={font=\small,inner sep=2pt},line width=.65pt]
\coordinate (A) at (0,6);
\coordinate (B) at (12,6);
\coordinate (C) at (12,0);
\coordinate (D) at (0,0);
\coordinate (F) at (3,0);
\coordinate (E) at (12,2);
\coordinate (M) at (6,6);
\coordinate (Q) at (6,0);
\coordinate (N) at (6,4);
\draw[inkiron] (A) -- (B) -- (C) -- (D) -- (A);
\draw[inkiron] (A) -- (E);
\draw[inkiron] (A) -- (F);
\draw[indigo,dashed] (M) -- (Q);
\draw[indigo,dashed] (F) -- (N);
\node[above left] at (A) {$A$};
\node[above right] at (B) {$B$};
\node[below right] at (C) {$C$};
\node[below left] at (D) {$D$};
\node[below] at (F) {$F$};
\node[right] at (E) {$E$};
\node[above] at (M) {$M$};
\node[below] at (Q) {$Q$};
\node[above right] at (N) {$N$};
\path (M) -- node[right,fill=white] {$x$} (N);
\path (N) -- node[right,fill=white] {$6-x$} (Q);
\end{tikzpicture}}

\newcommand{\figSeven}{%
\begin{tikzpicture}[scale=0.48,every node/.style={font=\small,inner sep=2pt},line width=.65pt]
\coordinate (A) at (0,6);
\coordinate (B) at (8,0);
\coordinate (C) at (0,0);
\coordinate (D) at (0,3);
\coordinate (E) at (2.72727273,0);
\coordinate (F) at (1.64383562,2.38356164);
\draw[inkiron] (A) -- (C) -- (B) -- (A);
\draw[inkiron] (D) -- (B);
\draw[inkiron] (A) -- (E);
\pic[draw=madder,angle radius=4.2mm,angle eccentricity=1.45,"$45^\circ$"] {angle=E--F--B};
\node[above] at (A) {$A$};
\node[below right] at (B) {$B$};
\node[below left] at (C) {$C$};
\node[left] at (D) {$D$};
\node[below] at (E) {$E$};
\node[above right] at (F) {$F$};
\end{tikzpicture}}

\newcommand{\figSevenSolve}{%
\begin{tikzpicture}[scale=0.43,every node/.style={font=\small,inner sep=2pt},line width=.65pt]
\coordinate (A) at (0,6);
\coordinate (B) at (8,0);
\coordinate (C) at (0,0);
\coordinate (D) at (0,3);
\coordinate (E) at (2.72727273,0);
\coordinate (F) at (1.64383562,2.38356164);
\coordinate (G) at (-2.25,0);
\draw[inkiron] (A) -- (C) -- (B) -- (A);
\draw[inkiron] (D) -- (B);
\draw[inkiron] (A) -- (E);
\draw[indigo,dashed] (A) -- (G) -- (C);
\pic[draw=madder,angle radius=4.2mm,angle eccentricity=1.45,"$45^\circ$"] {angle=G--A--E};
\node[above] at (A) {$A$};
\node[below right] at (B) {$B$};
\node[below left] at (C) {$C$};
\node[left] at (D) {$D$};
\node[below] at (E) {$E$};
\node[above right] at (F) {$F$};
\node[below left] at (G) {$G$};
\end{tikzpicture}}

\newcommand{\figEight}{%
\begin{tikzpicture}[scale=0.42,every node/.style={font=\small,inner sep=2pt},line width=.65pt]
\coordinate (A) at (0,8);
\coordinate (B) at (0,0);
\coordinate (C) at (9,0);
\coordinate (D) at (9,8);
\coordinate (E) at (4,0);
\coordinate (F) at (9,5);
\draw[inkiron] (A) -- (B) -- (C) -- (D) -- (A);
\draw[inkiron] (A) -- (E) -- (F) -- (A);
\pic[draw=madder,angle radius=4.2mm,angle eccentricity=1.45,"$45^\circ$"] {angle=E--A--F};
\node[above left] at (A) {$A$};
\node[below left] at (B) {$B$};
\node[below right] at (C) {$C$};
\node[above right] at (D) {$D$};
\node[below] at (E) {$E$};
\node[right] at (F) {$F$};
\path (E) -- node[below,fill=white] {$x$} (C);
\path (C) -- node[right,fill=white] {$x$} (F);
\end{tikzpicture}}

\newcommand{\figEightSolve}{%
\begin{tikzpicture}[scale=0.32,every node/.style={font=\small,inner sep=2pt},line width=.65pt]
\coordinate (A) at (0,8);
\coordinate (B) at (0,0);
\coordinate (C) at (9,0);
\coordinate (D) at (9,8);
\coordinate (E) at (4,0);
\coordinate (F) at (9,5);
\coordinate (G) at (8,12);
\coordinate (H) at (9,12);
\draw[inkiron] (A) -- (B) -- (C) -- (D) -- (A);
\draw[inkiron] (A) -- (E) -- (F) -- (A);
\draw[indigo,dashed] (A) -- (G) -- (H) -- (F);
\draw[indigo,dashed] (F) -- (G);
\pic[draw=madder,angle radius=4.2mm,angle eccentricity=1.45,"$45^\circ$"] {angle=E--A--F};
\pic[draw=madder,angle radius=4.2mm,angle eccentricity=1.45,"$45^\circ$"] {angle=F--A--G};
\node[above left] at (A) {$A$};
\node[below left] at (B) {$B$};
\node[below right] at (C) {$C$};
\node[above right] at (D) {$D$};
\node[below] at (E) {$E$};
\node[right] at (F) {$F$};
\node[above left] at (G) {$G$};
\node[above right] at (H) {$H$};
\end{tikzpicture}}

\newcommand{\figNineIn}{%
\begin{tikzpicture}[scale=0.68,every node/.style={font=\small,inner sep=2pt},line width=.65pt]
\coordinate (A) at (0,0);
\coordinate (B) at (3,0);
\coordinate (C) at (3,3);
\coordinate (D) at (0,3);
\coordinate (M) at (3,1);
\coordinate (N) at (1.5,3);
\draw[inkiron] (A) -- (B) -- (C) -- (D) -- (A);
\draw[inkiron] (A) -- (M) -- (N) -- (A);
\pic[draw=madder,angle radius=4.2mm,angle eccentricity=1.45,"$45^\circ$"] {angle=M--A--N};
\node[below left] at (A) {$A$};
\node[below right] at (B) {$B$};
\node[above right] at (C) {$C$};
\node[above left] at (D) {$D$};
\node[right] at (M) {$M$};
\node[above] at (N) {$N$};
\node[font=\small] at (current bounding box.south) [below=6pt] {（1）内部构型};
\end{tikzpicture}}

\newcommand{\figNineOut}{%
\begin{tikzpicture}[scale=0.43,every node/.style={font=\small,inner sep=2pt},line width=.65pt]
\coordinate (A) at (0,0);
\coordinate (B) at (3,0);
\coordinate (C) at (3,3);
\coordinate (D) at (0,3);
\coordinate (M) at (3,6);
\coordinate (N) at (-1,3);
\draw[inkiron] (A) -- (B) -- (C) -- (D) -- (A);
\draw[inkiron] (A) -- (M) -- (N) -- (A);
\draw[inkiron] (C) -- (M);
\draw[inkiron] (D) -- (N);
\pic[draw=madder,angle radius=4.2mm,angle eccentricity=1.45,"$45^\circ$"] {angle=M--A--N};
\node[below left] at (A) {$A$};
\node[below right] at (B) {$B$};
\node[above right] at (C) {$C$};
\node[above left] at (D) {$D$};
\node[right] at (M) {$M$};
\node[above left] at (N) {$N$};
\node[font=\small] at (current bounding box.south) [below=6pt] {（2）外部构型};
\end{tikzpicture}}

\newcommand{\figNine}{\figNineIn\hspace{1.2cm}\figNineOut}

\newcommand{\figNineInSolve}{%
\begin{tikzpicture}[scale=0.62,every node/.style={font=\small,inner sep=2pt},line width=.65pt]
\coordinate (A) at (0,0);
\coordinate (B) at (3,0);
\coordinate (C) at (3,3);
\coordinate (D) at (0,3);
\coordinate (M) at (3,1);
\coordinate (N) at (1.5,3);
\coordinate (P) at (-1,3);
\draw[inkiron] (A) -- (B) -- (C) -- (D) -- (A);
\draw[inkiron] (A) -- (M) -- (N) -- (A);
\draw[indigo,dashed] (A) -- (P) -- (D);
\pic[draw=madder,angle radius=4.2mm,angle eccentricity=1.45,"$45^\circ$"] {angle=M--A--N};
\node[below left] at (A) {$A$};
\node[below right] at (B) {$B$};
\node[above right] at (C) {$C$};
\node[above left] at (D) {$D$};
\node[right] at (M) {$M$};
\node[above] at (N) {$N$};
\node[above left] at (P) {$P$};
\node[font=\small] at (current bounding box.south) [below=6pt] {PN＝PD＋DN};
\end{tikzpicture}}

\newcommand{\figNineOutSolve}{%
\begin{tikzpicture}[scale=0.35,every node/.style={font=\small,inner sep=2pt},line width=.65pt]
\coordinate (A) at (0,0);
\coordinate (B) at (3,0);
\coordinate (C) at (3,3);
\coordinate (D) at (0,3);
\coordinate (M) at (3,6);
\coordinate (N) at (-1,3);
\coordinate (P) at (-6,3);
\draw[inkiron] (A) -- (B) -- (C) -- (D) -- (A);
\draw[inkiron] (A) -- (M) -- (N) -- (A);
\draw[inkiron] (C) -- (M);
\draw[inkiron] (D) -- (N);
\draw[indigo,dashed] (A) -- (P) -- (N);
\pic[draw=madder,angle radius=4.2mm,angle eccentricity=1.45,"$45^\circ$"] {angle=M--A--N};
\node[below left] at (A) {$A$};
\node[below right] at (B) {$B$};
\node[above right] at (C) {$C$};
\node[above left] at (D) {$D$};
\node[right] at (M) {$M$};
\node[above] at (N) {$N$};
\node[above left] at (P) {$P$};
\node[font=\small] at (current bounding box.south) [below=6pt] {PN＝PD−DN};
\end{tikzpicture}}

\newcommand{\figNineSolve}{\figNineInSolve\hspace{.7cm}\figNineOutSolve}

\newcommand{\figTen}{%
\begin{tikzpicture}[scale=0.48,every node/.style={font=\small,inner sep=2pt},line width=.65pt]
\coordinate (B) at (0,0);
\coordinate (D) at (4,0);
\coordinate (C) at (9,0);
\coordinate (A) at (6,6);
\draw[inkiron] (A) -- (B) -- (C) -- (A);
\draw[inkiron] (A) -- (D);
\pic[draw=madder,angle radius=4.2mm,angle eccentricity=1.45,"$45^\circ$"] {angle=D--A--C};
\pic[draw=madder,angle radius=4.2mm,angle eccentricity=1.45,"$45^\circ$"] {angle=C--B--A};
\node[below left] at (B) {$B$};
\node[below] at (D) {$D$};
\node[below right] at (C) {$C$};
\node[above] at (A) {$A$};
\end{tikzpicture}}

\newcommand{\figTenSmall}{%
\begin{tikzpicture}[scale=0.37,every node/.style={font=\small,inner sep=2pt},line width=.65pt]
\coordinate (B) at (0,0);
\coordinate (D) at (4,0);
\coordinate (C) at (9,0);
\coordinate (A) at (3,3);
\coordinate (H) at (3,0);
\coordinate (K) at (6,0);
\coordinate (G) at (6,4);
\draw[inkiron] (A) -- (B) -- (C) -- (A);
\draw[inkiron] (A) -- (D);
\draw[indigo,dashed] (A) -- (H);
\draw[indigo,dashed] (A) -- (G) -- (C);
\draw[indigo,dashed] (A) -- (K) -- (G);
\pic[draw=madder,angle radius=4.2mm,angle eccentricity=1.45,"$45^\circ$"] {angle=D--A--C};
\node[below left] at (B) {$B$};
\node[below] at (D) {$D$};
\node[below right] at (C) {$C$};
\node[above] at (A) {$A$};
\node[below] at (H) {$H$};
\node[below] at (K) {$K$};
\node[above] at (G) {$G$};
\node[font=\small] at (current bounding box.south) [below=6pt] {构型一 H 在 BD 内部};
\end{tikzpicture}}

\newcommand{\figTenLarge}{%
\begin{tikzpicture}[scale=0.37,every node/.style={font=\small,inner sep=2pt},line width=.65pt]
\coordinate (B) at (0,0);
\coordinate (D) at (4,0);
\coordinate (C) at (9,0);
\coordinate (A) at (6,6);
\coordinate (H) at (6,0);
\coordinate (K) at (12,0);
\coordinate (G) at (12,4);
\draw[inkiron] (A) -- (B) -- (C) -- (A);
\draw[inkiron] (A) -- (D);
\draw[indigo,dashed] (A) -- (H);
\draw[indigo,dashed] (A) -- (G) -- (C);
\draw[indigo,dashed] (A) -- (K) -- (G);
\pic[draw=madder,angle radius=4.2mm,angle eccentricity=1.45,"$45^\circ$"] {angle=D--A--C};
\node[below left] at (B) {$B$};
\node[below] at (D) {$D$};
\node[below right] at (C) {$C$};
\node[above] at (A) {$A$};
\node[below] at (H) {$H$};
\node[below] at (K) {$K$};
\node[above] at (G) {$G$};
\node[font=\small] at (current bounding box.south) [below=6pt] {构型二 H 在 DC 内部};
\end{tikzpicture}}

\newcommand{\figTenSolve}{\figTenSmall\hspace{.8cm}\figTenLarge}

% 重绘手拉手模型；辅助线为靛蓝虚线，数值长度留在题干中。
\newcommand{\hhOne}{%
\begin{tikzpicture}[scale=0.66,line width=.65pt,every node/.style={font=\small,inner sep=2pt}]
\coordinate (A) at (0,4);
\coordinate (B) at (0,0);
\coordinate (C) at (4,0);
\coordinate (D) at (-1,1);
\coordinate (E) at (3,3);
\coordinate (F) at (1,-1);
\draw[inkiron] (A) -- (B);
\draw[inkiron] (B) -- (C);
\draw[inkiron] (A) -- (C);
\draw[inkiron] (A) -- (D);
\draw[inkiron] (A) -- (E);
\draw[inkiron] (D) -- (E);
\draw[inkiron] (D) -- (F);
\draw[inkiron] (C) -- (F);
\draw[inkiron] (C) -- (E);
\pic[draw=selvage,angle radius=2.4mm] {right angle=D--A--E};
\pic[draw=selvage,angle radius=2.4mm] {right angle=A--B--C};
\draw[inkiron] ($(A)!0.5!(D)!3pt!90:(D)$) -- ($(A)!0.5!(D)!3pt!-90:(D)$);
\draw[inkiron] ($(A)!0.5!(E)!3pt!90:(E)$) -- ($(A)!0.5!(E)!3pt!-90:(E)$);
\draw[inkiron] ($(A)!0.4825!(B)!3pt!90:(B)$) -- ($(A)!0.4825!(B)!3pt!-90:(B)$);
\draw[inkiron] ($(A)!0.5175!(B)!3pt!90:(B)$) -- ($(A)!0.5175!(B)!3pt!-90:(B)$);
\draw[inkiron] ($(B)!0.4825!(C)!3pt!90:(C)$) -- ($(B)!0.4825!(C)!3pt!-90:(C)$);
\draw[inkiron] ($(B)!0.5175!(C)!3pt!90:(C)$) -- ($(B)!0.5175!(C)!3pt!-90:(C)$);
\draw[inkiron] ($(D)!0.46499999999999997!(B)!3pt!90:(B)$) -- ($(D)!0.46499999999999997!(B)!3pt!-90:(B)$);
\draw[inkiron] ($(D)!0.5!(B)!3pt!90:(B)$) -- ($(D)!0.5!(B)!3pt!-90:(B)$);
\draw[inkiron] ($(D)!0.535!(B)!3pt!90:(B)$) -- ($(D)!0.535!(B)!3pt!-90:(B)$);
\draw[inkiron] ($(B)!0.46499999999999997!(F)!3pt!90:(F)$) -- ($(B)!0.46499999999999997!(F)!3pt!-90:(F)$);
\draw[inkiron] ($(B)!0.5!(F)!3pt!90:(F)$) -- ($(B)!0.5!(F)!3pt!-90:(F)$);
\draw[inkiron] ($(B)!0.535!(F)!3pt!90:(F)$) -- ($(B)!0.535!(F)!3pt!-90:(F)$);
\node[above] at (A) {$A$};
\node[left] at (B) {$B$};
\node[right] at (C) {$C$};
\node[left] at (D) {$D$};
\node[right] at (E) {$E$};
\node[below] at (F) {$F$};
\end{tikzpicture}}

\newcommand{\hhOneExtend}{%
\begin{tikzpicture}[scale=0.5,line width=.65pt,every node/.style={font=\small,inner sep=2pt}]
\coordinate (A) at (0,4);
\coordinate (B) at (0,0);
\coordinate (C) at (4,0);
\coordinate (D) at (-1,1);
\coordinate (E) at (3,3);
\coordinate (F) at (1,-1);
\coordinate (G) at (0,-4);
\draw[inkiron] (A) -- (B);
\draw[inkiron] (B) -- (C);
\draw[inkiron] (A) -- (C);
\draw[inkiron] (A) -- (D);
\draw[inkiron] (A) -- (E);
\draw[inkiron] (D) -- (E);
\draw[inkiron] (D) -- (F);
\draw[inkiron] (C) -- (F);
\draw[inkiron] (C) -- (E);
\draw[indigo,dashed] (B) -- (G);
\draw[indigo,dashed] (G) -- (F);
\draw[indigo,dashed] (G) -- (C);
\pic[draw=selvage,angle radius=2.4mm] {right angle=D--A--E};
\pic[draw=selvage,angle radius=2.4mm] {right angle=A--B--C};
\pic[draw=selvage,angle radius=2.4mm] {right angle=A--C--G};
\draw[inkiron] ($(A)!0.5!(D)!3pt!90:(D)$) -- ($(A)!0.5!(D)!3pt!-90:(D)$);
\draw[inkiron] ($(A)!0.5!(E)!3pt!90:(E)$) -- ($(A)!0.5!(E)!3pt!-90:(E)$);
\draw[inkiron] ($(A)!0.4825!(B)!3pt!90:(B)$) -- ($(A)!0.4825!(B)!3pt!-90:(B)$);
\draw[inkiron] ($(A)!0.5175!(B)!3pt!90:(B)$) -- ($(A)!0.5175!(B)!3pt!-90:(B)$);
\draw[inkiron] ($(B)!0.4825!(C)!3pt!90:(C)$) -- ($(B)!0.4825!(C)!3pt!-90:(C)$);
\draw[inkiron] ($(B)!0.5175!(C)!3pt!90:(C)$) -- ($(B)!0.5175!(C)!3pt!-90:(C)$);
\draw[inkiron] ($(D)!0.46499999999999997!(B)!3pt!90:(B)$) -- ($(D)!0.46499999999999997!(B)!3pt!-90:(B)$);
\draw[inkiron] ($(D)!0.5!(B)!3pt!90:(B)$) -- ($(D)!0.5!(B)!3pt!-90:(B)$);
\draw[inkiron] ($(D)!0.535!(B)!3pt!90:(B)$) -- ($(D)!0.535!(B)!3pt!-90:(B)$);
\draw[inkiron] ($(B)!0.46499999999999997!(F)!3pt!90:(F)$) -- ($(B)!0.46499999999999997!(F)!3pt!-90:(F)$);
\draw[inkiron] ($(B)!0.5!(F)!3pt!90:(F)$) -- ($(B)!0.5!(F)!3pt!-90:(F)$);
\draw[inkiron] ($(B)!0.535!(F)!3pt!90:(F)$) -- ($(B)!0.535!(F)!3pt!-90:(F)$);
\node[above] at (A) {$A$};
\node[left] at (B) {$B$};
\node[right] at (C) {$C$};
\node[left] at (D) {$D$};
\node[right] at (E) {$E$};
\node[below] at (F) {$F$};
\node[below] at (G) {$G$};
\end{tikzpicture}}

\newcommand{\hhOneRotate}{%
\begin{tikzpicture}[scale=0.5,line width=.65pt,every node/.style={font=\small,inner sep=2pt}]
\coordinate (A) at (0,4);
\coordinate (B) at (0,0);
\coordinate (C) at (4,0);
\coordinate (D) at (-1,1);
\coordinate (E) at (3,3);
\coordinate (F) at (1,-1);
\coordinate (G) at (-4,0);
\draw[inkiron] (A) -- (B);
\draw[inkiron] (B) -- (C);
\draw[inkiron] (A) -- (C);
\draw[inkiron] (A) -- (D);
\draw[inkiron] (A) -- (E);
\draw[inkiron] (D) -- (E);
\draw[inkiron] (D) -- (F);
\draw[inkiron] (C) -- (F);
\draw[inkiron] (C) -- (E);
\draw[indigo,dashed] (B) -- (G);
\draw[indigo,dashed] (G) -- (A);
\draw[indigo,dashed] (G) -- (D);
\pic[draw=selvage,angle radius=2.4mm] {right angle=D--A--E};
\pic[draw=selvage,angle radius=2.4mm] {right angle=A--B--C};
\pic[draw=selvage,angle radius=2.4mm] {right angle=G--A--C};
\draw[inkiron] ($(A)!0.5!(D)!3pt!90:(D)$) -- ($(A)!0.5!(D)!3pt!-90:(D)$);
\draw[inkiron] ($(A)!0.5!(E)!3pt!90:(E)$) -- ($(A)!0.5!(E)!3pt!-90:(E)$);
\draw[inkiron] ($(A)!0.4825!(B)!3pt!90:(B)$) -- ($(A)!0.4825!(B)!3pt!-90:(B)$);
\draw[inkiron] ($(A)!0.5175!(B)!3pt!90:(B)$) -- ($(A)!0.5175!(B)!3pt!-90:(B)$);
\draw[inkiron] ($(B)!0.4825!(C)!3pt!90:(C)$) -- ($(B)!0.4825!(C)!3pt!-90:(C)$);
\draw[inkiron] ($(B)!0.5175!(C)!3pt!90:(C)$) -- ($(B)!0.5175!(C)!3pt!-90:(C)$);
\draw[inkiron] ($(D)!0.46499999999999997!(B)!3pt!90:(B)$) -- ($(D)!0.46499999999999997!(B)!3pt!-90:(B)$);
\draw[inkiron] ($(D)!0.5!(B)!3pt!90:(B)$) -- ($(D)!0.5!(B)!3pt!-90:(B)$);
\draw[inkiron] ($(D)!0.535!(B)!3pt!90:(B)$) -- ($(D)!0.535!(B)!3pt!-90:(B)$);
\draw[inkiron] ($(B)!0.46499999999999997!(F)!3pt!90:(F)$) -- ($(B)!0.46499999999999997!(F)!3pt!-90:(F)$);
\draw[inkiron] ($(B)!0.5!(F)!3pt!90:(F)$) -- ($(B)!0.5!(F)!3pt!-90:(F)$);
\draw[inkiron] ($(B)!0.535!(F)!3pt!90:(F)$) -- ($(B)!0.535!(F)!3pt!-90:(F)$);
\node[above] at (A) {$A$};
\node[left] at (B) {$B$};
\node[right] at (C) {$C$};
\node[left] at (D) {$D$};
\node[right] at (E) {$E$};
\node[below] at (F) {$F$};
\node[left] at (G) {$G$};
\end{tikzpicture}}

\newcommand{\hhTwo}{%
\begin{tikzpicture}[scale=0.59,line width=.65pt,every node/.style={font=\small,inner sep=2pt}]
\coordinate (B) at (0,0);
\coordinate (A) at (3,3);
\coordinate (C) at (6,0);
\coordinate (D) at (4.8,4.8);
\coordinate (E) at (0,3.6);
\draw[inkiron] (B) -- (D);
\draw[inkiron] (A) -- (C);
\draw[inkiron] (B) -- (C);
\draw[inkiron] (D) -- (C);
\draw[inkiron] (D) -- (E);
\draw[inkiron] (C) -- (E);
\draw[inkiron] (B) -- (E);
\pic[draw=selvage,angle radius=2.4mm] {right angle=B--A--C};
\pic[draw=selvage,angle radius=2.4mm] {right angle=C--D--E};
\draw[inkiron] ($(A)!0.5!(B)!3pt!90:(B)$) -- ($(A)!0.5!(B)!3pt!-90:(B)$);
\draw[inkiron] ($(A)!0.5!(C)!3pt!90:(C)$) -- ($(A)!0.5!(C)!3pt!-90:(C)$);
\draw[inkiron] ($(D)!0.4825!(C)!3pt!90:(C)$) -- ($(D)!0.4825!(C)!3pt!-90:(C)$);
\draw[inkiron] ($(D)!0.5175!(C)!3pt!90:(C)$) -- ($(D)!0.5175!(C)!3pt!-90:(C)$);
\draw[inkiron] ($(D)!0.4825!(E)!3pt!90:(E)$) -- ($(D)!0.4825!(E)!3pt!-90:(E)$);
\draw[inkiron] ($(D)!0.5175!(E)!3pt!90:(E)$) -- ($(D)!0.5175!(E)!3pt!-90:(E)$);
\node[below left] at (B) {$B$};
\node[below right] at (A) {$A$};
\node[below right] at (C) {$C$};
\node[above] at (D) {$D$};
\node[left] at (E) {$E$};
\end{tikzpicture}}

\newcommand{\hhTwoHands}{%
\begin{tikzpicture}[scale=0.46,line width=.65pt,every node/.style={font=\small,inner sep=2pt}]
\coordinate (B) at (0,0);
\coordinate (A) at (3,3);
\coordinate (C) at (6,0);
\coordinate (D) at (4.8,4.8);
\coordinate (E) at (0,3.6);
\coordinate (G) at (6,6);
\coordinate (F) at (9.6,6);
\draw[inkiron] (B) -- (D);
\draw[inkiron] (A) -- (C);
\draw[inkiron] (B) -- (C);
\draw[inkiron] (D) -- (C);
\draw[inkiron] (D) -- (E);
\draw[inkiron] (C) -- (E);
\draw[inkiron] (B) -- (E);
\draw[indigo,dashed] (D) -- (F);
\draw[indigo,dashed] (D) -- (G);
\draw[indigo,dashed] (C) -- (F);
\draw[indigo,dashed] (C) -- (G);
\draw[indigo,dashed] (G) -- (F);
\pic[draw=selvage,angle radius=2.4mm] {right angle=B--A--C};
\pic[draw=selvage,angle radius=2.4mm] {right angle=C--D--E};
\pic[draw=selvage,angle radius=2.4mm] {right angle=B--C--G};
\pic[draw=selvage,angle radius=2.4mm] {right angle=E--C--F};
\draw[inkiron] ($(A)!0.5!(B)!3pt!90:(B)$) -- ($(A)!0.5!(B)!3pt!-90:(B)$);
\draw[inkiron] ($(A)!0.5!(C)!3pt!90:(C)$) -- ($(A)!0.5!(C)!3pt!-90:(C)$);
\draw[inkiron] ($(D)!0.4825!(C)!3pt!90:(C)$) -- ($(D)!0.4825!(C)!3pt!-90:(C)$);
\draw[inkiron] ($(D)!0.5175!(C)!3pt!90:(C)$) -- ($(D)!0.5175!(C)!3pt!-90:(C)$);
\draw[inkiron] ($(D)!0.4825!(E)!3pt!90:(E)$) -- ($(D)!0.4825!(E)!3pt!-90:(E)$);
\draw[inkiron] ($(D)!0.5175!(E)!3pt!90:(E)$) -- ($(D)!0.5175!(E)!3pt!-90:(E)$);
\node[below left] at (B) {$B$};
\node[below right] at (A) {$A$};
\node[below right] at (C) {$C$};
\node[above] at (D) {$D$};
\node[left] at (E) {$E$};
\node[above] at (G) {$G$};
\node[right] at (F) {$F$};
\end{tikzpicture}}

\newcommand{\hhTwoParallel}{%
\begin{tikzpicture}[scale=0.46,line width=.65pt,every node/.style={font=\small,inner sep=2pt}]
\coordinate (B) at (0,0);
\coordinate (A) at (3,3);
\coordinate (C) at (6,0);
\coordinate (D) at (4.8,4.8);
\coordinate (E) at (0,3.6);
\coordinate (G) at (6,6);
\coordinate (F) at (9.6,6);
\coordinate (P) at (3.6,3.6);
\draw[inkiron] (B) -- (D);
\draw[inkiron] (A) -- (C);
\draw[inkiron] (B) -- (C);
\draw[inkiron] (D) -- (C);
\draw[inkiron] (D) -- (E);
\draw[inkiron] (C) -- (E);
\draw[inkiron] (B) -- (E);
\draw[indigo,dashed] (D) -- (F);
\draw[indigo,dashed] (D) -- (G);
\draw[indigo,dashed] (C) -- (F);
\draw[indigo,dashed] (C) -- (G);
\draw[indigo,dashed] (G) -- (F);
\draw[indigo,dashed] (E) -- (P);
\pic[draw=selvage,angle radius=2.4mm] {right angle=B--A--C};
\pic[draw=selvage,angle radius=2.4mm] {right angle=C--D--E};
\draw[inkiron] ($(A)!0.5!(B)!3pt!90:(B)$) -- ($(A)!0.5!(B)!3pt!-90:(B)$);
\draw[inkiron] ($(A)!0.5!(C)!3pt!90:(C)$) -- ($(A)!0.5!(C)!3pt!-90:(C)$);
\draw[inkiron] ($(D)!0.4825!(C)!3pt!90:(C)$) -- ($(D)!0.4825!(C)!3pt!-90:(C)$);
\draw[inkiron] ($(D)!0.5175!(C)!3pt!90:(C)$) -- ($(D)!0.5175!(C)!3pt!-90:(C)$);
\draw[inkiron] ($(D)!0.4825!(E)!3pt!90:(E)$) -- ($(D)!0.4825!(E)!3pt!-90:(E)$);
\draw[inkiron] ($(D)!0.5175!(E)!3pt!90:(E)$) -- ($(D)!0.5175!(E)!3pt!-90:(E)$);
\node[below left] at (B) {$B$};
\node[below right] at (A) {$A$};
\node[below right] at (C) {$C$};
\node[above] at (D) {$D$};
\node[left] at (E) {$E$};
\node[above] at (G) {$G$};
\node[right] at (F) {$F$};
\node[below right] at (P) {$P$};
\end{tikzpicture}}

% Same coordinates as the SVG/PNG diagrams.
\newcommand{\diReview}{%
\begin{tikzpicture}[scale=1.2,line width=.7pt,every node/.style={font=\small,inner sep=2pt}]
\coordinate (A) at (-1.00000000,0.00000000);
\coordinate (B) at (0.00000000,0.00000000);
\coordinate (C) at (2.00000000,0.00000000);
\coordinate (D) at (-1.00000000,2.00000000);
\coordinate (E) at (2.00000000,1.00000000);
\draw[inkiron] (A) -- (C);
\draw[inkiron] (A) -- (D);
\draw[inkiron] (B) -- (D);
\draw[inkiron] (B) -- (E);
\draw[inkiron] (C) -- (E);
\pic[draw=selvage,angle radius=2.2mm] {right angle=D--A--B};
\pic[draw=selvage,angle radius=2.2mm] {right angle=D--B--E};
\pic[draw=selvage,angle radius=2.2mm] {right angle=B--C--E};
\draw[inkiron] ($(B)!0.50000000!(D)!3pt!90:(D)$) -- ($(B)!0.50000000!(D)!3pt!-90:(D)$);
\draw[inkiron] ($(B)!0.50000000!(E)!3pt!90:(E)$) -- ($(B)!0.50000000!(E)!3pt!-90:(E)$);
\fill[inkiron] (A) circle (1.05pt);
\node at (-1.15000000,-0.27000000) {$A$};
\fill[inkiron] (B) circle (1.05pt);
\node at (0.00000000,-0.30000000) {$B$};
\fill[inkiron] (C) circle (1.05pt);
\node at (2.10000000,-0.27000000) {$C$};
\fill[inkiron] (D) circle (1.05pt);
\node at (-1.10000000,2.26000000) {$D$};
\fill[inkiron] (E) circle (1.05pt);
\node at (2.24000000,1.05000000) {$E$};
\end{tikzpicture}}

\newcommand{\diModel}{%
\begin{tikzpicture}[scale=0.47,line width=.7pt,every node/.style={font=\small,inner sep=2pt}]
\coordinate (B) at (0.00000000,0.00000000);
\coordinate (C) at (-6.00000000,2.00000000);
\coordinate (Cp) at (-2.00000000,-6.00000000);
\coordinate (D) at (3.00000000,1.00000000);
\coordinate (Dp) at (-1.00000000,3.00000000);
\fill[selvage!10] (C) -- (B) -- (Dp) -- cycle;
\fill[madder!10] (Cp) -- (B) -- (D) -- cycle;
\draw[inkiron] (B) -- (C);
\draw[inkiron] (B) -- (Cp);
\draw[inkiron] (B) -- (D);
\draw[inkiron] (B) -- (Dp);
\draw[inkiron] (C) -- (Dp);
\draw[inkiron] (Cp) -- (D);
\pic[draw=selvage,angle radius=2.2mm] {right angle=C--B--Cp};
\pic[draw=selvage,angle radius=2.2mm] {right angle=D--B--Dp};
\draw[inkiron] ($(B)!0.50000000!(C)!3pt!90:(C)$) -- ($(B)!0.50000000!(C)!3pt!-90:(C)$);
\draw[inkiron] ($(B)!0.50000000!(Cp)!3pt!90:(Cp)$) -- ($(B)!0.50000000!(Cp)!3pt!-90:(Cp)$);
\draw[inkiron] ($(B)!0.48260747!(D)!3pt!90:(D)$) -- ($(B)!0.48260747!(D)!3pt!-90:(D)$);
\draw[inkiron] ($(B)!0.51739253!(D)!3pt!90:(D)$) -- ($(B)!0.51739253!(D)!3pt!-90:(D)$);
\draw[inkiron] ($(B)!0.48260747!(Dp)!3pt!90:(Dp)$) -- ($(B)!0.48260747!(Dp)!3pt!-90:(Dp)$);
\draw[inkiron] ($(B)!0.51739253!(Dp)!3pt!90:(Dp)$) -- ($(B)!0.51739253!(Dp)!3pt!-90:(Dp)$);
\fill[inkiron] (B) circle (1.05pt);
\node at (-0.46000000,-0.45000000) {$B$};
\fill[inkiron] (C) circle (1.05pt);
\node at (-6.45000000,2.10000000) {$C$};
\fill[inkiron] (Cp) circle (1.05pt);
\node at (-2.20000000,-6.42000000) {$C'$};
\fill[inkiron] (D) circle (1.05pt);
\node at (3.36000000,1.12000000) {$D$};
\fill[inkiron] (Dp) circle (1.05pt);
\node at (-0.68000000,3.26000000) {$D'$};
\end{tikzpicture}}

\newcommand{\diArea}{%
\begin{tikzpicture}[scale=0.47,line width=.7pt,every node/.style={font=\small,inner sep=2pt}]
\coordinate (B) at (0.00000000,0.00000000);
\coordinate (C) at (-6.00000000,2.00000000);
\coordinate (Cp) at (-2.00000000,-6.00000000);
\coordinate (D) at (3.00000000,1.00000000);
\coordinate (Dp) at (-1.00000000,3.00000000);
\fill[selvage!10] (C) -- (B) -- (Dp) -- cycle;
\fill[madder!10] (Cp) -- (B) -- (D) -- cycle;
\draw[inkiron] (B) -- (C);
\draw[inkiron] (B) -- (Cp);
\draw[inkiron] (B) -- (D);
\draw[inkiron] (B) -- (Dp);
\draw[inkiron] (C) -- (Dp);
\draw[inkiron] (Cp) -- (D);
\pic[draw=selvage,angle radius=2.2mm] {right angle=C--B--Cp};
\pic[draw=selvage,angle radius=2.2mm] {right angle=D--B--Dp};
\draw[inkiron] ($(B)!0.50000000!(C)!3pt!90:(C)$) -- ($(B)!0.50000000!(C)!3pt!-90:(C)$);
\draw[inkiron] ($(B)!0.50000000!(Cp)!3pt!90:(Cp)$) -- ($(B)!0.50000000!(Cp)!3pt!-90:(Cp)$);
\draw[inkiron] ($(B)!0.48260747!(D)!3pt!90:(D)$) -- ($(B)!0.48260747!(D)!3pt!-90:(D)$);
\draw[inkiron] ($(B)!0.51739253!(D)!3pt!90:(D)$) -- ($(B)!0.51739253!(D)!3pt!-90:(D)$);
\draw[inkiron] ($(B)!0.48260747!(Dp)!3pt!90:(Dp)$) -- ($(B)!0.48260747!(Dp)!3pt!-90:(Dp)$);
\draw[inkiron] ($(B)!0.51739253!(Dp)!3pt!90:(Dp)$) -- ($(B)!0.51739253!(Dp)!3pt!-90:(Dp)$);
\coordinate (M) at (-1.80000000,0.60000000);
\coordinate (N) at (0.60000000,1.80000000);
\draw[indigo,dashed] (Dp) -- (M);
\draw[indigo,dashed] (D) -- (N);
\draw[indigo,dashed] (B) -- (N);
\pic[draw=selvage,angle radius=2.2mm] {right angle=B--M--Dp};
\pic[draw=selvage,angle radius=2.2mm] {right angle=B--N--D};
\fill[inkiron] (B) circle (1.05pt);
\node at (-0.46000000,-0.45000000) {$B$};
\fill[inkiron] (C) circle (1.05pt);
\node at (-6.45000000,2.10000000) {$C$};
\fill[inkiron] (Cp) circle (1.05pt);
\node at (-2.20000000,-6.42000000) {$C'$};
\fill[inkiron] (D) circle (1.05pt);
\node at (3.36000000,1.12000000) {$D$};
\fill[inkiron] (Dp) circle (1.05pt);
\node at (-0.68000000,3.26000000) {$D'$};
\fill[inkiron] (M) circle (1.05pt);
\node at (-2.13000000,0.38000000) {$M$};
\fill[inkiron] (N) circle (1.05pt);
\node at (0.64000000,2.16000000) {$N$};
\end{tikzpicture}}

\newcommand{\diMidpoint}{%
\begin{tikzpicture}[scale=0.45,line width=.7pt,every node/.style={font=\small,inner sep=2pt}]
\coordinate (B) at (0.00000000,0.00000000);
\coordinate (C) at (-6.00000000,2.00000000);
\coordinate (Cp) at (-2.00000000,-6.00000000);
\coordinate (D) at (3.00000000,1.00000000);
\coordinate (Dp) at (-1.00000000,3.00000000);
\fill[selvage!10] (C) -- (B) -- (Dp) -- cycle;
\fill[madder!10] (Cp) -- (B) -- (D) -- cycle;
\draw[inkiron] (B) -- (C);
\draw[inkiron] (B) -- (Cp);
\draw[inkiron] (B) -- (D);
\draw[inkiron] (B) -- (Dp);
\draw[inkiron] (C) -- (Dp);
\draw[inkiron] (Cp) -- (D);
\pic[draw=selvage,angle radius=2.2mm] {right angle=C--B--Cp};
\pic[draw=selvage,angle radius=2.2mm] {right angle=D--B--Dp};
\draw[inkiron] ($(B)!0.50000000!(C)!3pt!90:(C)$) -- ($(B)!0.50000000!(C)!3pt!-90:(C)$);
\draw[inkiron] ($(B)!0.50000000!(Cp)!3pt!90:(Cp)$) -- ($(B)!0.50000000!(Cp)!3pt!-90:(Cp)$);
\draw[inkiron] ($(B)!0.48260747!(D)!3pt!90:(D)$) -- ($(B)!0.48260747!(D)!3pt!-90:(D)$);
\draw[inkiron] ($(B)!0.51739253!(D)!3pt!90:(D)$) -- ($(B)!0.51739253!(D)!3pt!-90:(D)$);
\draw[inkiron] ($(B)!0.48260747!(Dp)!3pt!90:(Dp)$) -- ($(B)!0.48260747!(Dp)!3pt!-90:(Dp)$);
\draw[inkiron] ($(B)!0.51739253!(Dp)!3pt!90:(Dp)$) -- ($(B)!0.51739253!(Dp)!3pt!-90:(Dp)$);
\coordinate (E) at (-3.50000000,2.50000000);
\coordinate (G) at (-7.00000000,5.00000000);
\coordinate (H) at (1.51351351,-1.08108108);
\draw[indigo,dashed] (G) -- (H);
\draw[indigo,dashed] (C) -- (G);
\draw[indigo] ($(C)!0.45685445!(E)!3pt!90:(E)$) -- ($(C)!0.45685445!(E)!3pt!-90:(E)$);
\draw[indigo] ($(C)!0.50000000!(E)!3pt!90:(E)$) -- ($(C)!0.50000000!(E)!3pt!-90:(E)$);
\draw[indigo] ($(C)!0.54314555!(E)!3pt!90:(E)$) -- ($(C)!0.54314555!(E)!3pt!-90:(E)$);
\draw[indigo] ($(E)!0.45685445!(Dp)!3pt!90:(Dp)$) -- ($(E)!0.45685445!(Dp)!3pt!-90:(Dp)$);
\draw[indigo] ($(E)!0.50000000!(Dp)!3pt!90:(Dp)$) -- ($(E)!0.50000000!(Dp)!3pt!-90:(Dp)$);
\draw[indigo] ($(E)!0.54314555!(Dp)!3pt!90:(Dp)$) -- ($(E)!0.54314555!(Dp)!3pt!-90:(Dp)$);
\draw[madder] ($(B)!0.48721276!(E)!3pt!90:(E)$) -- ($(B)!0.48721276!(E)!3pt!-90:(E)$);
\draw[madder] ($(B)!0.51278724!(E)!3pt!90:(E)$) -- ($(B)!0.51278724!(E)!3pt!-90:(E)$);
\draw[madder] ($(E)!0.48721276!(G)!3pt!90:(G)$) -- ($(E)!0.48721276!(G)!3pt!-90:(G)$);
\draw[madder] ($(E)!0.51278724!(G)!3pt!90:(G)$) -- ($(E)!0.51278724!(G)!3pt!-90:(G)$);
\fill[inkiron] (B) circle (1.05pt);
\node at (-0.46000000,-0.45000000) {$B$};
\fill[inkiron] (C) circle (1.05pt);
\node at (-6.45000000,2.10000000) {$C$};
\fill[inkiron] (Cp) circle (1.05pt);
\node at (-2.20000000,-6.42000000) {$C'$};
\fill[inkiron] (D) circle (1.05pt);
\node at (3.36000000,1.12000000) {$D$};
\fill[inkiron] (Dp) circle (1.05pt);
\node at (-0.68000000,3.26000000) {$D'$};
\fill[inkiron] (E) circle (1.05pt);
\node at (-3.45000000,2.93000000) {$E$};
\fill[inkiron] (G) circle (1.05pt);
\node at (-7.20000000,5.35000000) {$G$};
\fill[inkiron] (H) circle (1.05pt);
\node at (1.81351351,-1.31108108) {$H$};
\end{tikzpicture}}

\newcommand{\diConverse}{%
\begin{tikzpicture}[scale=0.52,line width=.7pt,every node/.style={font=\small,inner sep=2pt}]
\coordinate (B) at (0.00000000,0.00000000);
\coordinate (C) at (-6.00000000,2.00000000);
\coordinate (Cp) at (-2.00000000,-6.00000000);
\coordinate (D) at (3.00000000,1.00000000);
\coordinate (Dp) at (-1.00000000,3.00000000);
\fill[selvage!10] (C) -- (B) -- (Dp) -- cycle;
\fill[madder!10] (Cp) -- (B) -- (D) -- cycle;
\draw[inkiron] (B) -- (C);
\draw[inkiron] (B) -- (Cp);
\draw[inkiron] (B) -- (D);
\draw[inkiron] (B) -- (Dp);
\draw[inkiron] (C) -- (Dp);
\draw[inkiron] (Cp) -- (D);
\pic[draw=selvage,angle radius=2.2mm] {right angle=C--B--Cp};
\pic[draw=selvage,angle radius=2.2mm] {right angle=D--B--Dp};
\draw[inkiron] ($(B)!0.50000000!(C)!3pt!90:(C)$) -- ($(B)!0.50000000!(C)!3pt!-90:(C)$);
\draw[inkiron] ($(B)!0.50000000!(Cp)!3pt!90:(Cp)$) -- ($(B)!0.50000000!(Cp)!3pt!-90:(Cp)$);
\draw[inkiron] ($(B)!0.48260747!(D)!3pt!90:(D)$) -- ($(B)!0.48260747!(D)!3pt!-90:(D)$);
\draw[inkiron] ($(B)!0.51739253!(D)!3pt!90:(D)$) -- ($(B)!0.51739253!(D)!3pt!-90:(D)$);
\draw[inkiron] ($(B)!0.48260747!(Dp)!3pt!90:(Dp)$) -- ($(B)!0.48260747!(Dp)!3pt!-90:(Dp)$);
\draw[inkiron] ($(B)!0.51739253!(Dp)!3pt!90:(Dp)$) -- ($(B)!0.51739253!(Dp)!3pt!-90:(Dp)$);
\coordinate (E) at (-3.50000000,2.50000000);
\coordinate (F) at (1.51351351,-1.08108108);
\coordinate (M) at (-4.91891892,3.51351351);
\coordinate (N) at (-2.08108108,1.48648649);
\draw[indigo,dashed] (M) -- (F);
\draw[indigo,dashed] (C) -- (M);
\draw[indigo,dashed] (Dp) -- (N);
\pic[draw=selvage,angle radius=2.2mm] {right angle=C--M--B};
\pic[draw=selvage,angle radius=2.2mm] {right angle=Dp--N--B};
\pic[draw=selvage,angle radius=2.2mm] {right angle=B--F--D};
\fill[inkiron] (B) circle (1.05pt);
\node at (-0.46000000,-0.45000000) {$B$};
\fill[inkiron] (C) circle (1.05pt);
\node at (-6.45000000,2.10000000) {$C$};
\fill[inkiron] (Cp) circle (1.05pt);
\node at (-2.20000000,-6.42000000) {$C'$};
\fill[inkiron] (D) circle (1.05pt);
\node at (3.36000000,1.12000000) {$D$};
\fill[inkiron] (Dp) circle (1.05pt);
\node at (-0.68000000,3.26000000) {$D'$};
\fill[inkiron] (E) circle (1.05pt);
\node at (-3.46000000,2.90000000) {$E$};
\fill[inkiron] (F) circle (1.05pt);
\node at (1.81351351,-1.31108108) {$F$};
\fill[inkiron] (M) circle (1.05pt);
\node at (-5.01891892,3.91351351) {$M$};
\fill[inkiron] (N) circle (1.05pt);
\node at (-2.20108108,1.12648649) {$N$};
\end{tikzpicture}}

% Generated from the same coordinates as PNG/SVG.
\newcommand{\esdiModel}{%
\begin{tikzpicture}[scale=.8,line width=.7pt,every node/.style={font=\small,inner sep=2pt}]
\coordinate (A) at (4.00000000,4.00000000);
\coordinate (B) at (0.00000000,0.00000000);
\coordinate (C) at (8.00000000,0.00000000);
\coordinate (D) at (7.20000000,2.40000000);
\coordinate (O) at (6.00000000,2.00000000);
\draw[inkiron] (A) -- (B);
\draw[inkiron] (B) -- (C);
\draw[inkiron] (C) -- (D);
\draw[inkiron] (D) -- (A);
\draw[inkiron] (A) -- (C);
\draw[inkiron] (B) -- (D);
\filldraw[fill=esblue!18,draw=esblue,line width=1pt] (A) -- (B) -- (O) -- cycle;
\filldraw[fill=esorange!18,draw=esorange,line width=1pt] (D) -- (C) -- (O) -- cycle;
\draw[selvage] (3.87979185,3.87979185) -- (4.00000000,3.75958369) -- (4.12020815,3.87979185);
\draw[selvage] (7.03872384,2.34624128) -- (7.09248256,2.18496512) -- (7.25375872,2.23872384);
\fill[inkiron] (A) circle (1pt);
\node at (3.85000000,4.28000000) {$A$};
\fill[inkiron] (B) circle (1pt);
\node at (-0.25000000,-0.25000000) {$B$};
\fill[inkiron] (C) circle (1pt);
\node at (8.25000000,-0.25000000) {$C$};
\fill[inkiron] (D) circle (1pt);
\node at (7.55000000,2.43000000) {$D$};
\fill[inkiron] (O) circle (1pt);
\node at (5.70000000,1.75000000) {$O$};
\end{tikzpicture}}

\newcommand{\esdiExterior}{%
\begin{tikzpicture}[scale=.8,line width=.7pt,every node/.style={font=\small,inner sep=2pt}]
\coordinate (A) at (4.00000000,4.00000000);
\coordinate (B) at (0.00000000,0.00000000);
\coordinate (C) at (8.00000000,0.00000000);
\coordinate (D) at (7.20000000,2.40000000);
\coordinate (X) at (6.08000000,5.76000000);
\draw[inkiron] (A) -- (B);
\draw[inkiron] (B) -- (C);
\draw[inkiron] (C) -- (D);
\draw[inkiron] (D) -- (A);
\draw[inkiron] (A) -- (C);
\draw[inkiron] (B) -- (D);
\draw[esblue,dashed] (D) -- (X);
\draw[esblue] (6.67822419,2.22607473) -- (6.67490997,2.23635264) -- (6.67179818,2.24669364) -- (6.66889003,2.25709374) -- (6.66618663,2.26754894) -- (6.66368902,2.27805520) -- (6.66139817,2.28860847) -- (6.65931497,2.29920468) -- (6.65744021,2.30983976) -- (6.65577461,2.32050959) -- (6.65431882,2.33121007) -- (6.65307341,2.34193706) -- (6.65203885,2.35268644) -- (6.65121553,2.36345406) -- (6.65060378,2.37423577) -- (6.65020384,2.38502742) -- (6.65001585,2.39582483) -- (6.65003989,2.40662386) -- (6.65027595,2.41742033) -- (6.65072394,2.42821008) -- (6.65138369,2.43898896) -- (6.65225493,2.44975281) -- (6.65333735,2.46049748) -- (6.65463051,2.47121882) -- (6.65613392,2.48191271) -- (6.65784701,2.49257502) -- (6.65976910,2.50320164) -- (6.66189947,2.51378848) -- (6.66423728,2.52433144) -- (6.66678164,2.53482648) -- (6.66953156,2.54526953) -- (6.67248599,2.55565658) -- (6.67564379,2.56598363) -- (6.67900374,2.57624668) -- (6.68256454,2.58644179) -- (6.68632482,2.59656502) -- (6.69028313,2.60661247) -- (6.69443795,2.61658027) -- (6.69878767,2.62646457) -- (6.70333062,2.63626157) -- (6.70806504,2.64596748);
\draw[esblue] (6.70806504,2.64596748) -- (6.71298912,2.65557856) -- (6.71810095,2.66509112) -- (6.72339855,2.67450148) -- (6.72887990,2.68380601) -- (6.73454287,2.69300113) -- (6.74038529,2.70208330) -- (6.74640490,2.71104900) -- (6.75259937,2.71989479) -- (6.75896633,2.72861725) -- (6.76550331,2.73721303) -- (6.77220780,2.74567880) -- (6.77907721,2.75401131) -- (6.78610890,2.76220734) -- (6.79330015,2.77026373) -- (6.80064819,2.77817738) -- (6.80815019,2.78594523) -- (6.81580325,2.79356430) -- (6.82360443,2.80103164) -- (6.83155071,2.80834437) -- (6.83963904,2.81549968) -- (6.84786630,2.82249480) -- (6.85622931,2.82932704) -- (6.86472485,2.83599378) -- (6.87334964,2.84249242) -- (6.88210037,2.84882048) -- (6.89097365,2.85497551) -- (6.89996607,2.86095514) -- (6.90907415,2.86675706) -- (6.91829440,2.87237904) -- (6.92762324,2.87781890) -- (6.93705710,2.88307456) -- (6.94659232,2.88814398) -- (6.95622523,2.89302522) -- (6.96595213,2.89771638) -- (6.97576925,2.90221566) -- (6.98567282,2.90652133) -- (6.99565902,2.91063173) -- (7.00572400,2.91454527) -- (7.01586387,2.91826044) -- (7.02607473,2.92177581);
\draw[selvage] (3.87979185,3.87979185) -- (4.00000000,3.75958369) -- (4.12020815,3.87979185);
\draw[selvage] (7.03872384,2.34624128) -- (7.09248256,2.18496512) -- (7.25375872,2.23872384);
\draw[inkiron] ($(A)!0.50000000!(B)!3pt!90:(B)$) -- ($(A)!0.50000000!(B)!3pt!-90:(B)$);
\draw[inkiron] ($(A)!0.30000000!(C)!3pt!90:(C)$) -- ($(A)!0.30000000!(C)!3pt!-90:(C)$);
\fill[inkiron] (A) circle (1pt);
\node at (3.85000000,4.28000000) {$A$};
\fill[inkiron] (B) circle (1pt);
\node at (-0.25000000,-0.25000000) {$B$};
\fill[inkiron] (C) circle (1pt);
\node at (8.25000000,-0.25000000) {$C$};
\fill[inkiron] (D) circle (1pt);
\node at (7.55000000,2.43000000) {$D$};
\fill[inkiron] (X) circle (1pt);
\node at (6.08000000,6.01000000) {$X$};
\end{tikzpicture}}

\newcommand{\esdiSharedHypotenuse}{%
\begin{tikzpicture}[scale=.8,line width=.7pt,every node/.style={font=\small,inner sep=2pt}]
\coordinate (A) at (4.00000000,4.00000000);
\coordinate (B) at (0.00000000,0.00000000);
\coordinate (C) at (8.00000000,0.00000000);
\coordinate (D) at (7.20000000,2.40000000);
\draw[inkiron] (A) -- (B);
\draw[inkiron] (B) -- (C);
\draw[inkiron] (C) -- (D);
\draw[inkiron] (D) -- (A);
\draw[inkiron] (A) -- (C);
\draw[inkiron] (B) -- (D);
\draw[selvage] (3.87979185,3.87979185) -- (4.00000000,3.75958369) -- (4.12020815,3.87979185);
\draw[selvage] (7.03872384,2.34624128) -- (7.09248256,2.18496512) -- (7.25375872,2.23872384);
\draw[inkiron] ($(A)!0.50000000!(B)!3pt!90:(B)$) -- ($(A)!0.50000000!(B)!3pt!-90:(B)$);
\draw[inkiron] ($(A)!0.30000000!(C)!3pt!90:(C)$) -- ($(A)!0.30000000!(C)!3pt!-90:(C)$);
\fill[inkiron] (A) circle (1pt);
\node at (3.85000000,4.28000000) {$A$};
\fill[inkiron] (B) circle (1pt);
\node at (-0.25000000,-0.25000000) {$B$};
\fill[inkiron] (C) circle (1pt);
\node at (8.25000000,-0.25000000) {$C$};
\fill[inkiron] (D) circle (1pt);
\node at (7.55000000,2.43000000) {$D$};
\end{tikzpicture}}

\newcommand{\esdiDoublePerpendicular}{%
\begin{tikzpicture}[scale=.8,line width=.7pt,every node/.style={font=\small,inner sep=2pt}]
\coordinate (A) at (4.00000000,4.00000000);
\coordinate (B) at (0.00000000,0.00000000);
\coordinate (C) at (8.00000000,0.00000000);
\coordinate (D) at (7.20000000,2.40000000);
\coordinate (M) at (4.80000000,1.60000000);
\coordinate (N) at (6.40000000,4.80000000);
\draw[inkiron] (A) -- (B);
\draw[inkiron] (B) -- (C);
\draw[inkiron] (C) -- (D);
\draw[inkiron] (D) -- (A);
\draw[inkiron] (A) -- (C);
\draw[inkiron] (B) -- (D);
\draw[esblue,dashed] (A) -- (M);
\draw[esblue,dashed] (A) -- (N);
\draw[esblue,dashed] (D) -- (N);
\draw[selvage] (4.74624128,1.76127616) -- (4.58496512,1.70751744) -- (4.63872384,1.54624128);
\draw[selvage] (6.23872384,4.74624128) -- (6.29248256,4.58496512) -- (6.45375872,4.63872384);
\fill[inkiron] (A) circle (1pt);
\node at (3.85000000,4.28000000) {$A$};
\fill[inkiron] (B) circle (1pt);
\node at (-0.25000000,-0.25000000) {$B$};
\fill[inkiron] (C) circle (1pt);
\node at (8.25000000,-0.25000000) {$C$};
\fill[inkiron] (D) circle (1pt);
\node at (7.55000000,2.43000000) {$D$};
\fill[inkiron] (M) circle (1pt);
\node at (4.80000000,1.25000000) {$M$};
\fill[inkiron] (N) circle (1pt);
\node at (6.65000000,5.05000000) {$N$};
\end{tikzpicture}}

\newcommand{\esdiExtendCd}{%
\begin{tikzpicture}[scale=.8,line width=.7pt,every node/.style={font=\small,inner sep=2pt}]
\coordinate (A) at (4.00000000,4.00000000);
\coordinate (B) at (0.00000000,0.00000000);
\coordinate (C) at (8.00000000,0.00000000);
\coordinate (D) at (7.20000000,2.40000000);
\coordinate (E) at (4.80000000,9.60000000);
\draw[inkiron] (A) -- (B);
\draw[inkiron] (B) -- (C);
\draw[inkiron] (C) -- (D);
\draw[inkiron] (D) -- (A);
\draw[inkiron] (A) -- (C);
\draw[inkiron] (B) -- (D);
\draw[esblue,dashed] (D) -- (E);
\draw[esblue,dashed] (A) -- (E);
\draw[inkiron] ($(D)!0.49275311!(E)!3pt!90:(E)$) -- ($(D)!0.49275311!(E)!3pt!-90:(E)$);
\draw[inkiron] ($(D)!0.50724689!(E)!3pt!90:(E)$) -- ($(D)!0.50724689!(E)!3pt!-90:(E)$);
\draw[inkiron] ($(D)!0.49275311!(B)!3pt!90:(B)$) -- ($(D)!0.49275311!(B)!3pt!-90:(B)$);
\draw[inkiron] ($(D)!0.50724689!(B)!3pt!90:(B)$) -- ($(D)!0.50724689!(B)!3pt!-90:(B)$);
\fill[inkiron] (A) circle (1pt);
\node at (3.85000000,4.28000000) {$A$};
\fill[inkiron] (B) circle (1pt);
\node at (-0.25000000,-0.25000000) {$B$};
\fill[inkiron] (C) circle (1pt);
\node at (8.25000000,-0.25000000) {$C$};
\fill[inkiron] (D) circle (1pt);
\node at (7.55000000,2.43000000) {$D$};
\fill[inkiron] (E) circle (1pt);
\node at (4.80000000,9.85000000) {$E$};
\end{tikzpicture}}

\newcommand{\esdiExtendBd}{%
\begin{tikzpicture}[scale=.8,line width=.7pt,every node/.style={font=\small,inner sep=2pt}]
\coordinate (A) at (4.00000000,4.00000000);
\coordinate (B) at (0.00000000,0.00000000);
\coordinate (C) at (8.00000000,0.00000000);
\coordinate (D) at (7.20000000,2.40000000);
\coordinate (F) at (9.60000000,3.20000000);
\draw[inkiron] (A) -- (B);
\draw[inkiron] (B) -- (C);
\draw[inkiron] (C) -- (D);
\draw[inkiron] (D) -- (A);
\draw[inkiron] (A) -- (C);
\draw[inkiron] (B) -- (D);
\draw[esblue,dashed] (D) -- (F);
\draw[esblue,dashed] (A) -- (F);
\draw[inkiron] ($(D)!0.47825934!(F)!3pt!90:(F)$) -- ($(D)!0.47825934!(F)!3pt!-90:(F)$);
\draw[inkiron] ($(D)!0.52174066!(F)!3pt!90:(F)$) -- ($(D)!0.52174066!(F)!3pt!-90:(F)$);
\draw[inkiron] ($(D)!0.47825934!(C)!3pt!90:(C)$) -- ($(D)!0.47825934!(C)!3pt!-90:(C)$);
\draw[inkiron] ($(D)!0.52174066!(C)!3pt!90:(C)$) -- ($(D)!0.52174066!(C)!3pt!-90:(C)$);
\fill[inkiron] (A) circle (1pt);
\node at (3.85000000,4.28000000) {$A$};
\fill[inkiron] (B) circle (1pt);
\node at (-0.25000000,-0.25000000) {$B$};
\fill[inkiron] (C) circle (1pt);
\node at (8.25000000,-0.25000000) {$C$};
\fill[inkiron] (D) circle (1pt);
\node at (7.55000000,2.43000000) {$D$};
\fill[inkiron] (F) circle (1pt);
\node at (9.60000000,3.45000000) {$F$};
\end{tikzpicture}}

\newcommand{\esdiHandLower}{%
\begin{tikzpicture}[scale=.8,line width=.7pt,every node/.style={font=\small,inner sep=2pt}]
\coordinate (A) at (4.00000000,4.00000000);
\coordinate (B) at (0.00000000,0.00000000);
\coordinate (C) at (8.00000000,0.00000000);
\coordinate (D) at (7.20000000,2.40000000);
\coordinate (P) at (2.40000000,0.80000000);
\draw[inkiron] (A) -- (B);
\draw[inkiron] (B) -- (C);
\draw[inkiron] (C) -- (D);
\draw[inkiron] (D) -- (A);
\draw[inkiron] (A) -- (C);
\draw[inkiron] (B) -- (D);
\draw[esblue,dashed] (A) -- (P);
\draw[selvage] (3.92397369,3.84794738) -- (4.07602631,3.77192107) -- (4.15205262,3.92397369);
\draw[inkiron] ($(A)!0.48462703!(P)!3pt!90:(P)$) -- ($(A)!0.48462703!(P)!3pt!-90:(P)$);
\draw[inkiron] ($(A)!0.51537297!(P)!3pt!90:(P)$) -- ($(A)!0.51537297!(P)!3pt!-90:(P)$);
\draw[inkiron] ($(A)!0.48462703!(D)!3pt!90:(D)$) -- ($(A)!0.48462703!(D)!3pt!-90:(D)$);
\draw[inkiron] ($(A)!0.51537297!(D)!3pt!90:(D)$) -- ($(A)!0.51537297!(D)!3pt!-90:(D)$);
\fill[inkiron] (A) circle (1pt);
\node at (3.85000000,4.28000000) {$A$};
\fill[inkiron] (B) circle (1pt);
\node at (-0.25000000,-0.25000000) {$B$};
\fill[inkiron] (C) circle (1pt);
\node at (8.25000000,-0.25000000) {$C$};
\fill[inkiron] (D) circle (1pt);
\node at (7.55000000,2.43000000) {$D$};
\fill[inkiron] (P) circle (1pt);
\node at (2.40000000,0.45000000) {$P$};
\end{tikzpicture}}

\newcommand{\esdiHandUpper}{%
\begin{tikzpicture}[scale=.8,line width=.7pt,every node/.style={font=\small,inner sep=2pt}]
\coordinate (A) at (4.00000000,4.00000000);
\coordinate (B) at (0.00000000,0.00000000);
\coordinate (C) at (8.00000000,0.00000000);
\coordinate (D) at (7.20000000,2.40000000);
\coordinate (Q) at (5.60000000,7.20000000);
\draw[inkiron] (A) -- (B);
\draw[inkiron] (B) -- (C);
\draw[inkiron] (C) -- (D);
\draw[inkiron] (D) -- (A);
\draw[inkiron] (A) -- (C);
\draw[inkiron] (B) -- (D);
\draw[esblue,dashed] (D) -- (Q);
\draw[esblue,dashed] (A) -- (Q);
\draw[selvage] (4.07602631,4.15205262) -- (4.22807893,4.07602631) -- (4.15205262,3.92397369);
\draw[inkiron] ($(A)!0.48462703!(Q)!3pt!90:(Q)$) -- ($(A)!0.48462703!(Q)!3pt!-90:(Q)$);
\draw[inkiron] ($(A)!0.51537297!(Q)!3pt!90:(Q)$) -- ($(A)!0.51537297!(Q)!3pt!-90:(Q)$);
\draw[inkiron] ($(A)!0.48462703!(D)!3pt!90:(D)$) -- ($(A)!0.48462703!(D)!3pt!-90:(D)$);
\draw[inkiron] ($(A)!0.51537297!(D)!3pt!90:(D)$) -- ($(A)!0.51537297!(D)!3pt!-90:(D)$);
\fill[inkiron] (A) circle (1pt);
\node at (3.85000000,4.28000000) {$A$};
\fill[inkiron] (B) circle (1pt);
\node at (-0.25000000,-0.25000000) {$B$};
\fill[inkiron] (C) circle (1pt);
\node at (8.25000000,-0.25000000) {$C$};
\fill[inkiron] (D) circle (1pt);
\node at (7.55000000,2.43000000) {$D$};
\fill[inkiron] (Q) circle (1pt);
\node at (5.60000000,7.45000000) {$Q$};
\end{tikzpicture}}

\newcommand{\esdiKPerpendicular}{%
\begin{tikzpicture}[scale=.8,line width=.7pt,every node/.style={font=\small,inner sep=2pt}]
\coordinate (A) at (4.00000000,4.00000000);
\coordinate (B) at (0.00000000,0.00000000);
\coordinate (C) at (8.00000000,0.00000000);
\coordinate (D) at (7.20000000,2.40000000);
\coordinate (M) at (2.40000000,4.80000000);
\coordinate (N) at (8.80000000,1.60000000);
\draw[inkiron] (A) -- (B);
\draw[inkiron] (B) -- (C);
\draw[inkiron] (C) -- (D);
\draw[inkiron] (D) -- (A);
\draw[inkiron] (A) -- (C);
\draw[inkiron] (B) -- (D);
\draw[esblue,dashed] (B) -- (M);
\draw[esblue,dashed] (M) -- (A);
\draw[esblue,dashed] (C) -- (N);
\draw[esblue,dashed] (D) -- (N);
\draw[selvage] (2.32397369,4.64794738) -- (2.47602631,4.57192107) -- (2.55205262,4.72397369);
\draw[selvage] (8.72397369,1.44794738) -- (8.57192107,1.52397369) -- (8.64794738,1.67602631);
\fill[inkiron] (A) circle (1pt);
\node at (3.85000000,4.28000000) {$A$};
\fill[inkiron] (B) circle (1pt);
\node at (-0.25000000,-0.25000000) {$B$};
\fill[inkiron] (C) circle (1pt);
\node at (8.25000000,-0.25000000) {$C$};
\fill[inkiron] (D) circle (1pt);
\node at (7.55000000,2.43000000) {$D$};
\fill[inkiron] (M) circle (1pt);
\node at (2.20000000,5.05000000) {$M$};
\fill[inkiron] (N) circle (1pt);
\node at (9.10000000,1.65000000) {$N$};
\end{tikzpicture}}

\newcommand{\esdiExample}{%
\begin{tikzpicture}[scale=.8,line width=.7pt,every node/.style={font=\small,inner sep=2pt}]
\coordinate (A) at (4.00000000,4.00000000);
\coordinate (B) at (0.00000000,0.00000000);
\coordinate (C) at (8.00000000,0.00000000);
\coordinate (D) at (5.65685425,2.34314575);
\coordinate (E) at (6.82842712,2.82842712);
\draw[inkiron] (A) -- (B);
\draw[inkiron] (B) -- (C);
\draw[inkiron] (A) -- (C);
\draw[inkiron] (B) -- (E);
\draw[inkiron] (C) -- (E);
\draw[selvage] (3.87979185,3.87979185) -- (4.00000000,3.75958369) -- (4.12020815,3.87979185);
\draw[selvage] (6.67136760,2.76337094) -- (6.73642379,2.60631142) -- (6.89348331,2.67136760);
\draw[inkiron] ($(A)!0.50000000!(B)!3pt!90:(B)$) -- ($(A)!0.50000000!(B)!3pt!-90:(B)$);
\draw[inkiron] ($(A)!0.50000000!(C)!3pt!90:(C)$) -- ($(A)!0.50000000!(C)!3pt!-90:(C)$);
\draw[esblue] (0.77781746,0.77781746) -- (0.78541606,0.77014389) -- (0.79293896,0.76239610) -- (0.80038543,0.75457482) -- (0.80775476,0.74668082) -- (0.81504624,0.73871485) -- (0.82225916,0.73067768) -- (0.82939283,0.72257009) -- (0.83644656,0.71439285) -- (0.84341968,0.70614676) -- (0.85031150,0.69783261) -- (0.85712137,0.68945120) -- (0.86384862,0.68100334) -- (0.87049262,0.67248985) -- (0.87705272,0.66391154) -- (0.88352828,0.65526923) -- (0.88991869,0.64656378) -- (0.89622333,0.63779600) -- (0.90244159,0.62896676) -- (0.90857287,0.62007689) -- (0.91461657,0.61112726) -- (0.92057213,0.60211872) -- (0.92643896,0.59305216) -- (0.93221649,0.58392843) -- (0.93790418,0.57474842) -- (0.94350147,0.56551302) -- (0.94900782,0.55622311) -- (0.95442271,0.54687959) -- (0.95974561,0.53748337) -- (0.96497600,0.52803533) -- (0.97011339,0.51853641) -- (0.97515728,0.50898751) -- (0.98010718,0.49938955) -- (0.98496261,0.48974346) -- (0.98972311,0.48005016) -- (0.99438822,0.47031060) -- (0.99895749,0.46052571) -- (1.00343048,0.45069643) -- (1.00780675,0.44082372) -- (1.01208589,0.43090851) -- (1.01626749,0.42095178);
\draw[esblue] (1.01626749,0.42095178) -- (1.02035113,0.41095447) -- (1.02433643,0.40091755) -- (1.02822300,0.39084199) -- (1.03201047,0.38072876) -- (1.03569847,0.37057884) -- (1.03928665,0.36039320) -- (1.04277466,0.35017282) -- (1.04616217,0.33991869) -- (1.04944884,0.32963180) -- (1.05263437,0.31931314) -- (1.05571844,0.30896371) -- (1.05870076,0.29858449) -- (1.06158104,0.28817650) -- (1.06435900,0.27774073) -- (1.06703438,0.26727820) -- (1.06960691,0.25678990) -- (1.07207636,0.24627685) -- (1.07444247,0.23574007) -- (1.07670503,0.22518056) -- (1.07886381,0.21459935) -- (1.08091861,0.20399746) -- (1.08286922,0.19337591) -- (1.08471547,0.18273572) -- (1.08645717,0.17207791) -- (1.08809416,0.16140352) -- (1.08962627,0.15071358) -- (1.09105337,0.14000910) -- (1.09237530,0.12929114) -- (1.09359195,0.11856071) -- (1.09470320,0.10781885) -- (1.09570894,0.09706661) -- (1.09660907,0.08630501) -- (1.09740350,0.07553508) -- (1.09809217,0.06475788) -- (1.09867500,0.05397444) -- (1.09915194,0.04318580) -- (1.09952294,0.03239299) -- (1.09978796,0.02159706) -- (1.09994699,0.01079905) -- (1.10000000,-0.00000000);
\fill[inkiron] (A) circle (1pt);
\node at (3.90000000,4.28000000) {$A$};
\fill[inkiron] (B) circle (1pt);
\node at (-0.25000000,-0.25000000) {$B$};
\fill[inkiron] (C) circle (1pt);
\node at (8.25000000,-0.25000000) {$C$};
\fill[inkiron] (D) circle (1pt);
\node at (6.00685425,2.69314575) {$D$};
\fill[inkiron] (E) circle (1pt);
\node at (7.12842712,3.00842712) {$E$};
\end{tikzpicture}}

\newcommand{\esdiExampleExtend}{%
\begin{tikzpicture}[scale=.8,line width=.7pt,every node/.style={font=\small,inner sep=2pt}]
\coordinate (A) at (4.00000000,4.00000000);
\coordinate (B) at (0.00000000,0.00000000);
\coordinate (C) at (8.00000000,0.00000000);
\coordinate (D) at (5.65685425,2.34314575);
\coordinate (E) at (6.82842712,2.82842712);
\coordinate (F) at (5.65685425,5.65685425);
\draw[inkiron] (A) -- (B);
\draw[inkiron] (B) -- (C);
\draw[inkiron] (A) -- (C);
\draw[inkiron] (B) -- (E);
\draw[inkiron] (C) -- (E);
\draw[selvage] (3.87979185,3.87979185) -- (4.00000000,3.75958369) -- (4.12020815,3.87979185);
\draw[selvage] (6.67136760,2.76337094) -- (6.73642379,2.60631142) -- (6.89348331,2.67136760);
\draw[inkiron] ($(A)!0.50000000!(B)!3pt!90:(B)$) -- ($(A)!0.50000000!(B)!3pt!-90:(B)$);
\draw[inkiron] ($(A)!0.50000000!(C)!3pt!90:(C)$) -- ($(A)!0.50000000!(C)!3pt!-90:(C)$);
\draw[esblue] (0.77781746,0.77781746) -- (0.78541606,0.77014389) -- (0.79293896,0.76239610) -- (0.80038543,0.75457482) -- (0.80775476,0.74668082) -- (0.81504624,0.73871485) -- (0.82225916,0.73067768) -- (0.82939283,0.72257009) -- (0.83644656,0.71439285) -- (0.84341968,0.70614676) -- (0.85031150,0.69783261) -- (0.85712137,0.68945120) -- (0.86384862,0.68100334) -- (0.87049262,0.67248985) -- (0.87705272,0.66391154) -- (0.88352828,0.65526923) -- (0.88991869,0.64656378) -- (0.89622333,0.63779600) -- (0.90244159,0.62896676) -- (0.90857287,0.62007689) -- (0.91461657,0.61112726) -- (0.92057213,0.60211872) -- (0.92643896,0.59305216) -- (0.93221649,0.58392843) -- (0.93790418,0.57474842) -- (0.94350147,0.56551302) -- (0.94900782,0.55622311) -- (0.95442271,0.54687959) -- (0.95974561,0.53748337) -- (0.96497600,0.52803533) -- (0.97011339,0.51853641) -- (0.97515728,0.50898751) -- (0.98010718,0.49938955) -- (0.98496261,0.48974346) -- (0.98972311,0.48005016) -- (0.99438822,0.47031060) -- (0.99895749,0.46052571) -- (1.00343048,0.45069643) -- (1.00780675,0.44082372) -- (1.01208589,0.43090851) -- (1.01626749,0.42095178);
\draw[esblue] (1.01626749,0.42095178) -- (1.02035113,0.41095447) -- (1.02433643,0.40091755) -- (1.02822300,0.39084199) -- (1.03201047,0.38072876) -- (1.03569847,0.37057884) -- (1.03928665,0.36039320) -- (1.04277466,0.35017282) -- (1.04616217,0.33991869) -- (1.04944884,0.32963180) -- (1.05263437,0.31931314) -- (1.05571844,0.30896371) -- (1.05870076,0.29858449) -- (1.06158104,0.28817650) -- (1.06435900,0.27774073) -- (1.06703438,0.26727820) -- (1.06960691,0.25678990) -- (1.07207636,0.24627685) -- (1.07444247,0.23574007) -- (1.07670503,0.22518056) -- (1.07886381,0.21459935) -- (1.08091861,0.20399746) -- (1.08286922,0.19337591) -- (1.08471547,0.18273572) -- (1.08645717,0.17207791) -- (1.08809416,0.16140352) -- (1.08962627,0.15071358) -- (1.09105337,0.14000910) -- (1.09237530,0.12929114) -- (1.09359195,0.11856071) -- (1.09470320,0.10781885) -- (1.09570894,0.09706661) -- (1.09660907,0.08630501) -- (1.09740350,0.07553508) -- (1.09809217,0.06475788) -- (1.09867500,0.05397444) -- (1.09915194,0.04318580) -- (1.09952294,0.03239299) -- (1.09978796,0.02159706) -- (1.09994699,0.01079905) -- (1.10000000,-0.00000000);
\draw[esblue,dashed] (A) -- (F);
\draw[esblue,dashed] (E) -- (F);
\fill[inkiron] (A) circle (1pt);
\node at (3.90000000,4.28000000) {$A$};
\fill[inkiron] (B) circle (1pt);
\node at (-0.25000000,-0.25000000) {$B$};
\fill[inkiron] (C) circle (1pt);
\node at (8.25000000,-0.25000000) {$C$};
\fill[inkiron] (D) circle (1pt);
\node at (6.00685425,2.69314575) {$D$};
\fill[inkiron] (E) circle (1pt);
\node at (7.12842712,3.00842712) {$E$};
\fill[inkiron] (F) circle (1pt);
\node at (5.65685425,5.90685425) {$F$};
\end{tikzpicture}}

\newcommand{\esdiExampleParallel}{%
\begin{tikzpicture}[scale=.8,line width=.7pt,every node/.style={font=\small,inner sep=2pt}]
\coordinate (A) at (4.00000000,4.00000000);
\coordinate (B) at (0.00000000,0.00000000);
\coordinate (C) at (8.00000000,0.00000000);
\coordinate (D) at (5.65685425,2.34314575);
\coordinate (E) at (6.82842712,2.82842712);
\coordinate (F) at (3.31370850,0.00000000);
\coordinate (G) at (2.82842712,1.17157288);
\draw[inkiron] (A) -- (B);
\draw[inkiron] (B) -- (C);
\draw[inkiron] (A) -- (C);
\draw[inkiron] (B) -- (E);
\draw[inkiron] (C) -- (E);
\draw[selvage] (3.87979185,3.87979185) -- (4.00000000,3.75958369) -- (4.12020815,3.87979185);
\draw[selvage] (6.67136760,2.76337094) -- (6.73642379,2.60631142) -- (6.89348331,2.67136760);
\draw[inkiron] ($(A)!0.50000000!(B)!3pt!90:(B)$) -- ($(A)!0.50000000!(B)!3pt!-90:(B)$);
\draw[inkiron] ($(A)!0.50000000!(C)!3pt!90:(C)$) -- ($(A)!0.50000000!(C)!3pt!-90:(C)$);
\draw[esblue] (0.77781746,0.77781746) -- (0.78541606,0.77014389) -- (0.79293896,0.76239610) -- (0.80038543,0.75457482) -- (0.80775476,0.74668082) -- (0.81504624,0.73871485) -- (0.82225916,0.73067768) -- (0.82939283,0.72257009) -- (0.83644656,0.71439285) -- (0.84341968,0.70614676) -- (0.85031150,0.69783261) -- (0.85712137,0.68945120) -- (0.86384862,0.68100334) -- (0.87049262,0.67248985) -- (0.87705272,0.66391154) -- (0.88352828,0.65526923) -- (0.88991869,0.64656378) -- (0.89622333,0.63779600) -- (0.90244159,0.62896676) -- (0.90857287,0.62007689) -- (0.91461657,0.61112726) -- (0.92057213,0.60211872) -- (0.92643896,0.59305216) -- (0.93221649,0.58392843) -- (0.93790418,0.57474842) -- (0.94350147,0.56551302) -- (0.94900782,0.55622311) -- (0.95442271,0.54687959) -- (0.95974561,0.53748337) -- (0.96497600,0.52803533) -- (0.97011339,0.51853641) -- (0.97515728,0.50898751) -- (0.98010718,0.49938955) -- (0.98496261,0.48974346) -- (0.98972311,0.48005016) -- (0.99438822,0.47031060) -- (0.99895749,0.46052571) -- (1.00343048,0.45069643) -- (1.00780675,0.44082372) -- (1.01208589,0.43090851) -- (1.01626749,0.42095178);
\draw[esblue] (1.01626749,0.42095178) -- (1.02035113,0.41095447) -- (1.02433643,0.40091755) -- (1.02822300,0.39084199) -- (1.03201047,0.38072876) -- (1.03569847,0.37057884) -- (1.03928665,0.36039320) -- (1.04277466,0.35017282) -- (1.04616217,0.33991869) -- (1.04944884,0.32963180) -- (1.05263437,0.31931314) -- (1.05571844,0.30896371) -- (1.05870076,0.29858449) -- (1.06158104,0.28817650) -- (1.06435900,0.27774073) -- (1.06703438,0.26727820) -- (1.06960691,0.25678990) -- (1.07207636,0.24627685) -- (1.07444247,0.23574007) -- (1.07670503,0.22518056) -- (1.07886381,0.21459935) -- (1.08091861,0.20399746) -- (1.08286922,0.19337591) -- (1.08471547,0.18273572) -- (1.08645717,0.17207791) -- (1.08809416,0.16140352) -- (1.08962627,0.15071358) -- (1.09105337,0.14000910) -- (1.09237530,0.12929114) -- (1.09359195,0.11856071) -- (1.09470320,0.10781885) -- (1.09570894,0.09706661) -- (1.09660907,0.08630501) -- (1.09740350,0.07553508) -- (1.09809217,0.06475788) -- (1.09867500,0.05397444) -- (1.09915194,0.04318580) -- (1.09952294,0.03239299) -- (1.09978796,0.02159706) -- (1.09994699,0.01079905) -- (1.10000000,-0.00000000);
\draw[esblue,dashed] (D) -- (F);
\draw[esblue,dashed] (F) -- (G);
\draw[selvage] (2.89348331,1.01451335) -- (2.73642379,0.94945717) -- (2.67136760,1.10651669);
\draw[inkiron] ($(B)!0.48340228!(F)!3pt!90:(F)$) -- ($(B)!0.48340228!(F)!3pt!-90:(F)$);
\draw[inkiron] ($(B)!0.51659772!(F)!3pt!90:(F)$) -- ($(B)!0.51659772!(F)!3pt!-90:(F)$);
\draw[inkiron] ($(D)!0.48340228!(F)!3pt!90:(F)$) -- ($(D)!0.48340228!(F)!3pt!-90:(F)$);
\draw[inkiron] ($(D)!0.51659772!(F)!3pt!90:(F)$) -- ($(D)!0.51659772!(F)!3pt!-90:(F)$);
\fill[inkiron] (A) circle (1pt);
\node at (3.90000000,4.28000000) {$A$};
\fill[inkiron] (B) circle (1pt);
\node at (-0.25000000,-0.25000000) {$B$};
\fill[inkiron] (C) circle (1pt);
\node at (8.25000000,-0.25000000) {$C$};
\fill[inkiron] (D) circle (1pt);
\node at (6.00685425,2.69314575) {$D$};
\fill[inkiron] (E) circle (1pt);
\node at (7.12842712,3.00842712) {$E$};
\fill[inkiron] (F) circle (1pt);
\node at (3.31370850,-0.35000000) {$F$};
\fill[inkiron] (G) circle (1pt);
\node at (2.47842712,1.27157288) {$G$};
\end{tikzpicture}}

\newcommand{\esdiExampleHand}{%
\begin{tikzpicture}[scale=.8,line width=.7pt,every node/.style={font=\small,inner sep=2pt}]
\coordinate (A) at (4.00000000,4.00000000);
\coordinate (B) at (0.00000000,0.00000000);
\coordinate (C) at (8.00000000,0.00000000);
\coordinate (D) at (5.65685425,2.34314575);
\coordinate (E) at (6.82842712,2.82842712);
\coordinate (F) at (2.82842712,1.17157288);
\draw[inkiron] (A) -- (B);
\draw[inkiron] (B) -- (C);
\draw[inkiron] (A) -- (C);
\draw[inkiron] (B) -- (E);
\draw[inkiron] (C) -- (E);
\draw[selvage] (3.87979185,3.87979185) -- (4.00000000,3.75958369) -- (4.12020815,3.87979185);
\draw[selvage] (6.67136760,2.76337094) -- (6.73642379,2.60631142) -- (6.89348331,2.67136760);
\draw[inkiron] ($(A)!0.50000000!(B)!3pt!90:(B)$) -- ($(A)!0.50000000!(B)!3pt!-90:(B)$);
\draw[inkiron] ($(A)!0.50000000!(C)!3pt!90:(C)$) -- ($(A)!0.50000000!(C)!3pt!-90:(C)$);
\draw[esblue] (0.77781746,0.77781746) -- (0.78541606,0.77014389) -- (0.79293896,0.76239610) -- (0.80038543,0.75457482) -- (0.80775476,0.74668082) -- (0.81504624,0.73871485) -- (0.82225916,0.73067768) -- (0.82939283,0.72257009) -- (0.83644656,0.71439285) -- (0.84341968,0.70614676) -- (0.85031150,0.69783261) -- (0.85712137,0.68945120) -- (0.86384862,0.68100334) -- (0.87049262,0.67248985) -- (0.87705272,0.66391154) -- (0.88352828,0.65526923) -- (0.88991869,0.64656378) -- (0.89622333,0.63779600) -- (0.90244159,0.62896676) -- (0.90857287,0.62007689) -- (0.91461657,0.61112726) -- (0.92057213,0.60211872) -- (0.92643896,0.59305216) -- (0.93221649,0.58392843) -- (0.93790418,0.57474842) -- (0.94350147,0.56551302) -- (0.94900782,0.55622311) -- (0.95442271,0.54687959) -- (0.95974561,0.53748337) -- (0.96497600,0.52803533) -- (0.97011339,0.51853641) -- (0.97515728,0.50898751) -- (0.98010718,0.49938955) -- (0.98496261,0.48974346) -- (0.98972311,0.48005016) -- (0.99438822,0.47031060) -- (0.99895749,0.46052571) -- (1.00343048,0.45069643) -- (1.00780675,0.44082372) -- (1.01208589,0.43090851) -- (1.01626749,0.42095178);
\draw[esblue] (1.01626749,0.42095178) -- (1.02035113,0.41095447) -- (1.02433643,0.40091755) -- (1.02822300,0.39084199) -- (1.03201047,0.38072876) -- (1.03569847,0.37057884) -- (1.03928665,0.36039320) -- (1.04277466,0.35017282) -- (1.04616217,0.33991869) -- (1.04944884,0.32963180) -- (1.05263437,0.31931314) -- (1.05571844,0.30896371) -- (1.05870076,0.29858449) -- (1.06158104,0.28817650) -- (1.06435900,0.27774073) -- (1.06703438,0.26727820) -- (1.06960691,0.25678990) -- (1.07207636,0.24627685) -- (1.07444247,0.23574007) -- (1.07670503,0.22518056) -- (1.07886381,0.21459935) -- (1.08091861,0.20399746) -- (1.08286922,0.19337591) -- (1.08471547,0.18273572) -- (1.08645717,0.17207791) -- (1.08809416,0.16140352) -- (1.08962627,0.15071358) -- (1.09105337,0.14000910) -- (1.09237530,0.12929114) -- (1.09359195,0.11856071) -- (1.09470320,0.10781885) -- (1.09570894,0.09706661) -- (1.09660907,0.08630501) -- (1.09740350,0.07553508) -- (1.09809217,0.06475788) -- (1.09867500,0.05397444) -- (1.09915194,0.04318580) -- (1.09952294,0.03239299) -- (1.09978796,0.02159706) -- (1.09994699,0.01079905) -- (1.10000000,-0.00000000);
\draw[esblue,dashed] (A) -- (F);
\draw[esblue,dashed] (A) -- (E);
\draw[selvage] (3.87754130,3.70435855) -- (4.17318275,3.58189985) -- (4.29564145,3.87754130);
\draw[inkiron] ($(A)!0.48203476!(F)!3pt!90:(F)$) -- ($(A)!0.48203476!(F)!3pt!-90:(F)$);
\draw[inkiron] ($(A)!0.51796524!(F)!3pt!90:(F)$) -- ($(A)!0.51796524!(F)!3pt!-90:(F)$);
\draw[inkiron] ($(A)!0.48203476!(E)!3pt!90:(E)$) -- ($(A)!0.48203476!(E)!3pt!-90:(E)$);
\draw[inkiron] ($(A)!0.51796524!(E)!3pt!90:(E)$) -- ($(A)!0.51796524!(E)!3pt!-90:(E)$);
\fill[inkiron] (A) circle (1pt);
\node at (3.90000000,4.28000000) {$A$};
\fill[inkiron] (B) circle (1pt);
\node at (-0.25000000,-0.25000000) {$B$};
\fill[inkiron] (C) circle (1pt);
\node at (8.25000000,-0.25000000) {$C$};
\fill[inkiron] (D) circle (1pt);
\node at (6.00685425,2.69314575) {$D$};
\fill[inkiron] (E) circle (1pt);
\node at (7.12842712,3.00842712) {$E$};
\fill[inkiron] (F) circle (1pt);
\node at (2.82842712,0.77157288) {$F$};
\end{tikzpicture}}

\newcommand{\esdiExampleReflect}{%
\begin{tikzpicture}[scale=.8,line width=.7pt,every node/.style={font=\small,inner sep=2pt}]
\coordinate (A) at (4.00000000,4.00000000);
\coordinate (B) at (0.00000000,0.00000000);
\coordinate (C) at (8.00000000,0.00000000);
\coordinate (D) at (5.65685425,2.34314575);
\coordinate (E) at (6.82842712,2.82842712);
\coordinate (F) at (5.65685425,5.65685425);
\draw[inkiron] (A) -- (B);
\draw[inkiron] (B) -- (C);
\draw[inkiron] (A) -- (C);
\draw[inkiron] (B) -- (E);
\draw[inkiron] (C) -- (E);
\draw[selvage] (3.87979185,3.87979185) -- (4.00000000,3.75958369) -- (4.12020815,3.87979185);
\draw[selvage] (6.67136760,2.76337094) -- (6.73642379,2.60631142) -- (6.89348331,2.67136760);
\draw[inkiron] ($(A)!0.50000000!(B)!3pt!90:(B)$) -- ($(A)!0.50000000!(B)!3pt!-90:(B)$);
\draw[inkiron] ($(A)!0.50000000!(C)!3pt!90:(C)$) -- ($(A)!0.50000000!(C)!3pt!-90:(C)$);
\draw[esblue] (0.77781746,0.77781746) -- (0.78541606,0.77014389) -- (0.79293896,0.76239610) -- (0.80038543,0.75457482) -- (0.80775476,0.74668082) -- (0.81504624,0.73871485) -- (0.82225916,0.73067768) -- (0.82939283,0.72257009) -- (0.83644656,0.71439285) -- (0.84341968,0.70614676) -- (0.85031150,0.69783261) -- (0.85712137,0.68945120) -- (0.86384862,0.68100334) -- (0.87049262,0.67248985) -- (0.87705272,0.66391154) -- (0.88352828,0.65526923) -- (0.88991869,0.64656378) -- (0.89622333,0.63779600) -- (0.90244159,0.62896676) -- (0.90857287,0.62007689) -- (0.91461657,0.61112726) -- (0.92057213,0.60211872) -- (0.92643896,0.59305216) -- (0.93221649,0.58392843) -- (0.93790418,0.57474842) -- (0.94350147,0.56551302) -- (0.94900782,0.55622311) -- (0.95442271,0.54687959) -- (0.95974561,0.53748337) -- (0.96497600,0.52803533) -- (0.97011339,0.51853641) -- (0.97515728,0.50898751) -- (0.98010718,0.49938955) -- (0.98496261,0.48974346) -- (0.98972311,0.48005016) -- (0.99438822,0.47031060) -- (0.99895749,0.46052571) -- (1.00343048,0.45069643) -- (1.00780675,0.44082372) -- (1.01208589,0.43090851) -- (1.01626749,0.42095178);
\draw[esblue] (1.01626749,0.42095178) -- (1.02035113,0.41095447) -- (1.02433643,0.40091755) -- (1.02822300,0.39084199) -- (1.03201047,0.38072876) -- (1.03569847,0.37057884) -- (1.03928665,0.36039320) -- (1.04277466,0.35017282) -- (1.04616217,0.33991869) -- (1.04944884,0.32963180) -- (1.05263437,0.31931314) -- (1.05571844,0.30896371) -- (1.05870076,0.29858449) -- (1.06158104,0.28817650) -- (1.06435900,0.27774073) -- (1.06703438,0.26727820) -- (1.06960691,0.25678990) -- (1.07207636,0.24627685) -- (1.07444247,0.23574007) -- (1.07670503,0.22518056) -- (1.07886381,0.21459935) -- (1.08091861,0.20399746) -- (1.08286922,0.19337591) -- (1.08471547,0.18273572) -- (1.08645717,0.17207791) -- (1.08809416,0.16140352) -- (1.08962627,0.15071358) -- (1.09105337,0.14000910) -- (1.09237530,0.12929114) -- (1.09359195,0.11856071) -- (1.09470320,0.10781885) -- (1.09570894,0.09706661) -- (1.09660907,0.08630501) -- (1.09740350,0.07553508) -- (1.09809217,0.06475788) -- (1.09867500,0.05397444) -- (1.09915194,0.04318580) -- (1.09952294,0.03239299) -- (1.09978796,0.02159706) -- (1.09994699,0.01079905) -- (1.10000000,-0.00000000);
\draw[esblue,dashed] (E) -- (F);
\draw[esblue,dashed] (B) -- (F);
\draw[inkiron] ($(C)!0.48203476!(E)!3pt!90:(E)$) -- ($(C)!0.48203476!(E)!3pt!-90:(E)$);
\draw[inkiron] ($(C)!0.51796524!(E)!3pt!90:(E)$) -- ($(C)!0.51796524!(E)!3pt!-90:(E)$);
\draw[inkiron] ($(E)!0.48203476!(F)!3pt!90:(F)$) -- ($(E)!0.48203476!(F)!3pt!-90:(F)$);
\draw[inkiron] ($(E)!0.51796524!(F)!3pt!90:(F)$) -- ($(E)!0.51796524!(F)!3pt!-90:(F)$);
\fill[inkiron] (A) circle (1pt);
\node at (3.90000000,4.28000000) {$A$};
\fill[inkiron] (B) circle (1pt);
\node at (-0.25000000,-0.25000000) {$B$};
\fill[inkiron] (C) circle (1pt);
\node at (8.25000000,-0.25000000) {$C$};
\fill[inkiron] (D) circle (1pt);
\node at (6.00685425,2.69314575) {$D$};
\fill[inkiron] (E) circle (1pt);
\node at (7.12842712,3.00842712) {$E$};
\fill[inkiron] (F) circle (1pt);
\node at (5.65685425,5.90685425) {$F$};
\end{tikzpicture}}

% END GENERATED: chapter-diagrams
\begin{document}
\loomcover{geometry Models}{}{Hence}{2026年10月}
\clearpage
{\small\tableofcontents}
\clearpage
% BEGIN GENERATED: chapter-body
\phantomsection
\part*{第一章\quad 夹半角}
\addcontentsline{toc}{part}{第一章 夹半角}
\renewcommand{\theHsection}{01-half-angle.\arabic{section}}
\setcounter{section}{0}
\section*{四种常见构型}
\phantomsection
\addcontentsline{toc}{section}{四种常见构型}
四图中的 \(ABCD\) 均为正方形，各动点在指定边的内部。识别时先找等长邻边和直角，再看 \(45^\circ\) 在顶点还是交点。原图与辅助线图并排，辅助线用虚线表示。

\block{图 A 顶点夹半角：完整证明}
\textbf{条件与结论}\quad \(E\in BC\)、\(F\in CD\)，\(\angle EAF=45^\circ\)。则 \(EF=BE+DF\)；反过来，\(EF=BE+DF\) 也能推出 \(\angle EAF=45^\circ\)。
\diagram{\figIntroA\hspace{1.2cm}\figIntroASolve}

\think 把 \(BE\) 搬到 \(CD\) 的延长线上，使 \(BE+DF\) 成为一条线段，再用半角证明这条线段等于 \(EF\)。
\solve 延长 \(CD\) 至 \(G\)，使 \(DG=BE\)，连接 \(AG\)。点序为 \(C,F,D,G\)。这相当于把 \(\triangle ABE\) 绕 \(A\) 逆时针旋转 \(90^\circ\) 到 \(\triangle ADG\)。
\begin{enumerate}
\item 在 \(\triangle ABE\) 与 \(\triangle ADG\) 中，\(AB=AD\)、\(BE=DG\)、\(\angle ABE=\angle ADG=90^\circ\)，所以 \(\triangle ABE\cong\triangle ADG\)（SAS）。于是 \(AE=AG\)，\(\angle BAE=\angle DAG\)。
\item 因为 \(\angle BAE+\angle EAD=90^\circ\)，所以
\[
\angle EAG=\angle EAD+\angle DAG=90^\circ.
\]
\(AF\) 在 \(\angle EAG\) 内，而 \(\angle EAF=45^\circ\)，故 \(\angle FAG=45^\circ\)。由 \(AE=AG\)、\(AF\) 公共、\(\angle EAF=\angle GAF\)，得 \(\triangle EAF\cong\triangle GAF\)（SAS），因此 \(EF=FG\)。
\item 由 \(F,D,G\) 的点序，\(EF=FG=FD+DG=DF+BE\)，结论成立。
\end{enumerate}
\insight{等长邻边提供第一次全等，半角提供第二次全等，最后按点序写线段和。}

\clearpage
\block{图 A 的逆命题与相应结论}
\textbf{逆命题的证明}\quad 仍作 \(DG=BE\)，第一组全等仍给出 \(AE=AG\)、\(\angle EAG=90^\circ\)。若 \(EF=BE+DF\)，则 \(EF=FG\)。结合 \(AF\) 公共，\(\triangle EAF\cong\triangle GAF\)（SSS），所以 \(\angle EAF=\angle GAF\)。两角相加为 \(90^\circ\)，故 \(\angle EAF=45^\circ\)。

设 \(AB=s\)、\(BE=e\)、\(DF=f\)，其中 \(0<e,f<s\)。由线段和可以继续得到：
\begin{enumerate}
\item \textbf{定周长。} \(CE=s-e\)、\(CF=s-f\)，所以
\[
EF=e+f\quad\Longrightarrow\quad
C_{\triangle CEF}=(s-e)+(s-f)+(e+f)=2s.
\]
\item \textbf{长度关系。} 在 \(\mathrm{Rt}\triangle CEF\) 中，由勾股定理，
\[
(s-e)^2+(s-f)^2=(e+f)^2.
\]
\item \textbf{已知 \(AB\)、\(BE\) 求另两段。} 由上式解得
\[
\boxed{DF=\frac{s(s-e)}{s+e}},\qquad
\boxed{CF=s-DF=\frac{2se}{s+e}}.
\]
\end{enumerate}

\par\bigskip
\noindent\begin{minipage}{\linewidth}
\block{图 B 交叉角：平行转移}
\(E\in BC\)、\(F\in AB\)，\(AE\) 与 \(CF\) 交于 \(P\)，\(\angle EPC=45^\circ\)。过 \(A\) 作 \(AG\parallel CF\)，交 \(CD\) 于 \(G\)，并连接 \(EG\)。
\diagram{\figIntroB\hspace{1.2cm}\figIntroBSolve}
\[
AG\parallel CF\quad\Longrightarrow\quad\angle EAG=\angle EPC=45^\circ.
\]
\[
\left.\begin{gathered}
AD=CB,\quad\angle ADG=\angle CBF=90^\circ\\
\angle DAG=\angle BCF
\end{gathered}\right\}
\ \Longrightarrow\ \triangle ADG\cong\triangle CBF\ (\mathrm{ASA})
\ \Longrightarrow\ DG=BF.
\]
\[
\text{转为图 A}\quad\Longrightarrow\quad
\boxed{EG=BE+DG=BE+BF}.
\]
若 \(AB=s\)、\(BE=e\)、\(BF=f\)，则由图 A 的勾股等式化简得 \(\boxed{ef+s(e+f)=s^2}\)。
\end{minipage}

\clearpage
\noindent\begin{minipage}{\linewidth}
\block{图 C 交叉角：垂直转移}
\(E\in BC\)、\(F\in AD\)，\(AE\) 与 \(BF\) 交于 \(P\)，\(\angle APF=45^\circ\)。过 \(A\) 作 \(AG\perp BF\)，交 \(CD\) 于 \(G\)，并连接 \(EG\)。
\diagram{\figIntroC\hspace{1.2cm}\figIntroCSolve}
\[
AG\perp BF,\quad\angle APF=45^\circ
\quad\Longrightarrow\quad\angle EAG=90^\circ-45^\circ=45^\circ.
\]
\[
\left.\begin{gathered}
AD=BA,\quad\angle ADG=\angle BAF=90^\circ\\
\angle DAG=\angle ABF
\end{gathered}\right\}
\ \Longrightarrow\ \triangle ADG\cong\triangle BAF\ (\mathrm{ASA})
\ \Longrightarrow\ DG=AF.
\]
其中 \(\angle DAG=\angle ABF\)，因为这两个锐角的两边分别垂直。
\[
\text{转为图 A}\quad\Longrightarrow\quad
\boxed{EG=BE+DG=BE+AF}.
\]
若 \(AB=s\)、\(BE=e\)、\(AF=f\)，则 \(\boxed{ef+s(e+f)=s^2}\)。
\end{minipage}

\par\bigskip
\noindent\begin{minipage}{\linewidth}
\block{图 D 两条线同时平移}
\(E\in AD\)、\(F\in BC\)、\(M\in AB\)、\(N\in CD\)，\(EF\) 与 \(MN\) 交于 \(P\)，\(\angle FPN=45^\circ\)。按图示位置，\(BF>AE\)、\(DN>AM\)。过 \(A\) 作 \(AT\parallel EF\)、\(AU\parallel MN\)，分别交 \(BC\)、\(CD\) 于 \(T,U\)，并连接 \(TU\)。
\diagram{\figIntroD\hspace{1.2cm}\figIntroDSolve}
\[
\begin{aligned}
AE\parallel TF,\ AT\parallel EF
&\ \Longrightarrow\ AEFT\text{ 为平行四边形}\\
&\ \Longrightarrow\ TF=AE\ \Longrightarrow\ BT=BF-AE;\\[5pt]
AM\parallel NU,\ AU\parallel MN
&\ \Longrightarrow\ AMNU\text{ 为平行四边形}\\
&\ \Longrightarrow\ NU=AM\ \Longrightarrow\ DU=DN-AM.
\end{aligned}
\]
\[
AT\parallel PF,\ AU\parallel PN
\quad\Longrightarrow\quad\angle TAU=45^\circ
\quad\Longrightarrow\quad\text{转为图 A}.
\]
\[
\boxed{TU=BT+DU=(BF-AE)+(DN-AM)}.
\]
若 \(AB=s\)、\(BT=e\)、\(DU=f\)，仍有 \(ef+s(e+f)=s^2\)。点序改变时，要重新判断射线方向及线段加减。
\end{minipage}
\clearpage

\section{一 正方形内的旋转}
\begin{problem}{题 1 旋转拼出线段和}
正方形 \(ABCD\) 中，\(E\)、\(F\) 分别在边 \(BC\)、\(CD\) 的内部，\(\angle EAF=45^\circ \)。证明 \(EF=BE+DF\)。再思考：若只知道 \(EF=BE+DF\)，能否推出 \(\angle EAF=45^\circ \)？
\diagram{\figOne}
\hint{延长 \(CD\) 至 \(G\)，使 \(DG=BE\)，再比较 \(AE\) 与 \(AG\)。}
\end{problem}
\begin{problem}{题 2 四个条件相互推导}
正方形 \(ABCD\) 中，\(M\)、\(N\) 分别在 \(AD\)、\(CD\) 的内部。证明下列四个条件彼此等价：\((1) \angle MBN=45^\circ \)；\((2) MN=AM+CN\)；\((3) \angle AMB=\angle BMN\)；\((4) \angle BNM=\angle BNC\)。
\diagram{\figTwo}
\hint{先由\((1)\)证明\((2)\)。再从 \(B\) 向 \(MN\) 作垂线，观察等角与等距的联系。}
\end{problem}
\clearpage
\begin{problem}{题 3 从等角恢复半角}
正方形 \(ABCD\) 中，\(M\in AD\)、\(N\in CD\)，\(\angle BMN=\angle MBC\)，\(AB=9\)，\(DM=6\)，求 \(CN\)。
\diagram{\figThree}
\hint{由 \(AD\parallel BC\)，把已知等角转写为题 \(2\) 中的一个条件。}
\end{problem}
\begin{problem}{题 4 对角线上的垂直与等长}
正方形 \(ABCD\) 中，\(E\) 为对角线 \(AC\) 内部的一点，\(N\) 在 \(CD\) 上，且 \(BE\perp EN\)。证明 \(BE=EN\)。
\diagram{\figFour}
\hint{连接 \(DE\)。关于对角线 \(AC\) 对称的两个顶点是哪两个？}
\end{problem}
\clearpage
\section{二 从三角形补出模型}
\begin{problem}{题 5 同一道题的三种构造}
\(\triangle ABC\) 中，\(D\) 为 \(BC\) 内部的一点，\(AD\perp BC\)，\(\angle BAC=45^\circ \)，\(BD=3\)，\(AD=6\)，求 \(CD\)。分别尝试旋转、两次翻折、延长构造等腰直角三角形。
\diagram{\figFive}
\hint{在直线 \(BC\) 上取 \(E\)、\(F\)，使 \(DE=DF=AD\)，且 \(E\)、\(B\)、\(D\)、\(C\)、\(F\) 按此顺序排列。}
\end{problem}
\section{三 矩形与交叉角}
\begin{problem}{题 6 矩形中截出正方形}
矩形 \(ABCD\) 中，\(AD=6\)，\(AB=12\)，\(F\in CD\)，\(DF=3\)，\(E\in BC\)，\(\angle EAF=45^\circ \)，求 \(BE\)。
\diagram{\figSix}
\hint{取 \(AB\) 的中点 \(M\)，过 \(M\) 作 \(MQ\parallel AD\)，\(Q\in DC\)。\(AE\) 与 \(MQ\) 交于 \(N\)。}
\end{problem}
\clearpage
\begin{problem}{题 7 把交叉角搬到顶点}
直角三角形 \(ABC\) 中，\(\angle C=90^\circ \)，\(AC=6\)，\(BC=8\)，\(D\) 是 \(AC\) 的中点，\(E\in BC\)，\(BD\) 与 \(AE\) 交于 \(F\)。若 \(\angle BFE=45^\circ \)，求 \(CE\)。
\diagram{\figSeven}
\hint{过 \(A\) 作 \(AG\perp BD\)，\(G\) 在 \(BC\) 向 \(C\) 外的延长线上，把交叉角转为 \(\angle GAE\)。}
\end{problem}
\begin{problem}{题 8 相等的两段补足旋转}
矩形 \(ABCD\) 中，\(E\)、\(F\) 分别在 \(BC\)、\(CD\) 的内部，\(CE=CF\)，\(DF=3\)，\(BE=4\)，\(\angle EAF=45^\circ \)，求 \(AB\)。
\diagram{\figEight}
\hint{把 \(AE\) 绕 \(A\) 逆时针旋转 \(90^\circ \)。旋转后的端点与 \(F\) 的水平距离、竖直距离能否直接算出？}
\end{problem}
\clearpage
\section{四 位置变化与分类讨论}
\begin{problem}{题 9 从线段和到线段差}
四边形 \(ABCD\) 的四条边相等，四个角都是直角，连接 \(AM\)、\(AN\)、\(MN\)。

\((1)M\)、\(N\) 分别在 \(BC\)、\(CD\) 的内部，\(\angle MAN=45^\circ \)。证明 \(MN=BM+DN\)。（原题 \(6\) 分）

\((2)M\) 在 \(BC\) 向 \(C\) 外的延长线上，\(N\) 在 \(CD\) 向 \(D\) 外的延长线上，仍有 \(\angle MAN=45^\circ \)。判断 \(BM\)、\(DN\)、\(MN\) 的数量关系，并证明。（原题 \(6\) 分）
\diagram{\figNine}
\hint{两问采用同一次旋转，再观察三个共线点的顺序。}
\end{problem}
\begin{problem}{题 10 一个条件对应两种构型}
\(\triangle ABC\) 中，\(D\) 在 \(BC\) 上，\(\angle ABC=45^\circ \)，\(\angle CAD=45^\circ \)，\(BD=4\)，\(BC=9\)，求 \(AB\) 的长。
\diagram{\figTen}
\hint{过 \(A\) 作 \(AH\perp BC\)，垂足 \(H\) 在直线 \(BC\) 上。先利用 \(\angle B=45^\circ \)，再讨论 \(H\) 与 \(D\) 的位置。}
\end{problem}

\clearpage
\section{集中解析}

\needspace{10\baselineskip}
\begin{solution}{题 1 · 解析}
\think 用旋转连接半角与线段和，用垂线连接等角与等距。
\diagram{\figTwoEquivalence}

\solve 将 \(M\) 绕 \(B\) 顺时针旋转 \(90^\circ\) 到 \(P\)，得 \(CP=AM\)、\(BP=BM\)、\(\angle MBP=90^\circ\)；另作 \(BQ\perp MN\)。两种辅助线分图表示。
\diagram{\figTwoSolve}

\solvehead{旋转：半角与线段和相互推出}
\((1)\) 使 \(BN\) 平分 \(\angle MBP\)，因此
\[
\triangle MBN\cong\triangle PBN\ (\mathrm{SAS})
\ \Longrightarrow\ \relax
\begin{cases}
MN=PN=AM+CN &(2),\\
\angle BMN=\angle BPN=\angle AMB &(3),\\
\angle BNM=\angle BNP=\angle BNC &(4).
\end{cases}
\]
反过来，\((2)\) 给出 \(MN=PN\)，结合 \(BM=BP\)、\(BN\) 公共，
\[
\triangle MBN\cong\triangle PBN\ (\mathrm{SSS})
\quad\Longrightarrow\quad
\angle MBN=\tfrac12\angle MBP=45^\circ\quad(1).
\]

\solvehead{垂线：两个等角条件相互推出}
\(B\) 在 \(\angle AMN\)、\(\angle CNM\) 内，角平分线可用到两边的等距来判断。
\[
\begin{aligned}
(3)&\ \Longrightarrow\ \mathrm{Rt}\triangle BQM\cong\mathrm{Rt}\triangle BAM
\ \Longrightarrow\ BQ=BA=BC\ \Longrightarrow\ (4);\\
(4)&\ \Longrightarrow\ BQ=BC=BA
\ \Longrightarrow\ BM\text{ 平分 }\angle AMN\ \Longrightarrow\ (3).
\end{aligned}
\]
第一行的全等使用公共斜边 \(BM\) 与 \(M\) 处相等的锐角。

\solvehead{等角：恢复顶点处的半角}
令 \(\alpha=\angle ABM\)、\(\beta=\angle NBC\)、\(\theta=\angle MBN\)。由 \((3)(4)\)，
\[
\left.\begin{aligned}
\theta&=\alpha+\beta &&(\triangle BMN),\\
\theta+\alpha+\beta&=90^\circ &&(\angle ABC)
\end{aligned}\right\}
\quad\Longrightarrow\quad \theta=45^\circ\quad(1).
\]
\insight{旋转看全等，垂线看等距；最后用三角形内角和恢复半角。}
\end{solution}
\needspace{10\baselineskip}
\begin{solution}{题 2 · 解析}
\think 题面没有写 \(45^\circ \)，但已给出题 \(1\) 的等价条件。

\solve 因 \(AD\parallel BC\)，\(\angle AMB=\angle MBC\)。已知 \(\angle BMN=\angle MBC\)，故 \(\angle AMB=\angle BMN\)，由题 \(1\)，\(\angle MBN=45^\circ \)且 \(MN=AM+CN\)。\(AM=9-6=3\)。设 \(CN=x\)，则 \(DN=9-x\)，\(MN=3+x\)。在 \(\mathrm{Rt}\triangle MDN\) 中，\(6^2+(9-x)^2=(3+x)^2\)，解得 \(x=\frac{9}{2}\)。所以 \(CN=\frac{9}{2}\)。

\insight{没有出现半角数值时，先找模型的等价条件。}
\end{solution}
\needspace{10\baselineskip}
\begin{solution}{题 3 · 解析}
\think \(E\) 在正方形对角线上，天然具有对称性；先把 \(BE\) 变成 \(DE\)，再用垂直关系证明 \(DE=EN\)。

\diagram{\figFourSolve}
\solve 连接 \(DE\)。正方形关于 \(AC\) 对称，\(B\) 与 \(D\) 互为对称点，\(E\) 位于对称轴上，所以 \(BE=DE\)、\(\angle ABE=\angle ADE\)。由 \(BE\perp EN\)、\(BC\perp ND\)，且按图中射线方向，有 \(\angle DNE=\angle EBC\)。又 \(\angle EDN=90^\circ -\angle ADE=90^\circ -\angle ABE=\angle EBC\)，所以 \(\angle EDN=\angle DNE\)，\(DE=EN\)。因此 \(BE=EN\)。

\insight{对角线给出等长，垂直给出等角，合起来构成等腰三角形。}
\end{solution}
\needspace{10\baselineskip}
\begin{solution}{题 5 · 解析}
\think 长边是短边的两倍，可以截出一个正方形，先在小图中求长，再回到原矩形。

\diagram{\figSixSolve}
\solve 取 \(AB\) 中点 \(M\)，作 \(MQ\parallel AD\)，交 \(DC\) 于 \(Q\)，\(AE\) 与 \(MQ\) 交于 \(N\)。\(AMQD\) 是边长 \(6\) 的正方形，\(DQ=6\)、\(FQ=3\)。设 \(MN=x\)，图 A 给出 \(FN=DF+MN=3+x\)，而 \(NQ=6-x\)。在 \(\mathrm{Rt}\triangle FQN\) 中，\((3+x)^2=3^2+(6-x)^2\)，得 \(x=2\)。\(\triangle AMN\sim \triangle ABE\)，\(\frac{AB}{AM}=2\)，因此 \(BE=2MN=4\)。

\insight{矩形先截正方形，局部模型求出后再用相似放回整体。}
\end{solution}
\needspace{16\baselineskip}
\begin{solution}{题 6 · 解析}
\think 交点处的 \(45^\circ \)可通过一条垂线搬到 \(A\)，随后变成例题 \(4\) 的底边高模型。

\diagram{\figSevenSolve}
\solve 过 \(A\) 作 \(AG\perp BD\)，交 \(BC\) 向 \(C\) 外的延长线于 \(G\)。\(\angle BFE=45^\circ \)，所以按 \(G\)、\(C\)、\(E\) 的点序，\(\angle GAE=45^\circ \)。\(\mathrm{Rt}\triangle ACG\) 与 \(\mathrm{Rt}\triangle BCD\) 相似，\(\frac{CG}{AC}=\frac{CD}{BC}=\frac{3}{8}\)，故 \(CG=\frac{9}{4}\)。设 \(CE=x\)。在 \(\triangle AGE\) 中，\(AC\) 为高，\(AC=6\)、\(CG=\frac{9}{4}\)。按例题 \(4\) 的延长旋转构造，\((\frac{9}{4}+x)^2=(6-\frac{9}{4})^2+(6-x)^2\)，展开并合并同类项，得 \(x=\frac{30}{11}\)。所以 \(CE=\frac{30}{11}\)。

\insight{交叉角优先用平行或垂直转移，再进入熟悉的模型。}
\end{solution}
\clearpage
\begin{solution}{题 7 · 解析}
\think 先用旋转得到相等的斜边，再从图上读出“一差一和”，用勾股定理求值。

\solve 设 \(CE=CF=x\)，则 \(AB=x+3\)、\(AD=x+4\)。将 \(E\) 绕 \(A\) 逆时针旋转 \(90^\circ\) 到 \(G\)。作 \(GH\parallel AD\)，交 \(CD\) 的延长线于 \(H\)；作 \(GT\perp AD\)，垂足为 \(T\)。
\diagram{\figEightSolve}

\solvehead{旋转全等：两条斜边相等}
\[
\left.\begin{gathered}
AE=AG,\quad AF\text{ 公共}\\
\angle EAF=\angle GAF=45^\circ
\end{gathered}\right\}
\ \xRightarrow{\mathrm{SAS}}\ \triangle EAF\cong\triangle GAF
\ \Longrightarrow\ EF=GF.
\]

\solvehead{读图：横向取差，竖向取和}
旋转给出 \(AT=AB=x+3\)、\(GT=BE=4\)；\(GTDH\) 为矩形，故
\[
\begin{aligned}
GH&=TD=AD-AT=(x+4)-(x+3)=1,\\
FH&=FD+DH=3+4=7.
\end{aligned}
\]

\solvehead{勾股：同一斜边的平方相等}
\[
\left.\begin{aligned}
EF^2&=CE^2+CF^2=2x^2,\\
GF^2&=GH^2+FH^2=1^2+7^2=50
\end{aligned}\right\}
\ \xRightarrow{EF=GF}\ 2x^2=50.
\]
\[
x>0\quad\Longrightarrow\quad x=5
\quad\Longrightarrow\quad \boxed{AB=x+3=8}.
\]
\insight{相等的 \(CE\)、\(CF\) 让横向的 \(x\) 抵消，竖向直接相加。}
\end{solution}
\needspace{16\baselineskip}
\begin{solution}{题 8 · 解析}
\think 辅助线完全相同，位置改变只影响共线线段的加减。

\diagram{\figNineSolve}
\solve \((1)\)将 \(M\) 绕 \(A\) 逆时针旋转 \(90^\circ \)到 \(P\)，则 \(P\) 在 \(CD\) 向 \(D\) 外的延长线上，\(DP=BM\)、\(AP=AM\)。\(\angle MAP=90^\circ \)，\(AN\) 平分此角，所以 \(\triangle MAN\cong \triangle PAN(\mathrm{SAS})\)，\(MN=PN\)。由 \(P\)、\(D\)、\(N\) 的点序，\(PN=PD+DN=BM+DN\)。

\((2)\)同样旋转 \(M\) 得到 \(P\)。此时 \(P\)、\(N\)、\(D\)、\(C\) 按此顺序共线。\(\angle MAP=90^\circ \)、\(\angle MAN=45^\circ \)，\(AN\) 仍在该直角内部并平分它，所以 \(\triangle MAN\cong \triangle PAN\)，\(MN=PN\)。由点序，\(PN=PD-ND=BM-DN\)，即 \(BM=MN+DN\)。这里 \(BM>DN\)，不能把两条正长度的和照搬到外部构型。

\insight{先证明全等，再根据点序写加法或减法。}
\end{solution}
\clearpage
\begin{solution}{题 9 · 解析}
\think 旋转后得到一个 \(3\!:\!4\!:\!5\) 的直角三角形，再按 \(K\) 在 \(C\) 的哪一侧分两种情况。

\solve 作 \(AH\perp BC\)。\(\angle B=45^\circ\) 使 \(H\) 在 \(B\) 向 \(C\) 的射线上，设 \(AH=BH=t>0\)。将 \(\triangle ABD\) 绕 \(A\) 逆时针旋转 \(90^\circ\) 到 \(\triangle AKG\)，\(B,D\) 分别对应 \(K,G\)。
\diagram{\figTenSolve}

\solvehead{旋转：搬出已知的直角边}
\[
\begin{aligned}
\triangle ABH\text{ 为等腰直角三角形}
&\ \Longrightarrow\ BH=HK=t\ \Longrightarrow\ BK=2t;\\
BD\xrightarrow[\text{绕 }A]{90^\circ}KG
&\ \Longrightarrow\ KG\perp BC,\quad KG=BD=4.
\end{aligned}
\]

\solvehead{半角全等：搬出斜边}
\[
\left.\begin{gathered}
AD=AG,\quad AC\text{ 公共}\\
\angle DAC=\angle CAG=45^\circ
\end{gathered}\right\}
\ \xRightarrow{\mathrm{SAS}}\ \triangle DAC\cong\triangle GAC
\ \Longrightarrow\ GC=DC=5.
\]

\solvehead{勾股与分类：先求一段，再分加减}
\[
\mathrm{Rt}\triangle GKC\quad\Longrightarrow\quad
CK^2=GC^2-KG^2=5^2-4^2=9\quad\Longrightarrow\quad CK=3.
\]
由 \(BK=2t\)、\(AB=\sqrt2\,t\)，两种位置分别给出：
\[
\begin{array}{c@{\qquad\qquad}c}
K\text{ 在 }C\text{ 左侧}&K\text{ 在 }C\text{ 右侧}\\
\Downarrow&\Downarrow\\
BK=9-3=6&BK=9+3=12\\
\Downarrow&\Downarrow\\
t=3&t=6\\
\Downarrow&\Downarrow\\
\boxed{AB=3\sqrt2}&\boxed{AB=6\sqrt2}
\end{array}
\]
两种构型都满足原条件；若限定 \(H\) 在 \(DC\) 内部，才只保留 \(6\sqrt2\)。
\insight{先用旋转固定 \(CK=3\)，再由点序决定 \(BK\) 用加法还是减法。}
\end{solution}

\clearpage
\section{模型总结与拓展}
\ifteachernotes
\block{方法卡片}

长什么样：同一个顶点处有两条等长邻边夹 \(90^\circ \)，其中另两条射线夹 \(45^\circ \)；\(45^\circ \)也可能藏在两线交点处。

怎么想：先把 \(45^\circ \)搬到合适的顶点，再把分散的长度搬到同一直线上。

几步模板：确认等长与角度条件 → 转移交叉角或补出正方形 → 旋转或翻折得到全等 → 按点序写线段和差 → 用勾股或相似求值。

易错警示：矩形缺少邻边相等的条件，必须先截形或另作构造；点在延长线上时，需要重判线段加减。

常用添线：旋转适合等长邻边与半角；翻折适合半角顶点到同一直线有垂线；延长适合复制已知高，构造等腰直角三角形；矩形优先寻找局部正方形。

\block{拓展一 四个条件与定高}

在题 \(2\) 的内部构型中，四个条件任一成立，则 \(B\) 到 \(MN\) 的距离等于正方形边长 \(s\)。证明见题 \(2\)：从角平分关系得两组直角三角形全等。因此 \(S\triangle BMN=s\cdot MN/2\)。这是由斜边长求面积的快捷方法。

\block{拓展二 定周长与乘积关系}

仍用题 \(2\) 的构型，设 \(AM=m\)、\(CN=n\)、正方形边长为 \(s\)。由 \(MN=m+n\)，\(\triangle DMN\) 的周长为 \((s-m)+(s-n)+(m+n)=2s\)。

在 \(\mathrm{Rt}\triangle DMN\) 中，\((s-m)^2+(s-n)^2=(m+n)^2\)，整理得 \(mn+s(m+n)=s^2\)，即 \((s+m)(s+n)=2s^2\)。已知 \(m\) 可求 \(n=\frac{s(s-m)}{s+m}\)。该公式也可由题 \(2\) 逆向恢复半角，但使用时必须满足 \(0<m,n<s\)。

再设 \(L=MN=m+n\)。由 \(mn=s^2-sL>0\)，可得 \(L<s\)；由 \((m-n)^2=(L+2s)^2-8s^2\ge 0\)，得到 \(L\ge 2(\sqrt{2}-1)s\)。等号恰在 \(m=n=(\sqrt{2}-1)s\) 时成立。

\clearpage
\block{拓展三 底边高型的通式}

\(\triangle ABC\) 中，\(AD\perp BC\)，\(D\) 在 \(BC\) 内部，\(\angle BAC=45^\circ \)。设 \(AD=h\)、\(BD=p\)、\(CD=q\)，则 \(p,q>0\)，且 \(p,q<h\)。按题 \(5\) 的延长与旋转构造，\((p+q)^2=(h-p)^2+(h-q)^2\)，整理为 \(pq+h(p+q)=h^2\)，故 \(q=\frac{h(h-p)}{h+p}\)。反之，若 \(h,p,q>0\) 且该等式成立，同一构造可由边长恢复 \(45^\circ \)。

例如题 \(5\) 中 \(h=6\)、\(p=3\)，得 \(q=2\)；题 \(7\) 转化后 \(h=6\)、\(p=\frac{9}{4}\)，得 \(q=\frac{30}{11}\)。这里 \(D\) 必须在底边内部，不能把它直接套到垂足在外的题 \(10\) 第一种构型。

\block{拓展四 怎样判断应写和还是差}

题 \(9\) 的两问都先证明 \(MN=PN\)，再按点序表达 \(PN\)。\(P\)、\(D\)、\(N\) 时用 \(PD+DN\)；\(P\)、\(N\)、\(D\) 时用 \(PD-DN\)。解题时可以先写“目标线段＝辅助线段”，最后一步才换成题目中给出的长度。这样即使把图形转动、翻转，或把点移到延长线上，也不容易误套公式。

顶点夹半角构型中，若 \(AB=s\)、\(BE=e\)，那么 \(DF=\frac{s(s-e)}{s+e}\)，\(CF=\frac{2se}{s+e}\)，适用范围为 \(0<e<s\)。这两式把“已知 \(AB\)、\(BE\) 求 \(DF\)、\(CF\)”的手写提示补成了可直接使用的结论；其证明与拓展二的乘积关系相同。

\block{易错提醒}

\begin{itemize}
\item 矩形没有 \(AB=AD\)，不能直接套正方形的线段和。
\item 所有旋转都应注明中心、方向与角度，旋转点落在边内还是延长线上也要说明。
\item “半角”必须是两条指定射线的角，不能用交点附近任意一个 \(45^\circ \)替换。
\item 出现线段和或线段差，先判断共线点的顺序。
\item 图上看起来相等、垂直或居中，只有题目给出或证明过才能使用。
\item 求出数值后，检查长度为正、点在线段内部以及角度条件。
\end{itemize}
\else
\block{我的方法卡片}\workspace[10]
\fi


\clearpage
\phantomsection
\part*{第二章\quad 手拉手模型}
\addcontentsline{toc}{part}{第二章 手拉手模型}
\renewcommand{\theHsection}{02-hand-in-hand.\arabic{section}}
\setcounter{section}{-1}
\section{从一幅图开始}
本章从两道等腰直角三角形的题目出发，学习怎样把“手拉脚”“脚拉脚”转化为“手拉手”。先做题，再阅读解析；重点观察辅助点为什么取在这个位置，以及中点条件怎样参与第二次全等。题目与解析分开。



题 1 中有两个等腰直角三角形，但它们的直角顶点分别是 \(A\)、\(B\)。公共点 \(A\) 在一个三角形中是直角顶点，在另一个三角形中却是底角顶点。两个图形暂时不能绕同一个直角顶点一起旋转。

再看条件 \(BF=BD\)：\(B\) 是 \(DF\) 的中点。延长一条经过 \(B\) 的线段，往往能让 \(B\) 同时成为另一条线段的中点。这样既能出现“八字形”全等，又能补出新的等腰直角三角形。

本章的入手顺序是：\textbf{先认出两个等腰直角三角形，再检查公共点的身份；身份不一致，就用倍长或旋转补成一致。}

\section{手拉脚与中点}
\begin{problem}{题 1 共用一个底角顶点与一个直角顶点}
如图，在 \(\triangle ADE\) 和 \(\triangle ABC\) 中，\(AD=AE\)，\(AB=BC\)，\(\angle DAE=\angle ABC=90^\circ\)。延长 \(DB\) 到 \(F\)，使 \(BF=DB\)，连接 \(CF\)、\(CE\)。求证：\(CF=CE\)，\(CF\perp CE\)。

各点方向按图取定。尝试两种构造：延长 \(AB\)；或者把 \(AC\) 绕 \(A\) 旋转 \(90^\circ\)。

\diagram{\hhOne}

\hint{\(B\) 已经是 \(DF\) 的中点，能否让它也成为一条新线段的中点？}

\end{problem}
\clearpage
\section{脚拉脚与两次倍长}
\begin{problem}{题 2 两个底角顶点重合}
如图，\(\triangle ABC\) 和 \(\triangle DCE\) 都是等腰直角三角形，\(\angle BAC=\angle CDE=90^\circ\)，即 \(AB=AC\)、\(DC=DE\)。\(D\) 在 \(BA\) 向 \(A\) 外的延长线上，\(E\) 与 \(C\) 位于直线 \(AB\) 的两侧。连接 \(BE\)，求证：\(BE\perp BC\)。

\diagram{\hhTwo}

\hint{分别倍长 \(ED\)、\(BA\)，使公共点 \(C\) 成为两个新等腰直角三角形的直角顶点。}

\end{problem}
\clearpage

\clearpage
\section{集中解析}
\begin{solution}{题 1 · 解析}
\think

原图的两个等腰直角三角形是“手拉脚”：\(A\) 是 \(\triangle ADE\) 的直角顶点，却是 \(\triangle ABC\) 的底角顶点。\(BF=BD\) 提供了中心 \(B\) 的对称关系。把另一个点也关于 \(B\) 对称，就能利用一组“八字形”全等搬运线段。

下面两种方法使用不同的 \(G\)，分别作图，不共用辅助点。

\block{方法一 倍长 \(AB\)，把公共直角顶点搬到 \(C\)}

\diagram{\hhOneExtend}

\solve

延长 \(AB\) 到 \(G\)，使 \(BG=AB\)，连接 \(FG\)、\(CG\)。

第一步，用中点搬运 \(AD\)。

\(AB=GB\)，\(BD=BF\)，\(\angle ABD=\angle GBF\)，所以

\[
\triangle ABD\cong\triangle GBF\quad(\mathrm{SAS}).
\]

因此 \(AD=GF\)，且 \(AD\parallel GF\)。这里的平行由对应角相等得到；也可以看成把 \(\triangle ABD\) 绕 \(B\) 旋转 \(180^\circ\)，得到 \(\triangle GBF\)。

又 \(AD=AE\)，\(AD\perp AE\)，故

\[
GF=AE,\qquad GF\perp AE.
\]

第二步，补出以 \(C\) 为直角顶点的等腰直角三角形。

\(AB=BC=BG\)，\(BC\perp AG\)，所以 \(\triangle ABC\)、\(\triangle GBC\) 都是等腰直角三角形。于是

\[
CA=CG,\qquad
\angle ACG=\angle ACB+\angle BCG=45^\circ+45^\circ=90^\circ.
\]

第三步，证明第二组全等。

\(CA\perp CG\)，\(AE\perp GF\)。按图中射线的位置，\(\angle CAE\)、\(\angle CGF\) 是两边分别垂直的锐角，因此相等。结合 \(CA=CG\)、\(AE=GF\)，有

\[
\triangle ACE\cong\triangle GCF\quad(\mathrm{SAS}).
\]

所以 \(CE=CF\)，\(\angle ACE=\angle GCF\)。再按图中射线的次序，有

\[
\angle ECF
=\angle ACE+\angle ACG-\angle GCF
=90^\circ.
\]

故 \(CF=CE\)，\(CF\perp CE\)。

\block{方法二 旋转 \(AC\)，在 \(A\) 处直接补成手拉手}

\diagram{\hhOneRotate}

\solve

过 \(A\) 作 \(AG\perp AC\)，交 \(CB\) 向 \(B\) 外的延长线于 \(G\)，连接 \(DG\)。

第一步，说明 \(G\) 为什么正好是旋转后的 \(C\)。

\(\triangle ABC\) 是等腰直角三角形，\(\angle BAC=45^\circ\)。由 \(AG\perp AC\)，得 \(\angle BAG=45^\circ\)；又 \(\angle ABG=90^\circ\)，所以 \(\triangle ABG\) 也是等腰直角三角形。因此

\[
BG=AB=BC,\qquad AG=AC.
\]

可见“过 \(A\) 作 \(AC\) 的垂线”与“延长 \(CB\)，使 \(BG=BC\)”确定的是同一个 \(G\)。

第二步，建立旋转型全等。

把 \(\triangle AEC\) 绕 \(A\) 顺时针旋转 \(90^\circ\)。因为 \(AE=AD\)、\(AE\perp AD\)，\(E\) 落到 \(D\)；因为 \(AC=AG\)、\(AC\perp AG\)，\(C\) 落到 \(G\)。所以旋转后的三角形是 \(\triangle ADG\)，从而

\[
DG=CE,\qquad DG\perp CE.
\]

若写成全等证明：\(AD=AE\)，\(AG=AC\)，且按图有 \(\angle DAG=\angle EAC\)，故 \(\triangle ADG\cong \triangle AEC\)（\(\mathrm{SAS}\)）。垂直关系则由上述旋转对应得到。

第三步，再用中点把 \(DG\) 搬成 \(CF\)。

\(BD=BF\)，\(BG=BC\)，\(\angle DBG=\angle FBC\)，故

\[
\triangle DBG\cong\triangle FBC\quad(\mathrm{SAS}).
\]

这组全等也是中心 \(B\) 的半转：\(D\) 对应 \(F\)，\(G\) 对应 \(C\)。因此 \(DG=CF\)，\(DG\parallel CF\)。结合 \(DG=CE\)、\(DG\perp CE\)，得到

\[
CF=CE,\qquad CF\perp CE.
\]

如果采用对应角相加的写法：第一组全等给出 \(\angle ECA=\angle DGA\)，第二组全等给出 \(\angle BCF=\angle BGD\)。按图，\(GD\) 在 \(\angle BGA\) 内，所以

\[
\angle ECA+\angle BCF
=\angle DGA+\angle BGD
=\angle BGA=45^\circ.
\]

再加上 \(\angle ACB=45^\circ\)，同样得到 \(\angle ECF=90^\circ\)。

\insight{中点先提供一次半转，等腰直角三角形再提供一次四分之一转；两次搬运把待证的线段连接起来。}

\end{solution}
\clearpage
\begin{solution}{题 2 · 解析}
\think

\(C\) 是两个原三角形的底角顶点，属于“脚拉脚”。从底角顶点到直角顶点的那条腰，延长一倍后，连接新的端点与 \(C\)，就能把 \(C\) 补成直角顶点。本题需要对两个三角形各做一次。

初二可沿“倍长、旋转型全等、中点构造、等腰倒角”完成。初三学过相似后，则可以直接比较两组相差 \(\sqrt{2}\) 倍的边。

\block{方法一 两次倍长，再作平行线}

\diagram{\hhTwoHands}

\solve

第一步，分别补出两个新等腰直角三角形。

延长 \(ED\) 到 \(F\)，使 \(DF=DE\)，连接 \(CF\)。因为 \(DC=DE=DF\)，\(CD\perp EF\)，\(\triangle CDE\)、\(\triangle CDF\) 全等，得

\[
CE=CF,\qquad
\angle ECF=\angle ECD+\angle DCF=45^\circ+45^\circ=90^\circ.
\]

延长 \(BA\) 到 \(G\)，使 \(AG=AB\)，连接 \(CG\)。同理，由 \(CA=AB=AG\)、\(CA\perp BG\)，得到

\[
CB=CG,\qquad \angle BCG=90^\circ.
\]

此时 \(\triangle ECF\)、\(\triangle BCG\) 的直角顶点都是 \(C\)，“脚拉脚”已经补成“手拉手”。

第二步，用旋转型全等得到 \(BE=GF\)。

以 \(C\) 为中心顺时针旋转 \(90^\circ\)，\(B\) 落到 \(G\)，\(E\) 落到 \(F\)。因此 \(\triangle BCE\) 对应 \(\triangle GCF\)，得到

\[
\triangle BCE\cong\triangle GCF\quad(\mathrm{SAS}),
\]

于是

\[
BE=GF,\qquad \angle EBC=\angle CGF.
\]

第三步，把 \(GF\) 的长度搬到 \(E\)，构造等腰三角形。

\diagram{\hhTwoParallel}

将 \(G\) 关于 \(D\) 对称到 \(P\)，连接 \(EP\)。\(D\) 是 \(EF\) 的中点，半转使 \(F\) 对应 \(E\)、\(G\) 对应 \(P\)，因此

\[
EP=GF,\qquad EP\parallel GF.
\]

这也就是过 \(E\) 作 \(GF\) 的平行线，与直线 \(BD\) 相交的构造。按所画的内部位置，可写成 \(\triangle EPD\cong \triangle FGD\)（\(\mathrm{ASA}\)）；用半转说明则能同时处理其他 \(D\) 的位置。

\(A\) 是 \(BG\) 的中点，\(D\) 在 \(BA\) 向 \(A\) 外的延长线上，故 \(P\) 在射线 \(BA\) 上，并且 \(BP=2AD>0\)。由 \(BE=GF=EP\)，\(\triangle BEP\) 为等腰三角形。设

\[
\angle EBP=\angle BPE=\alpha.
\]

第四步，用两种角度表达求 \(\alpha \)。

\(P\) 在射线 \(BA\) 上，\(\angle PBC=45^\circ\)，故

\[
\angle EBC=\alpha+45^\circ.
\]

又 \(EP\parallel GF\)，\(PB\) 与 \(GB\) 同向，所以 \(\angle BGF=180^\circ-\alpha \)。\(\triangle BCG\) 为等腰直角三角形，\(\angle BGC=45^\circ\)；按对应的射线方向，\(GC\) 在 \(\angle BGF\) 内。因此

\[
\angle CGF
=\angle BGF-\angle BGC
=180^\circ-\alpha-45^\circ
=135^\circ-\alpha.
\]

由第二步的对应角相等，得

\[
\alpha+45^\circ=135^\circ-\alpha.
\]

所以 \(\alpha =45^\circ\)，\(\angle EBC=90^\circ\)，即 \(BE\perp BC\)。

\textbf{位置补充}

题目没有规定 \(AD<AB\)，所以 \(D\) 不一定在 \(AG\) 内。上面的半转构造仍然适用：

\begin{itemize}
\item \(D\) 在 \(AG\) 内时，\(BP=BD-DG=2BD-BG=2AD\)。
\item \(D\) 在 \(G\) 及其外侧时，\(BP=BD+DG=2BD-BG=2AD\)。
\end{itemize}

两种情况下 \(P\) 都在射线 \(BA\) 上，倒角时使用的是 \(\angle BGF\)，避免把 \(\angle DGF\) 误当成同一个角。若 \(D=G\)，则 \(P=D\)；此时 \(\triangle EPD\)、\(\triangle FGD\) 退化，不能再称为全等三角形，但半转仍给出 \(EP=GF\)、\(EP\parallel GF\)，其余推理照常成立。

\block{方法二 用 \(\sqrt{2}\) 倍关系直接证明相似}

\solve

\(\triangle ABC\) 是等腰直角三角形，所以 \(BC=\sqrt{2}\cdot AC\)；\(\triangle DCE\) 是等腰直角三角形，所以 \(CE=\sqrt{2}\cdot CD\)。因此

\[
\frac{AC}{BC}=\frac{CD}{CE}=\frac1{\sqrt2}.
\]

又 \(\angle BCA=\angle DCE=45^\circ\)。按原图的方向，\(CA\)、\(CE\) 都位于 \(\angle BCD\) 内，并分别从它的两侧截去 \(45^\circ\)，所以

\[
\angle ACD=\angle BCD-45^\circ
=\angle BCE.
\]

由两边成比例且夹角相等，得

\[
\triangle ADC\sim\triangle BEC.
\]

对应顶点为 \(A\leftrightarrow B\)、\(D\leftrightarrow E\)、\(C\leftrightarrow C\)，因此 \(\angle EBC=\angle DAC\)。\(D\) 在 \(BA\) 的延长线上，而 \(AC\perp AB\)，所以 \(\angle DAC=90^\circ\)。故 \(BE\perp BC\)。

\insight{初二用倍长把两个底角补成直角，初三用共同的 \(\sqrt{2}\) 倍关系直接建立相似。}

\end{solution}

\section{模型总结与迁移}
\ifteachernotes
\block{一 先判断公共点的身份}

这里把等腰三角形两腰所夹的顶角叫作“手”，底边两端的底角顶点叫作“脚”。在本章的等腰直角三角形中，“手”就是直角顶点。

{\small\renewcommand{\arraystretch}{1.25}
\begin{tabularx}{\linewidth}{@{}p{1.2cm} X X p{2cm}@{}}
\toprule
构型 & 公共点的身份 & 优先考虑的构造 & 本章对应 \\
\midrule
手拉手 & 两个三角形的直角顶点 & 绕公共顶点同向旋转 \(90^\circ\)，配对端点 & 两道题构造后的新图形 \\
手拉脚 & 一个直角顶点、一个底角顶点 & 倍长经过直角顶点的腰，或旋转另一条边 & 题 1 \\
脚拉脚 & 两个底角顶点 & 各倍长一条腰，分别把公共点补成直角顶点 & 题 2 \\
\bottomrule
\end{tabularx}}

名称只是帮助认图。判断时要写出实际的等长条件、直角位置和旋转对应，不能仅凭图形像不像。

\block{二 手拉手为什么能同时得到等长与垂直}

设 \(\triangle OAB\)、\(\triangle OCD\) 是共用直角顶点 \(O\) 的等腰直角三角形，\(OA=OB\)，\(OC=OD\)。若同一个方向的 \(90^\circ\)旋转使 \(A\) 对应 \(B\)，同时使 \(C\) 对应 \(D\)，那么整条 \(AC\) 就会对应 \(BD\)。

当 \(A\)、\(C\) 不重合时，旋转保持长度，又把线段方向转过 \(90^\circ\)，所以

\[
AC=BD,\qquad AC\perp BD.
\]

用全等来写，就是 \(\triangle OAC\cong \triangle OBD\)（\(\mathrm{SAS}\)）；用旋转来理解，还能立即说明垂直性。这也解释了为什么写证明时必须明确端点的对应。

题 1 方法二中，以 \(A\) 为中心，\(E\to D\)、\(C\to G\)，得到 \(CE=DG\)、\(CE\perp DG\)；题 2 中，以 \(C\) 为中心，\(B\to G\)、\(E\to F\)，得到 \(BE=GF\)，并得到对应角相等。随后还要利用中点，把这个新结论送到题目要求的线段上。

\block{三 倍长为什么能把脚补成手}

设 \(\triangle OXY\) 在 \(X\) 处为直角，\(XO=XY\)。公共点 \(O\) 原来是一个底角顶点。延长 \(YX\) 到 \(Z\)，使 \(XZ=XY\)，连接 \(OZ\)。

由于 \(XO=XY=XZ\)、\(XO\perp YZ\)，两侧小三角形都是等腰直角三角形，故

\[
OY=OZ,\qquad
\angle YOZ=45^\circ+45^\circ=90^\circ.
\]

新三角形 \(\triangle YOZ\) 的直角顶点便是 \(O\)。操作时应倍长经过原直角顶点 \(X\) 的腰，连接原底角顶点 \(O\) 与新的端点 \(Z\)。

题 1 方法一用 \(\triangle ABC\)、\(\triangle GBC\) 在 \(C\) 处拼出直角；方法二用 \(\triangle ABC\)、\(\triangle ABG\) 在 \(A\) 处拼出直角。题 2 则分别在 \(A\)、\(D\) 两个位置倍长，最后都在 \(C\) 处拼出直角。

\block{四 中点怎样帮助完成第二次搬运}

\textbf{两条线段共用中点。} 若 \(B\) 同时是 \(DF\)、\(AG\) 的中点，绕 \(B\) 半转便有 \(D\to F\)、\(A\to G\)，得到对应线段相等且平行。题 1 方法一采用的就是这个构造。

\textbf{中点配平行线。} 若 \(D\) 是 \(EF\) 的中点，\(G\) 关于 \(D\) 的对称点是 \(P\)，则半转有 \(F\to E\)、\(G\to P\)，得到 \(GF=EP\)、\(GF\parallel EP\)。题 2 把旋转型全等给出的 \(BE=GF\) 接到这里，才构造出 \(BE=EP\) 的等腰三角形。

遇到中点，不必机械地延长待证线段。先确定需要搬运哪条边，再决定应当倍长或作平行线的位置。

\block{五 与夹半角的联系}

题 1 的 \(A\) 处既有 \(\angle DAE=90^\circ\)，又有 \(\angle BAC=45^\circ\)。这提示可以把相关三角形绕 \(A\) 旋转 \(90^\circ\)，把原本分开的边放在一起。旋转之后，等腰直角三角形保证端点对应，中点条件完成另一组全等。

两类模型都要看等长和角度，但本章的结论来自具体的旋转对应；每一次换构型，都应重新检查对应端点和射线方向。

\block{易错提醒}

\begin{itemize}
\item 共用一个点，不等于共用直角顶点。先标明公共点在两个三角形中分别是什么身份。
\item 两个等腰直角三角形必须按同一旋转方向配对端点。把其中一个三角形翻到另一侧后，原来的连线结论可能改变。
\item 全等判定要根据已有条件填写。本章关键步骤使用两边及夹角，判定是 \(\mathrm{SAS}\)；平行线构造中的两角及夹边可以用 \(\mathrm{ASA}\)。
\item 中点带来的等长、平行必须由全等或半转说明，不能仅因图中出现“八字形”就直接使用。
\item 等长与垂直要分别交代依据。只证明了 \(CE=CF\)，还不能据此断定 \(CE\perp CF\)。
\item 题 2 的“如图”包含两个三角形的方位。本章将它明确为 \(E\)、\(C\) 在 \(AB\) 两侧；若把 \(E\) 换到另一侧，\(BE\perp BC\) 不再是同一个命题。
\item 点序改变后，重新判断角的加减；辅助三角形退化时，改用仍然有效的旋转或对称关系。
\end{itemize}
\else
请从公共点的身份、同向旋转、倍长与中点搬运四个方面，自主整理本章方法。
\fi


\clearpage
\phantomsection
\part*{第三章\quad 双等腰模型（婆罗摩笈多模型）}
\addcontentsline{toc}{part}{第三章 双等腰模型（婆罗摩笈多模型）}
\renewcommand{\theHsection}{03-double-isosceles.\arabic{section}}
\setcounter{section}{-1}
\section{一线三等角：快速回顾}
先看一幅熟悉的图：\(A\)、\(B\)、\(C\) 在同一直线上，\(B\) 在 \(A\)、\(C\) 之间；\(D\)、\(E\) 在直线 \(AC\) 的同侧，且
\[
\angle DAB=\angle DBE=\angle BCE=90^\circ,\qquad BD=BE.
\]
\diagram{\diReview}

直线上的三个角满足
\[
\angle DBA+\angle DBE+\angle EBC=180^\circ,
\]
所以 \(\angle DBA+\angle EBC=90^\circ\)。而在直角三角形 \(DAB\) 中，
\(\angle BDA+\angle DBA=90^\circ\)，故 \(\angle BDA=\angle EBC\)。

再结合 \(\angle DAB=\angle BCE\) 和 \(BD=BE\)，得到
\[
\triangle DBA\cong\triangle BEC\quad(\mathrm{AAS}),
\]
于是 \(\boldsymbol{AB=CE}\)、\(\boldsymbol{AD=BC}\)。注意对应顺序：
\(D\leftrightarrow B\)、\(B\leftrightarrow E\)、\(A\leftrightarrow C\)。

\begin{strand}
\textbf{识别顺序}\quad 一条直线 \(\longrightarrow\) 三个直角
\(\longrightarrow\) 等斜边 \(\longrightarrow\) 全等后交换对应直角边。

本章只回顾这一种直角全等型，下面“作双垂”的证明会连续用到它两次。
\end{strand}
\clearpage

\section{双等腰模型的三项结论}
\block{共用构型}
如图，\(\triangle CBC'\) 和 \(\triangle DBD'\) 都是以 \(B\) 为直角顶点的等腰直角三角形，即
\[
BC=BC',\quad BD=BD',\quad
\angle CBC'=\angle DBD'=90^\circ.
\]
各点方向按图取定：射线 \(BD\)、\(BD'\)、\(BC\)、\(BC'\) 绕 \(B\) 依次排列，\(\angle CBD'\) 为锐角。从 \(BC\) 到 \(BC'\)、从 \(BD\) 到 \(BD'\) 的 \(90^\circ\) 转动方向相同。以下三题共用这一构型，分别作辅助线。
\diagram{\diModel}

\begin{problem}{题 1 等面积从哪里来}
求证：\(S_{\triangle CBD'}=S_{\triangle C'BD}\)。
\hint{已经有 \(BC=BC'\)。分别以 \(BC\)、\(BC'\) 为底，试着比较两条高。}
\end{problem}

\begin{problem}{题 2 已知中点，证明垂直}
\(E\) 是 \(CD'\) 的中点。连接 \(BE\)，求证：\(BE\perp C'D\)。
\hint{中点条件能否通过倍长 \(BE\)，转化成一组八字形全等？}
\end{problem}

\begin{problem}{题 3 已知垂直，证明中点}
过 \(B\) 作 \(BF\perp C'D\)，\(F\) 为垂足；直线 \(BF\) 与 \(CD'\) 相交于 \(E\)。求证：\(CE=ED'\)。
\hint{还不知道 \(E\) 是中点。分别从 \(C\)、\(D'\) 向直线 \(BF\) 作垂线，看看会出现几组一线三等角。}
\end{problem}
\clearpage

\section{集中解析}
\begin{solution}{题 1 · 解析}
\think 两个三角形不一定全等，等面积也不必比较三边。先用已知的 \(BC=BC'\) 作等底，再把问题转成等高。

\solve 过 \(D'\) 作 \(D'M\perp BC\)，\(M\) 为垂足；过 \(D\) 作 \(DN\perp BC'\)，\(N\) 为垂足。本图中 \(M\) 在射线 \(BC\) 上，\(N\) 在 \(C'B\) 向 \(B\) 外的延长线上。
\diagram{\diArea}

由两个公共顶角都是直角，按图有
\[
\angle CBD'=180^\circ-\angle C'BD.
\]
由于 \(BM\) 与 \(BC\) 同向、\(BN\) 与 \(BC'\) 反向，
\[
\angle MBD'=\angle CBD',\qquad
\angle NBD=180^\circ-\angle C'BD,
\]
所以 \(\angle MBD'=\angle NBD\)。结合 \(\angle BMD'=\angle BND=90^\circ\)、\(BD'=BD\)，得到
\[
\triangle BMD'\cong\triangle BND\quad(\mathrm{AAS}),\qquad D'M=DN.
\]
从而
\[
S_{\triangle CBD'}=\frac12 BC\cdot D'M
=\frac12 BC'\cdot DN=S_{\triangle C'BD}.
\]

\insight{等面积先找等底；高不现成，就作垂线，把等高交给直角三角形全等。}
\end{solution}
\clearpage

\begin{solution}{题 2 · 解析}
\think \(E\) 是 \(CD'\) 的中点，但这个等长条件还不能和 \(BC=BC'\)、\(BD=BD'\) 放进同一组全等。倍长 \(BE\)，把 \(BD'\) 搬到 \(C\) 处，再比较一个大三角形。

\solve 延长 \(BE\) 到 \(G\)，使 \(EG=EB\)，连接 \(CG\)。于是 \(B\)、\(E\)、\(G\) 共线，\(CE=ED'\)，\(\angle CEG=\angle D'EB\)。
\diagram{\diMidpoint}

\solvehead{第一步：用八字形全等搬运一条边}
\[
\triangle CEG\cong\triangle D'EB\quad(\mathrm{SAS}).
\]
因此 \(CG=D'B=DB\)，且 \(\angle ECG=\angle ED'B\)，故 \(CG\parallel BD'\)。按图，射线 \(CG\) 与 \(BD'\) 同向。

\solvehead{第二步：比较一组大三角形}
由于 \(CB\) 与 \(BC\) 反向，\(CG\) 与 \(BD'\) 同向，
\[
\angle GCB=180^\circ-\angle CBD'=\angle DBC'.
\]
结合 \(CG=DB\)、\(CB=BC'\)，得到
\[
\triangle GCB\cong\triangle DBC'\quad(\mathrm{SAS}).
\]
于是 \(\angle GBC=\angle DC'B\)。

\solvehead{第三步：用对应角推出垂直}
延长 \(EB\) 过 \(B\)，与 \(C'D\) 相交于 \(H\)。按图，\(BH\)、\(BG\) 是反向射线，\(\angle CBC'=90^\circ\)，所以
\[
\angle C'BH+\angle GBC=90^\circ.
\]
又因 \(H\) 在 \(C'D\) 上，\(\angle BC'H=\angle DC'B=\angle GBC\)。在 \(\triangle BC'H\) 中，\(B\)、\(C'\) 两个角之和为 \(90^\circ\)，故 \(\angle BHC'=90^\circ\)，即 \(BE\perp C'D\)。

\solvehead{顺手得到的长度关系}
第二次全等还给出 \(GB=DC'\)。因为 \(GB=2BE\)，故
\[
\boxed{2BE=C'D}.
\]

\insight{倍长中线先搬运等长边，再用大三角形全等同时取得角度与长度关系。}
\end{solution}
\clearpage

\begin{solution}{题 3 · 解析}
\think 这一题和题 2 的已知、结论交换了位置。现在 \(CE=ED'\) 还是待证，不能把它放进倍长后的全等条件。题目给出了垂直，便从 \(C\)、\(D'\) 向 \(BF\) 作双垂，连续用两次章首模型。

\solve 过 \(C\) 作 \(CM\perp BF\)，过 \(D'\) 作 \(D'N\perp BF\)，\(M\)、\(N\) 为垂足。本题的 \(M\)、\(N\) 与题 1 的垂足不同。
\diagram{\diConverse}

\solvehead{第一步：用 \(BC=BC'\)，得到 \(CM=BF\)}
\(M\)、\(B\)、\(F\) 共线，\(\angle CMB=\angle CBC'=\angle BFC'=90^\circ\)。由直线上的角和，
\[
\angle CBM+\angle C'BF=90^\circ.
\]
而在直角三角形 \(C'BF\) 中，\(\angle BC'F+\angle C'BF=90^\circ\)，故 \(\angle CBM=\angle BC'F\)。再结合 \(BC=BC'\)，得到
\[
\triangle BCM\cong\triangle C'BF\quad(\mathrm{AAS}),\qquad CM=BF.
\]
这正是章首的一线三等角：等斜边推出对应直角边相等。

\solvehead{第二步：用 \(BD'=BD\)，得到 \(D'N=BF\)}
\(N\)、\(B\)、\(F\) 共线，\(\angle D'NB=\angle D'BD=\angle BFD=90^\circ\)。同样，
\begin{align*}
\angle D'BN+\angle DBF&=90^\circ,\\
\angle BDF+\angle DBF&=90^\circ,
\end{align*}
所以 \(\angle D'BN=\angle BDF\)。结合 \(BD'=BD\)，得到
\[
\triangle BD'N\cong\triangle DBF\quad(\mathrm{AAS}),\qquad D'N=BF.
\]
于是 \(CM=D'N\)。

\solvehead{第三步：接上八字形全等，得到中点}
按图，\(C\)、\(E\)、\(D'\) 共线，\(M\)、\(E\)、\(N\) 共线，\(\angle CME=\angle D'NE=90^\circ\)，\(\angle CEM=\angle D'EN\)。结合 \(CM=D'N\)，得到
\[
\triangle CME\cong\triangle D'NE\quad(\mathrm{AAS}),\qquad CE=ED'.
\]
所以 \(E\) 是 \(CD'\) 的中点。

若 \(M\)、\(N\) 恰好都与 \(E\) 重合，则 \(CM=D'N\) 直接给出 \(CE=ED'\)，不再使用退化三角形的全等。

\insight{要从垂直证明中点，先用双垂把同一条 \(BF\) 搬到两侧，再用八字形全等收尾。}
\end{solution}
\clearpage

\section{模型总结与迁移}
\ifteachernotes
\block{三个条件，三条路线}
\begin{center}\small
\renewcommand{\arraystretch}{1.3}
\begin{tabularx}{\linewidth}{@{}p{.26\linewidth}p{.22\linewidth}X@{}}
\toprule
已知／目标 & 辅助线 & 真正起作用的关系\\
\midrule
两等腰直角，证明等面积 & 以等长边作底，分别作高 & 等底＋直角三角形全等得到等高\\
\addlinespace
已知中点，证明垂直 & 倍长 \(BE\) & 八字形全等搬运边，再比较大三角形\\
\addlinespace
已知垂直，证明中点 & 从 \(C\)、\(D'\) 向 \(BE\) 作双垂 & 两次一线三等角，再接八字形全等\\
\bottomrule
\end{tabularx}
\end{center}

在本章构型下，可以把题 2、题 3 合起来记：
\[
E\in CD':\qquad CE=ED'\ \Longleftrightarrow\ BE\perp C'D.
\]
题 2 还告诉我们：当 \(E\) 是 \(CD'\) 的中点时，\textbf{\(BE\) 既垂直于 \(C'D\)，又等于 \(C'D\) 的一半}。

\block{再看一次等面积}
题 3 的两组全等给出 \(CM=D'N=BF\)。以 \(BE\) 为底，\(\triangle CBE\) 和 \(\triangle D'BE\) 的高都是 \(BF\)，所以
\begin{align*}
S_{\triangle CBD'}
&=S_{\triangle CBE}+S_{\triangle D'BE}\\
&=BE\cdot BF.
\end{align*}
而题 2 已经得到 \(2BE=C'D\)，故
\[
S_{\triangle CBD'}=\frac12 C'D\cdot BF=S_{\triangle C'BD}.
\]
这给出等面积的另一种看法：倍长中线得到的长度关系，和双垂得到的高度关系，最终在面积公式里相遇。

\block{易错提醒}
\begin{enumerate}
\item 两个等腰直角三角形的方向要按图取定。只记两组等长和两个直角，不看方位，不能直接套用本章结论。
\item 高是到\textbf{直线}的距离，垂足可能在边的延长线上；题 1 的 \(N\) 就在 \(BC'\) 的延长线上。
\item 题 3 中的中点是结论，不能预先使用 \(CE=ED'\) 来证明全等。
\item “两块面积始终相等”不表示“每块面积是一个固定值”；构型变化时，两块面积可以一起变化。
\item 三题分别作图。题 1、题 3 的 \(M\)、\(N\) 含义不同；题 2 的 \(G\) 是倍长点，\(H\) 是交点。
\end{enumerate}
\else
\block{学生自主总结}
在每道题旁记录：已知条件、第一条辅助线、第一次全等搬运了什么、第二次全等得到什么。再尝试写出中点、垂直与长度之间的关系。
\vspace{8\baselineskip}
\fi


\clearpage
\phantomsection
\part*{第四章\quad 八字形模型}
\addcontentsline{toc}{part}{第四章 八字形模型}
\renewcommand{\theHsection}{04-eight-shape.\arabic{section}}
\begingroup
\setcounter{secnumdepth}{0}
\setlength{\abovedisplayskip}{7pt plus 2pt minus 1pt}
\setlength{\belowdisplayskip}{7pt plus 2pt minus 1pt}
\setlength{\abovedisplayshortskip}{4pt plus 1pt}
\setlength{\belowdisplayshortskip}{4pt plus 1pt}
% Synced from content.md by sync_latex.py.
\section{\texorpdfstring{关键的八字形}{关键的八字形}}

$\triangle ABC$ 与 $\triangle DBC$ 有公共边 $BC$，连接 $AD$，$AC$ 与 $BD$ 相交于 $O$。

\esdiagram{\esdiModel}

由八字形得

\[
\angle BAC+\angle ABD=\angle BDC+\angle ACD.
\]

因此，$\angle BAC=\angle BDC$ 与 $\angle ABD=\angle ACD$ 可以相互推出。图中两个顶角为直角时，就能得到 $\angle ABD=\angle ACD$。

% Synced from content.md by sync_latex.py.
\section{\texorpdfstring{例 $1$ 四个条件，知三推一}{例 1 四个条件，知三推一}}

按图取定方位：$A$、$D$ 在 $BC$ 的同侧，$AC$ 与 $BD$ 在内部相交。延长 $CD$ 至 $X$，连接 $AD$，$DA$ 位于 $\angle BDX$ 内部。有四个条件：

\esdiagram{\esdiExterior}

\begin{center}
\small
\begin{tabularx}{\linewidth}{lL}
\toprule
条件 & 内容 \\
\midrule
\escond{1} & $AB=AC$ \\
\escond{2} & $\angle BAC=90^\circ$ \\
\escond{3} & $\angle BDC=90^\circ$ \\
\escond{4} & $DA$ 平分 $\angle BDX$，即 $\angle BDA=\angle ADX$ \\
\bottomrule
\end{tabularx}
\end{center}

\textbf{知道其中任意三个，就能推出剩下的一个。} 

\needspace{6\baselineskip}\subsection{\texorpdfstring{（$1$）$\text{\escond{1}\escond{2}\escond{3}}\Rightarrow\text{\escond{4}}$：连接直角顶点，得到外角平分线}{（1）①②③⇒④：连接直角顶点，得到外角平分线}}

已知 $AB=AC$，$\angle BAC=\angle BDC=90^\circ$，说明 $DA$ 平分 $\angle BDX$。

\think

两个直角先给出 $\angle ABD=\angle ACD$。接下来有两个方向：从 $A$ 向 $BD$、$CD$ 作双垂，把斜边 $AB$、$AC$ 配成直角三角形；或者围绕 $A$ 补出另一个等腰直角三角形，把上一章的手拉手用起来。

由 $\angle BAC=\angle BDC$，利用八字形倒角，得

\[
\angle ABD=\angle ACD.
\]

\needspace{6\baselineskip}\textbf{方法一：从 $A$ 作双垂。}

过 $A$ 作 $AM\perp BD$，垂足为 $M$；作 $AN\perp CD$，垂足为 $N$。本图中 $M$ 在线段 $BD$ 上，$N$ 在 $CD$ 经过 $D$ 的延长线上。

\esdiagram{\esdiDoublePerpendicular}

在 $\mathrm{Rt}\triangle ABM$ 与 $\mathrm{Rt}\triangle ACN$ 中，$AB=AC$，$\angle ABM=\angle ACN$，且两垂足处都是直角，故

\[
\triangle ABM\cong\triangle ACN\quad(\mathrm{AAS}),
\qquad AM=AN.
\]

再看 $\mathrm{Rt}\triangle ADM$ 与 $\mathrm{Rt}\triangle ADN$：斜边 $AD$ 公共，$AM=AN$，所以

\[
\triangle ADM\cong\triangle ADN\quad(\mathrm{HL}).
\]

因此 $\angle MDA=\angle NDA$，即 $\angle BDA=\angle ADX$。又因为 $\angle BDA+\angle ADX=90^\circ$，所以

\[
\boxed{\angle BDA=\angle ADX=45^\circ.}
\]

\needspace{6\baselineskip}\textbf{方法二：过 $A$ 添直角，构造手拉手。}

过 $A$ 作 $AP\perp AD$，交 $BD$ 于 $P$。

\esdiagram{\esdiHandLower}

由 $AB\perp AC$、$AP\perp AD$，按图有 $\angle BAP=\angle CAD$；又有 $\angle ABP=\angle ACD$、$AB=AC$，故

\[
\triangle ABP\cong\triangle ACD\quad(\mathrm{ASA}),
\qquad AP=AD.
\]

因此 $\triangle APD$ 为等腰直角三角形，$\angle PDA=45^\circ$。$P$ 在 $DB$ 上，故 $\angle BDA=45^\circ$，$\angle ADX=90^\circ-45^\circ=45^\circ$。

也可以先在 $BD$ 上取 $P$，使 $BP=CD$。由 $\angle ABP=\angle ACD$、$AB=AC$，先以 $\mathrm{SAS}$ 证明同一组全等，得到 $AP=AD$、$\angle BAP=\angle CAD$，再推出 $AP\perp AD$。两种作法落到同一幅图上：一种先添角，另一种先添边。

\insight{共斜边先倒角，再把“等腰”转成一组全等；外角平分线由等距离或等腰直角三角形得到。}

\needspace{6\baselineskip}\subsection{\texorpdfstring{（$2$）$\text{\escond{1}\escond{2}\escond{4}}\Rightarrow\text{\escond{3}}$：已知外角平分线，反推 $D$ 处直角}{（2）①②④⇒③：已知外角平分线，反推 D 处直角}}

已知 $AB=AC$，$\angle BAC=90^\circ$，$DA$ 平分 $\angle BDX$，说明 $\angle BDC=90^\circ$。

\think

这次不能先用两个直角倒角，因为 $D$ 处的直角正是待证结论。$\angle BDA=\angle ADX$ 能先制造对称：作双垂得到 $AM=AN$，或截取等长边制造 $\mathrm{SAS}$ 全等。把对称带来的等长再接上 $AB=AC$，便能推出 $\angle ABD=\angle ACD$。

\needspace{6\baselineskip}\textbf{方法一：作双垂，先用角平分线。}

沿用例 $1$（$1$）的 $AM$、$AN$。由 $\angle BDA=\angle ADX$、公共边 $AD$ 和两个直角，得

\[
\triangle ADM\cong\triangle ADN\quad(\mathrm{AAS}),
\qquad AM=AN.
\]

再结合 $AB=AC$，以 $\mathrm{HL}$ 得 $\mathrm{Rt}\triangle ABM\cong \mathrm{Rt}\triangle ACN$，所以 $\angle ABD=\angle ACD$。由八字形倒角，

\[
\angle BDC=\angle BAC=90^\circ.
\]

\needspace{6\baselineskip}\textbf{方法二：向上延长 $CD$，截取 $DE=DB$。}

延长 $CD$ 至 $E$，使 $DE=DB$，连接 $AE$。

\esdiagram{\esdiExtendCd}

由 $\angle BDA=\angle ADE$、公共边 $AD$、$DB=DE$，得

\[
\triangle ADB\cong\triangle ADE\quad(\mathrm{SAS}).
\]

所以 $AB=AE$，$\angle ABD=\angle AED$。又因为 $AB=AC$，所以 $AE=AC$，即 $\angle AEC=\angle ACE$。因此$\angle AEC=\angle ACE=\angle ABD$

最后八字形得 $\angle BDC=90^\circ$。

\needspace{6\baselineskip}\textbf{方法三：向右延长 $BD$，截取 $DF=DC$。}

延长 $BD$ 至 $F$，使 $DF=DC$，连接 $AF$。

\esdiagram{\esdiExtendBd}



\[
\angle ADF=180^\circ-\angle BDA,
\qquad \angle ADC=180^\circ-\angle ADX.
\]

等角的补角相等，所以 $\angle ADF=\angle ADC$。结合 $AD$ 公共、$DF=DC$，得

\[
\triangle ADF\cong\triangle ADC\quad(\mathrm{SAS}),
\qquad AF=AC,
\quad \angle AFD=\angle ACD.
\]

故 $AB=AC=AF$，从而 $\angle ABD=\angle AFD=\angle ACD$。最后八字形得 $\angle BDC=90^\circ$。

\insight{已知角平分线，可以先围绕它制造全等或者等腰。}

\needspace{6\baselineskip}\subsection{\texorpdfstring{（$3$）$\text{\escond{2}\escond{3}\escond{4}}\Rightarrow\text{\escond{1}}$：保留两个直角，反推等腰}{（3）②③④⇒①：保留两个直角，反推等腰}}

已知 $\angle BAC=\angle BDC=90^\circ$，$DA$ 平分 $\angle BDX$，说明 $AB=AC$。

\think

两个直角给 $\angle ABD=\angle ACD$；外角平分线给 $AM=AN$。这一次 $AB=AC$ 是待证结论，全等不能再用 $\mathrm{HL}$。已有两个角和一条直角边，可以改用 $\mathrm{AAS}$。

由 $\angle BAC=\angle BDC=90^\circ$，八字形倒角得 $\angle ABD=\angle ACD$。

从 $A$ 向 $BD$、$CD$ 作双垂，记垂足为 $M$、$N$。由 $\angle BDA=\angle ADX$，得 $\mathrm{Rt}\triangle ADM\cong \mathrm{Rt}\triangle ADN$，从而 $AM=AN$。

再由 $\angle ABM=\angle ACN$、两个直角、$AM=AN$，得

\[
\triangle ABM\cong\triangle ACN\quad(\mathrm{AAS}),
\qquad AB=AC.
\]

\textbf{另外两种对称全等的写法。}

\begin{itemize}
\item 按（$2$）方法二作 $E$，得到 $AB=AE$、$\angle ABD=\angle AED$。由 $\angle ABD=\angle ACD$，得 $\angle AEC=\angle ACE$，所以 $AE=AC$，进而 $AB=AC$。
\item 按（$2$）方法三作 $F$，得到 $AF=AC$、$\angle AFD=\angle ACD$。由 $\angle ABD=\angle ACD$，得 $\angle AFB=\angle ABF$，所以 $AB=AF$，进而 $AB=AC$。
\end{itemize}

\insight{条件与结论交换后，要重新选择全等判定；等边待证时，用等角和等距离搭桥。}

\needspace{6\baselineskip}\subsection{\texorpdfstring{（$4$）$\text{\escond{1}\escond{3}\escond{4}}\Rightarrow\text{\escond{2}}$：保留等边，反推 $A$ 处直角}{（4）①③④⇒②：保留等边，反推 A 处直角}}

已知 $AB=AC$，$\angle BDC=90^\circ$，$DA$ 平分 $\angle BDX$，说明 $\angle BAC=90^\circ$。

\think

$ABC$ 是否为直角三角形还不知道，不能提前使用手拉手或 $K$ 型全等。现在真正可用的是外角平分线和 $AB=AC$：先获得等距离，再用 $\mathrm{HL}$ 换出 $\angle ABD=\angle ACD$，最后倒角。

从 $A$ 向 $BD$、$CD$ 作双垂，垂足分别为 $M$、$N$。由 $\angle BDA=\angle ADX$、$AD$ 公共和两个直角，得 $\mathrm{Rt}\triangle ADM\cong \mathrm{Rt}\triangle ADN$，所以 $AM=AN$。

又 $AB=AC$，故

\[
\triangle ABM\cong\triangle ACN\quad(\mathrm{HL}),
\qquad \angle ABD=\angle ACD.
\]

由八字形倒角，得

\[
\angle BAC=\angle BDC=90^\circ.
\]

也可以照例 $1$（$2$）的两种延长作法：先由对称全等和 $AB=AC$ 得到 $\angle ABD=\angle ACD$，再把已知的 $\angle BDC=90^\circ$ 换回 $\angle BAC=90^\circ$。这两种作法在得出 $\angle ABD=\angle ACD$ 之前，都没有用到 $A$ 处是直角。

\insight{先清点已知条件；通过等距离与等斜边得到等角，再把直角搬回另一顶点。}

\needspace{6\baselineskip}\subsection{\texorpdfstring{单角变式：把\escond{4}换成 $45^\circ$ 或 $135^\circ$}{单角变式：把④换成 45° 或 135°}}

\needspace{8\baselineskip}\subsubsection*{\texorpdfstring{变式一：只给 $\angle BDA=45^\circ$}{变式一：只给 ∠BDA＝45°}}

已知 $AB=AC$，$\angle BAC=90^\circ$，$\angle BDA=45^\circ$，说明 $\angle BDC=90^\circ$。

\think

只给 $\angle BDA=45^\circ$，还没有给 $\angle ADX$。此时“到角两边距离相等”不能直接用。应让这个 $45^\circ$ 和新作的直角出现在同一个三角形里，先得到等腰直角三角形，再用手拉手；另一条路是把 $B$、$C$ 向 $AD$ 作双垂，用 $K$ 型全等搬运两条直角边。

\needspace{6\baselineskip}\textbf{方法一：过 $A$ 作 $AP\perp AD$，交 $BD$ 于 $P$。}

\esdiagram{\esdiHandLower}

由于 $\angle ADP=45^\circ$、$\angle PAD=90^\circ$，$\triangle APD$ 是等腰直角三角形，故 $AP=AD$。

由 $AB=AC$、$AP=AD$、$\angle BAP=\angle CAD$，得

\[
\triangle ABP\cong\triangle ACD\quad(\mathrm{SAS}).
\]

所以 $\angle ABP=\angle ACD$，即 $\angle ABD=\angle ACD$。由八字形倒角，$\angle BDC=\angle BAC=90^\circ$。

与例 $1$（$1$）对照：那里先全等，再证明 $AP=AD$；这里先由 $45^\circ$ 得 $AP=AD$，再全等。辅助线相同，证明顺序由已知条件决定。

\needspace{6\baselineskip}\textbf{方法二：从 $B$、$C$ 向 $AD$ 作双垂。}

过 $B$ 作 $BM\perp AD$，过 $C$ 作 $CN\perp AD$，垂足分别为 $M$、$N$。本构型中，沿 $AD$ 的方向，点的顺序为 $M$、$A$、$D$、$N$。

\esdiagram{\esdiKPerpendicular}

在 $\mathrm{Rt}\triangle ABM$ 与 $\mathrm{Rt}\triangle CAN$ 中，$\angle BAM=\angle ACN$，$AB=CA$，所以

\[
\triangle ABM\cong\triangle CAN\quad(\mathrm{AAS}),
\qquad BM=AN,\quad AM=CN.
\]

这里的角相等来自 $AB\perp AC$、$AM\perp CN$，按图取相等的锐角。

由于 $\angle BDM=45^\circ$，$\mathrm{Rt}\triangle BMD$ 为等腰直角三角形，所以 $BM=DM$。于是

\[
AN=DM=AM+AD,
\qquad DN=AN-AD=AM=CN.
\]

故 $\mathrm{Rt}\triangle CDN$ 也是等腰直角三角形，$\angle CDN=45^\circ$。$DN$ 与 $DA$ 反向，所以 $\angle ADX=45^\circ$。又 $\angle BDA=45^\circ$，故 $\angle BDC=180^\circ-45^\circ-45^\circ=90^\circ$。

\insight{只给单个 $45^\circ$ 时，先补等腰直角三角形；用 $K$ 型全等时，等长常藏在“加减公共部分”中。}

\needspace{8\baselineskip}\subsubsection*{\texorpdfstring{变式二：换成 $\angle ADC=135^\circ$}{变式二：换成 ∠ADC＝135°}}

已知 $AB=AC$，$\angle BAC=90^\circ$，$\angle ADC=135^\circ$，说明 $\angle BDC=90^\circ$。

\think

$\angle ADC=135^\circ$ 等价于 $\angle ADX=45^\circ$。与变式一对照，现在方便补出的是 $AD$ 与 $CD$ 延长线之间的等腰直角三角形。$K$ 型全等也仍可用，只是先得到等腰直角的变成了 $\triangle CDN$。

\needspace{6\baselineskip}\textbf{方法一：向上补等腰直角三角形。}

过 $A$ 作 $AQ\perp AD$，交 $CD$ 经过 $D$ 的延长线于 $Q$。

\esdiagram{\esdiHandUpper}

由 $\angle ADQ=180^\circ-135^\circ=45^\circ$，得 $AQ=AD$。又 $AB=AC$，且由两对垂直关系，按图有 $\angle BAD=\angle CAQ$，故

\[
\triangle ABD\cong\triangle ACQ\quad(\mathrm{SAS}).
\]

所以 $\angle ABD=\angle ACQ=\angle ACD$。再倒角得 $\angle BDC=90^\circ$。

\needspace{6\baselineskip}\textbf{方法二：从 $B$、$C$ 向 $AD$ 作双垂。}

沿用变式一的 $K$ 型图，由 $\mathrm{Rt}\triangle ABM\cong \mathrm{Rt}\triangle CAN$，有 $BM=AN$、$AM=CN$。

因为 $\angle CDN=180^\circ-\angle CDA=45^\circ$，所以 $\mathrm{Rt}\triangle CDN$ 为等腰直角三角形，$CN=DN$。于是

\[
AM=DN,
\qquad DM=AM+AD=DN+AD=AN=BM.
\]

故 $\mathrm{Rt}\triangle BMD$ 为等腰直角三角形，$\angle BDA=45^\circ$。已知 $\angle ADX=45^\circ$，所以 $\angle BDC=90^\circ$。

\insight{$135^\circ$ 先看邻补角；与 $45^\circ$ 变式比较，交换的是先出现的等腰直角三角形。}

\newpage
% Synced from content.md by sync_latex.py.
\section{\texorpdfstring{例 $2$ 角平分线上的二倍关系}{例 2 角平分线上的二倍关系}}

如图，在 $\triangle ABC$ 中，$\angle A=90^\circ$，$AB=AC$。$BD$ 平分 $\angle ABC$，交 $AC$ 于 $D$。过 $C$ 作 $CE\perp BD$，交 $BD$ 的延长线于 $E$。求证：$BD=2$$CE$。

\esdiagram{\esdiExample}

\think

$\angle ABC=\angle ACB=45^\circ$，而 $BD$ 平分 $\angle ABC$，所以 $\angle ABD=\angle DBC=22.5^\circ$。目标 $BD=2$$CE$ 可以分别拆成四条路线：

\begin{enumerate}
\item 延长相交：把 $BD$ 换成一条整段 $CF$，再用 $E$ 是 $CF$ 的中点。
\item 平行补形：找到 $BD$ 的中点 $G$，证明 $DG=CE$。
\item 手拉手补形：找到 $BD$ 的中点 $F$，证明 $BF=CE$。
\item 补短：先把 $CE$ 补成整段 $CF$，再证明 $CF=BD$。
\end{enumerate}

下面各方法中的 $F$、$G$ 都独立作出。

\needspace{6\baselineskip}\textbf{方法一：延长 $BA$，与 $CE$ 的延长线相交。}

延长 $BA$、$CE$，交于 $F$。

\esdiagram{\esdiExampleExtend}

由于 $AB\perp AC$、$BD\perp CF$，按图有 $\angle ABD=\angle ACF$。又 $\angle BAD=\angle CAF=90^\circ$、$AB=AC$，故

\[
\triangle ABD\cong\triangle ACF\quad(\mathrm{ASA}),
\qquad BD=CF.
\]

在 $\triangle BFC$ 中，$\angle FBC=45^\circ$，$\angle FCB=90^\circ-22.5^\circ=67.5^\circ$，所以 $\angle BFC=67.5^\circ$，从而 $BF=BC$。

因为 $BE\perp FC$，在等腰三角形 $BFC$ 中，底边上的高也是中线，故 $FE=EC$，$CF=2$$CE$。因此 $BD=2$$CE$。

\needspace{6\baselineskip}\textbf{方法二：过 $D$ 作平行线，把二倍关系拆成两半。}

过 $D$ 作 $DF\parallel AB$，交 $BC$ 于 $F$；过 $F$ 作 $FG\perp BD$，垂足为 $G$。

\esdiagram{\esdiExampleParallel}

由 $DF\parallel AB$、$BD$ 平分 $\angle ABC$，得 $\angle DBF=\angle BDF=22.5^\circ$，所以 $BF=DF$。

又 $DF\perp DC$，$\angle DCF=\angle ACB=45^\circ$，故 $\triangle DCF$ 是等腰直角三角形，$DF=DC$。

在 $\mathrm{Rt}\triangle BFG$ 与 $\mathrm{Rt}\triangle DFG$ 中，$BF=DF$，$FG$ 公共，故两三角形全等，$BG=DG$。因此 $BD=2$$DG$。

由于 $DF\perp DC$、$DG\perp CE$，按图有 $\angle FDG=\angle DCE$。结合两个直角、$DF=DC$，得

\[
\triangle DFG\cong\triangle CDE\quad(\mathrm{AAS}),
\qquad DG=CE.
\]

故 $BD=2$$DG=2$$CE$。

\needspace{6\baselineskip}\textbf{方法三：在 $A$ 处添直角，补出手拉手。}

过 $A$ 作 $AF\perp AE$，交 $BD$ 于 $F$。

\esdiagram{\esdiExampleHand}

由 $AB\perp AC$、$AF\perp AE$，按图有 $\angle BAF=\angle CAE$。又 $\angle ABF=22.5^\circ$，由 $AC\perp AB$、$CE\perp BD$，得 $\angle ACE=22.5^\circ$。结合 $AB=AC$，有

\[
\triangle ABF\cong\triangle ACE\quad(\mathrm{ASA}),
\qquad AF=AE,\quad BF=CE.
\]

所以 $\triangle AFE$ 是等腰直角三角形，$\angle AFE=45^\circ$。同时 $\angle BAF=\angle CAE=22.5^\circ$，因此在 $\triangle ABF$ 中，$AF=BF$。

$F$、$D$、$E$ 依次共线，所以 $\angle AFD=45^\circ$；又 $\angle FAD=90^\circ-22.5^\circ=67.5^\circ$，故 $\angle ADF=67.5^\circ$，从而 $FD=AF$。

因此 $BF=AF=FD$，$F$ 是 $BD$ 的中点。于是 $BD=2$$BF=2$$CE$。

\needspace{6\baselineskip}\textbf{方法四：先补短，把 $CE$ 倍长。}

延长 $CE$ 至 $F$，使 $EF=EC$，连接 $BF$。

\esdiagram{\esdiExampleReflect}

因为 $BE\perp CF$、$CE=EF$、$BE$ 公共，

\[
\triangle BEC\cong\triangle BEF\quad(\mathrm{SAS}).
\]

所以 $\angle EBF=\angle EBC=22.5^\circ$。$F$、$C$ 在 $BE$ 的两侧，故 $\angle CBF=45^\circ=\angle CBA$，即 $B$、$A$、$F$ 共线，且 $F$ 在 $BA$ 经过 $A$ 的延长线上。

再由 $AB=AC$、$\angle BAD=\angle CAF=90^\circ$、$\angle ABD=\angle ACF$，得 $\triangle ABD\cong \triangle ACF$（$\mathrm{ASA}$），所以 $BD=CF$。

又 $CF=CE+EF=2$$CE$，故 $BD=2$$CE$。

方法一和方法四最终落到同一个辅助点 $F$：方法一先通过相交找到它，再证明 $CE=EF$；方法四先由 $CE=EF$ 作出它，再证明 $B$、$A$、$F$ 共线。作图顺序不同，关键构型相同。

\insight{看到“二倍”，先决定是倍长短段，还是平分长段；再用全等把整段或半段换到目标线段上。}

% Synced from content.md by sync_latex.py.
\ifteachernotes
\section{\texorpdfstring{辅助线总结}{辅助线总结}}

\begin{center}
\small
\begin{tabularx}{\linewidth}{LLL}
\toprule
看见的条件或目标 & 可以先作的线 & 希望得到什么 \\
\midrule
已知外角平分线 & 从 $A$ 向 $BD$、$CD$ 作双垂 & 到角两边的距离相等 \\
外角相等，想造对称全等 & 延长 $CD$，取 $DE=DB$；或延长 $BD$，取 $DF=DC$ & $\mathrm{SAS}$ 全等，再借等腰换角 \\
单个 $45^\circ$ 或 $135^\circ$ & 过 $A$ 作 $AD$ 的垂线，交相应边或延长线 & 等腰直角三角形，接手拉手 \\
等腰直角 $ABC$，需搬运长度 & 从 $B$、$C$ 向 $AD$ 作双垂 & $K$ 型全等，直角边交叉对应 \\
目标为二倍关系 & 倍长短段，或给长段找中点 & 将二倍关系转为整段相等或半段相等 \\
例 $2$ 的角平分线 & 平行线补等腰；垂线取中点 & 把角平分线条件转成中点与等长 \\
\bottomrule
\end{tabularx}
\end{center}

两种“双垂”要分清：\textbf{$A$ 向 $BD$、$CD$ 作双垂，是利用角平分线的等距离；$B$、$C$ 向 $AD$ 作双垂，是利用等腰直角 $ABC$ 的 $K$ 型全等。} 它们的出发点和全等对应顺序都不同。

\needspace{4\baselineskip}\subsection{\texorpdfstring{使用时注意}{使用时注意}}

\begin{enumerate}
\item $DA$ 平分的是 $\angle BDX$，参与比较的两个角是 $\angle BDA$ 与 $\angle ADX$。
\item 条件与结论交换后，重新检查全等判定。等边待证时不能用它作 $\mathrm{HL}$ 条件；直角待证时不能提前用手拉手。
\item $K$ 型全等推出的是 $BM=AN$、$AM=CN$。写完对应边后，再按点的顺序加减公共部分，不能把同一条垂线上的两段凭图认成相等。
\end{enumerate}
\fi

\endgroup

% END GENERATED: chapter-body
\end{document}
